% !TeX spellcheck = ru_RU
% !TeX encoding = UTF-8
%&preformat-disser
\RequirePackage[l2tabu,orthodox]{nag} % Раскомментировав, можно в логе получать рекомендации относительно правильного использования пакетов и предупреждения об устаревших и нерекомендуемых пакетах
% Формат А4, 14pt (ГОСТ Р 7.0.11-2011, 5.3.6)
\documentclass[a4paper,14pt,oneside,openany]{memoir}

%%%%%%%%%%%%%%%%%%%%%%%%%%%%%%%%%%%%%%%%%%%%%%%%%%%%%%%%%%%%%%%%%%%%%%%%%%%%%%%%
%%%% Файл упрощённых настроек шаблона, общих для диссертации и автореферата %%%%
%%%%%%%%%%%%%%%%%%%%%%%%%%%%%%%%%%%%%%%%%%%%%%%%%%%%%%%%%%%%%%%%%%%%%%%%%%%%%%%%

%%% Режим черновика %%%
\makeatletter
\@ifundefined{c@draft}{
  \newcounter{draft}
  \setcounter{draft}{0}  % 0 --- чистовик (максимальное соблюдение ГОСТ)
                         % 1 --- черновик (отклонения от ГОСТ, но быстрая
                         %       сборка итоговых PDF)
}{}
\makeatother

%%% Использование в pdflatex шрифтов не по-умолчанию %%%
\makeatletter
\@ifundefined{c@usealtfont}{
  \newcounter{usealtfont}
  \setcounter{usealtfont}{2}    % 0 --- шрифты на базе Computer Modern
                                % 1 --- использовать пакет pscyr, при его
                                %       наличии
                                % 2 --- использовать пакет XCharter, при наличии
                                %       подходящей версии
}{}
\makeatother

%%% Использование в xelatex и lualatex семейств шрифтов %%%
\makeatletter
\@ifundefined{c@fontfamily}{
  \newcounter{fontfamily}
  \setcounter{fontfamily}{1}  % 0 --- CMU семейство. Используется как fallback;
                              % 1 --- Шрифты от MS (Times New Roman и компания)
                              % 2 --- Семейство Liberation
}{}
\makeatother

%%% Библиография %%%
\makeatletter
\@ifundefined{c@bibliosel}{
  \newcounter{bibliosel}
  \setcounter{bibliosel}{0}   % 0 --- встроенная реализация с загрузкой файла
                              %       через движок bibtex8;
                              % 1 --- реализация пакетом biblatex через движок
                              %       biber
}{}
\makeatother

%%% Предкомпиляция tikz рисунков для ускорения работы %%%
\makeatletter
\@ifundefined{c@imgprecompile}{
  \newcounter{imgprecompile}
  \setcounter{imgprecompile}{0}   % 0 --- без предкомпиляции;
                                  % 1 --- пользоваться предварительно
                                  %       скомпилированными pdf вместо генерации
                                  %       заново из tikz
}{}
\makeatother
            % общие настройки шаблона
%%% Проверка используемого TeX-движка %%%
\newif\ifxetexorluatex   % определяем новый условный оператор (http://tex.stackexchange.com/a/47579)
\ifxetex
    \xetexorluatextrue
\else
    \ifluatex
        \xetexorluatextrue
    \else
        \xetexorluatexfalse
    \fi
\fi

\newif\ifsynopsis           % Условие, проверяющее, что документ --- автореферат

\usepackage{etoolbox}[2015/08/02]               % Для продвинутой проверки разных условий
\providebool{presentation}

%%% Поля и разметка страницы %%%
\usepackage{pdflscape}                              % Для включения альбомных страниц
\usepackage{geometry}                               % Для последующего задания полей

%%% Математические пакеты %%%
\usepackage{amsthm,amsmath,amscd}   % Математические дополнения от AMS
\usepackage{amsfonts,amssymb}       % Математические дополнения от AMS
\usepackage{mathtools}              % Добавляет окружение multlined
\usepackage{xfrac}                  % Красивые дроби
\usepackage[
    locale = DE,
    list-separator       = {;\,},
    list-final-separator = {;\,},
    list-pair-separator  = {;\,},
    range-phrase={\text{\ensuremath{-}}},
    % quotient-mode        = fraction, % красивые дроби могут не соответствовать ГОСТ
    fraction-function    = \sfrac,
    separate-uncertainty,
    ]{siunitx}                      % Размерности SI
\sisetup{inter-unit-product = \ensuremath{{}\cdot{}}}

% Кириллица в нумерации subequations
% Для правильной работы требуется выполнение сразу после загрузки пакетов
\patchcmd{\subequations}{\def\theequation{\theparentequation\alph{equation}}}
{\def\theequation{\theparentequation\asbuk{equation}}}
{\typeout{subequations patched}}{\typeout{subequations not patched}}

%%%% Установки для размера шрифта 14 pt %%%%
%% Формирование переменных и констант для сравнения (один раз для всех подключаемых файлов)%%
%% должно располагаться до вызова пакета fontspec или polyglossia, потому что они сбивают его работу
\newlength{\curtextsize}
\newlength{\bigtextsize}
\setlength{\bigtextsize}{13.9pt}

\makeatletter
%\show\f@size                                       % неплохо для отслеживания, но вызывает стопорение процесса, если документ компилируется без команды  -interaction=nonstopmode
\setlength{\curtextsize}{\f@size pt}
\makeatother

%%% Кодировки и шрифты %%%
\ifxetexorluatex
    \PassOptionsToPackage{no-math}{fontspec}        % https://tex.stackexchange.com/a/26295/104425
    \usepackage{polyglossia}[2014/05/21]            % Поддержка многоязычности (fontspec подгружается автоматически)
\else
   %%% Решение проблемы копирования текста в буфер кракозябрами
    \ifnumequal{\value{usealtfont}}{0}{}{
        \input glyphtounicode.tex
        \input glyphtounicode-cmr.tex %from pdfx package
        \pdfgentounicode=1
    }
    \usepackage{cmap}                               % Улучшенный поиск русских слов в полученном pdf-файле
    \ifnumequal{\value{usealtfont}}{2}{}{
        \defaulthyphenchar=127                      % Если стоит до fontenc, то переносы не впишутся в выделяемый текст при копировании его в буфер обмена
    }
    \usepackage{textcomp}
    \usepackage[T1,T2A]{fontenc}                    % Поддержка русских букв
    \ifnumequal{\value{usealtfont}}{1}{% Используется pscyr, при наличии
        \IfFileExists{pscyr.sty}{\usepackage{pscyr}}{}  % Подключение pscyr
    }{}
    \usepackage[utf8]{inputenc}[2014/04/30]         % Кодировка utf8
    \usepackage[english, russian]{babel}[2014/03/24]% Языки: русский, английский
    \ifnumequal{\value{usealtfont}}{2}{
        % http://dxdy.ru/post1238763.html#p1238763
        \usepackage[scaled=0.960]{XCharter}[2017/12/19] % Подключение русифицированных шрифтов XCharter
        \usepackage[charter, vvarbb, scaled=1.048]{newtxmath}[2017/12/14]
        \ifpresentation
        \else
            \setDisplayskipStretch{-0.078}
        \fi
    }{}
\fi

%%% Оформление абзацев %%%
\usepackage{indentfirst}                            % Красная строка

%%% Цвета %%%
\ifpresentation
\else
    \usepackage[dvipsnames, table, hyperref]{xcolor} % Совместимо с tikz
\fi

%%% Таблицы %%%
\usepackage{longtable,ltcaption} % Длинные таблицы
\usepackage{multirow,makecell}   % Улучшенное форматирование таблиц
\usepackage{tabu, tabulary}      % таблицы с автоматически подбирающейся
                                 % шириной столбцов (tabu обязательно
                                 % до hyperref вызывать)
\usepackage{threeparttable}      % автоматический подгон ширины подписи таблицы

%%% Общее форматирование
\usepackage{soulutf8}                               % Поддержка переносоустойчивых подчёркиваний и зачёркиваний
\usepackage{icomma}                                 % Запятая в десятичных дробях

%%% Оптимизация расстановки переносов и длины последней строки абзаца
\IfFileExists{impnattypo.sty}{% проверка установленности пакета impnattypo
    \ifluatex
        \ifnumequal{\value{draft}}{1}{% Черновик
            \usepackage[hyphenation, lastparline, nosingleletter, homeoarchy,
            rivers, draft]{impnattypo}
        }{% Чистовик
            \usepackage[hyphenation, lastparline, nosingleletter]{impnattypo}
        }
    \else
        \usepackage[hyphenation, lastparline]{impnattypo}
    \fi
}{}

%% Векторная графика

\usepackage{tikz}                   % Продвинутый пакет векторной графики
\usetikzlibrary{chains}             % Для примера tikz рисунка
\usetikzlibrary{shapes.geometric}   % Для примера tikz рисунка
\usetikzlibrary{shapes.symbols}     % Для примера tikz рисунка
\usetikzlibrary{arrows}             % Для примера tikz рисунка

%%% Гиперссылки %%%
\usepackage{hyperref}[2012/11/06]

%%% Изображения %%%
\usepackage{graphicx}[2014/04/25]                   % Подключаем пакет работы с графикой

%%% Счётчики %%%
\usepackage[figure,table]{totalcount}               % Счётчик рисунков и таблиц
\usepackage{totcount}                               % Пакет создания счётчиков на основе последнего номера подсчитываемого элемента (может требовать дважды компилировать документ)
\usepackage{totpages}                               % Счётчик страниц, совместимый с hyperref (ссылается на номер последней страницы). Желательно ставить последним пакетом в преамбуле

%%% Продвинутое управление групповыми ссылками (пока только формулами) %%%
\ifpresentation
\else
    \usepackage[russian]{cleveref} % cleveref имеет сложности со считыванием
    % языка из babel. Такое решение русификации вывода выбрано вместо
    % определения в documentclass из опасности что-то лишнее передать во все
    % остальные пакеты, включая библиографию.
    \creflabelformat{equation}{#2#1#3} % Формат по умолчанию ставил круглые
    % скобки вокруг каждого номера ссылки, теперь просто номера ссылок без
    % какого-либо дополнительного оформления
    \crefrangelabelformat{equation}{#3#1#4\cyrdash#5#2#6} % Интервалы в русском
    % языке принято делать через тире, если иное не оговорено

    % решение проблемы с "и" в \labelcref
    % https://tex.stackexchange.com/a/455124/104425
    \ifxetexorluatex
        \DeclareTextSymbol{\cyri}\UnicodeEncodingName{"0438} % и
    \fi

    % Добавление возможности использования пробелов в \labelcref
    % https://tex.stackexchange.com/a/340502/104425
    \usepackage{kvsetkeys}
    \makeatletter
    \let\org@@cref\@cref
    \renewcommand*{\@cref}[2]{%
        \edef\process@me{%
            \noexpand\org@@cref{#1}{\zap@space#2 \@empty}%
        }\process@me
    }
    \makeatother

    \newcommand{\eqrefs}[1]{(\labelcref{#1})}
    \newcommand{\refs}[1]{\labelcref{#1}}
\fi

\ifnumequal{\value{draft}}{1}{% Черновик
    \usepackage[firstpage]{draftwatermark}
    \SetWatermarkText{DRAFT}
    \SetWatermarkFontSize{14pt}
    \SetWatermarkScale{15}
    \SetWatermarkAngle{45}
}{}

%%% Исправление положения якорей подписей (под)рисунков %%%
% Без hypcap и патча, при клике по ссылке на подрисунок, просмотрщик pdf прыгает "к подписи" а не "к рисунку".
% Подробнее: https://github.com/AndreyAkinshin/Russian-Phd-LaTeX-Dissertation-Template/issues/238
% (!) Даже с патчем, если мешать в одной фиге разные типы подфиг (subbottom и subcaption) - ссылки всё равно будут работать неправильно  (см. https://www.overleaf.com/read/czmbmmtnqrrg ).
\ifpresentation
\else
    \usepackage[all]{hypcap}

    \makeatletter
    \ltx@ifclasslater{memoir}{2018/12/13}{
        % Предполагается, что в следующей версии класс будет исправлен
        \typeout{Assuming this version of memoir is free from the jumping-to-caption bug.}
    }{
        \usepackage{xpatch}

        \newcommand\mem@step@subcounter{\refstepcounter{sub\@captype}\@contkeep}

        \xpatchcmd{\@memsubbody}%
        {\refstepcounter{sub\@captype}\@contkeep}% search pattern
        {}% replacement
        {\typeout{@memsubbody is patched}}%
        {\typeout{@memsubbody is NOT patched}}%

        \xpatchcmd{\@memcontsubbody}%
        {\refstepcounter{sub\@captype}\@contkeep}% pattern
        {}% replacement
        {\typeout{@memcontsubbody is patched}}%
        {\typeout{@memcontsubbody is NOT patched}}%

        \xpatchcmd{\@memsubfloat}%
        {\vbox\bgroup}% search pattern
        {\vbox\bgroup\mem@step@subcounter}% replacement
        {\typeout{@memsubfloat patch is ok}}%
        {\typeout{@memsubfloat patch is NOT ok}}%

        \xpatchcmd{\subcaption}%
        {\refstepcounter{sub\@captype}}% search pattern
        {\H@refstepcounter{sub\@captype}}% replacement
        {\typeout{subcaption second patch is ok}}%
        {\typeout{subcaption second patch is NOT ok}}%
    }
    \makeatother
\fi

%%% Цитата, не приводимая в автореферате:
% возможно, актуальна только для biblatex
%\newcommand{\citeinsynopsis}[1]{\ifsynopsis\else ~\cite{#1} \fi}

% если текущий процесс запущен библиотекой tikz-external, то прекомпиляция должна быть включена
\ifdefined\tikzexternalrealjob
    \setcounter{imgprecompile}{1}
\fi

\ifnumequal{\value{imgprecompile}}{1}{% Только если у нас включена предкомпиляция
    \usetikzlibrary{external}   % подключение возможности предкомпиляции
    \tikzexternalize[prefix=images/cache/] % activate! % здесь можно указать отдельную папку для скомпилированных файлов
    \ifxetex
        \tikzset{external/up to date check={diff}}
    \fi
}{}

%% --> ScrewDriver te Blackheadborough 
%%%%%%%%%%%%%%%%%%%%%%%%%%%%%%%%%%%%%%%%%%%%%%%%%%%%%%%%%%%%%%%%%%%%%%%%%%%%%%%%%%%%
%% Пакет, обеспечивающий адекватную расстановку "русских" кавычек
\usepackage[russian]{quotmark}
%% Теперь для точной расстановки кавычек можно использовать вложенные команды tqt:
%% \tqt{а \tqt{б} в} = ,,а <<б>> в''
%%%%%%%%%%%%%%%%%%%%%%%%%%%%%%%%%%%%%%%%%%%%%%%%%%%%%%%%%%%%%%%%%%%%%%%%%%%%%%%%%%%%
%% Химические пакеты
\usepackage[version=4]{mhchem}
\usepackage{chemfig}
\usepackage{chemnum}
\makeatletter
\def\CF@node@content{\expandafter\expandafter\expandafter
    \printatom\expandafter\expandafter\expandafter
    {\csname atom@\number\CF@cnt@atomnumber\endcsname}%
    \ensuremath{\CF@node@strut}%
}
\makeatother
%\renewcommand*{\printatom}[1]{{\ce{#1}}}
\renewcommand*\printatom[1]{\ensuremath{\mathsf{#1}}} % from the ChemFig manual
%%%%%%%%%%%%%%%%%%%%%%%%%%%%%%%%%%%%%%%%%%%%%%%%%%%%%%%%%%%%%%%%%%%%%%%%%%%%%%%%%%%%
%% Relied upon \compoundstyle from *chemcompounds*
\usepackage{chemcompounds}
\newcommand{\ConfName}[1]{{\compoundstyle{#1}}}
%
\newcommand{\BB}[1]{{#1}\ConfName{ВВ}}
\newcommand{\BC}[1]{{#1}\ConfName{ВК}}
\newcommand{\CB}[1]{{#1}\ConfName{КВ}}
\newcommand{\CC}[1]{{#1}\ConfName{КК}}
\newcommand{\TT}[1]{{#1}\ConfName{ТТ}}
%%
%%%%%%%%%%%%%%%%%%%%%%%%%%%%%%%%%%%%%%%%%%%%%%%%%%%%%%%%%%%%%%%%%%%%%%%%%%%%%%%%%%%%
%% chemfig setup
\newcommand{\boldbondwidth}{3pt}
\setchemfig{bond join=true}
\setchemfig{atom sep=1.75em}
\setchemfig{bond style={line width=1.15pt}}
\setchemfig{bond offset=2pt}
\setchemfig{cram width =\boldbondwidth}
\setchemfig{cram dash width=1pt}
\setchemfig{cram dash sep=2pt}
\setchemfig{double bond sep=2.5pt}
\setchemfig{angle increment=15}
\setchemfig{cycle radius coeff=0.5}
\setchemfig{compound sep=3em}
\setchemfig{arrow offset=10pt}
%%
%%%%%%%%%%%%%%%%%%%%%%%%%%%%%%%%%%%%%%%%%%%%%%%%%%%%%%%%%%%%%%%%%%%%%%%%%%%%%%%%%%%%
% подпись на связях в молекуле, нарисованной в пакете *chemfig*
\newcommand\namebond[4][7.5pt]{\chemmove{\path(#2)--(#3)node[midway,sloped,yshift=#1]{#4};}}

\newcommand\arcbetweennodes[3]{%
  \pgfmathanglebetweenpoints{\pgfpointanchor{#1}{center}}{\pgfpointanchor{#2}{center}}%
  \let#3\pgfmathresult}
%%
\newcommand\arclabel[6][stealth-stealth,shorten <=1pt,shorten >=1pt]{%
  \chemmove{%
    \arcbetweennodes{#4}{#3}\anglestart \arcbetweennodes{#4}{#5}\angleend
    \draw[#1]([shift=(\anglestart:#2)]#4)arc(\anglestart:\angleend:#2);
    \pgfmathparse{(\anglestart+\angleend)/2}\let\anglestart\pgfmathresult
    \node[shift=(\anglestart:#2+1pt)#4,anchor=\anglestart+180,rotate=\anglestart+90,inner sep=0pt, outer sep=0pt]at(#4){#6};}
}
%%
\newcommand{\ChemPicture}[1]{\begin{tikzpicture}\chemfig{#1}\end{tikzpicture}}
%%
%%%%%%%%%%%%%%%%%%%%%%%%%%%%%%%%%%%%%%%%%%%%%%%%%%%%%%%%%%%%%%%%%%%%%%%%%%%%%%%%%%%%%%
%% Немного математических символов
\newcommand{\AGroup}[1]{{\ensuremath{\mathbb {#1}}}}
\newcommand{\SymGroup}[2]{\ensuremath{\AGroup{#1}_{#2}}}
%%
%%%%%%%%%%%%%%%%%%%%%%%%%%%%%%%%%%%%%%%%%%%%%%%%%%%%%%%%%%%%%%%%%%%%%%%%%%%%%%%%%%%%%%
         % Пакеты общие для диссертации и автореферата
\synopsisfalse                      % Этот документ --- не автореферат
\input{Dissertation/dispackages}    % Пакеты для диссертации
\input{Dissertation/userpackages}   % Пакеты для специфических пользовательских задач

%%%%%%%%%%%%%%%%%%%%%%%%%%%%%%%%%%%%%%%%%%%%%%%%%%%%%%%%%%%%%%%%%%%%%%%%%%%%%%%%
%%%% Файл упрощённых настроек шаблона, общих для диссертации и автореферата %%%%
%%%%%%%%%%%%%%%%%%%%%%%%%%%%%%%%%%%%%%%%%%%%%%%%%%%%%%%%%%%%%%%%%%%%%%%%%%%%%%%%

%%% Режим черновика %%%
\makeatletter
\@ifundefined{c@draft}{
  \newcounter{draft}
  \setcounter{draft}{0}  % 0 --- чистовик (максимальное соблюдение ГОСТ)
                         % 1 --- черновик (отклонения от ГОСТ, но быстрая
                         %       сборка итоговых PDF)
}{}
\makeatother

%%% Использование в pdflatex шрифтов не по-умолчанию %%%
\makeatletter
\@ifundefined{c@usealtfont}{
  \newcounter{usealtfont}
  \setcounter{usealtfont}{2}    % 0 --- шрифты на базе Computer Modern
                                % 1 --- использовать пакет pscyr, при его
                                %       наличии
                                % 2 --- использовать пакет XCharter, при наличии
                                %       подходящей версии
}{}
\makeatother

%%% Использование в xelatex и lualatex семейств шрифтов %%%
\makeatletter
\@ifundefined{c@fontfamily}{
  \newcounter{fontfamily}
  \setcounter{fontfamily}{1}  % 0 --- CMU семейство. Используется как fallback;
                              % 1 --- Шрифты от MS (Times New Roman и компания)
                              % 2 --- Семейство Liberation
}{}
\makeatother

%%% Библиография %%%
\makeatletter
\@ifundefined{c@bibliosel}{
  \newcounter{bibliosel}
  \setcounter{bibliosel}{0}   % 0 --- встроенная реализация с загрузкой файла
                              %       через движок bibtex8;
                              % 1 --- реализация пакетом biblatex через движок
                              %       biber
}{}
\makeatother

%%% Предкомпиляция tikz рисунков для ускорения работы %%%
\makeatletter
\@ifundefined{c@imgprecompile}{
  \newcounter{imgprecompile}
  \setcounter{imgprecompile}{0}   % 0 --- без предкомпиляции;
                                  % 1 --- пользоваться предварительно
                                  %       скомпилированными pdf вместо генерации
                                  %       заново из tikz
}{}
\makeatother
      % Упрощённые настройки шаблона

\input{common/newnames}         % Новые переменные, для всего проекта

% !TeX spellcheck = ru_RU
%%% Основные сведения %%%
\newcommand{\thesisAuthorLastName}{Писарев}
\newcommand{\thesisAuthorOtherNames}{Сергей Анатольевич}
\newcommand{\thesisAuthorInitials}{С.\:А.}
\newcommand{\thesisAuthor}             % Диссертация, ФИО автора
{%
  \texorpdfstring{% \texorpdfstring takes two arguments and uses the first for (La)TeX and the second for pdf
    \thesisAuthorLastName~\thesisAuthorOtherNames% так будет отображаться на титульном листе или в тексте, где будет использоваться переменная
  }{%
    \thesisAuthorLastName, \thesisAuthorOtherNames% эта запись для свойств pdf-файла. В таком виде, если pdf будет обработан программами для сбора библиографических сведений, будет правильно представлена фамилия.
  }
}
\newcommand{\thesisAuthorShort}        % Диссертация, ФИО автора инициалами
{\thesisAuthorInitials~\thesisAuthorLastName}
%\newcommand{\thesisUdk}                % Диссертация, УДК
%{\todo{xxx.xxx}}
\newcommand{\thesisTitle}              % Диссертация, название
{Количественный конформационный анализ производных бицикло{[3.3.1]}нонана и внутримолекулярные невалентные взаимодействия}
\newcommand{\thesisSpecialtyNumber}    % Диссертация, специальность, номер
{02.00.03}
\newcommand{\thesisSpecialtyTitle}     % Диссертация, специальность, название (название взято с сайта ВАК для примера)
{Органическая химия}
\newcommand{\thesisSpecialtyTwoNumber} % Диссертация, вторая специальность, номер
{02.00.17}
\newcommand{\thesisSpecialtyTwoTitle}  % Диссертация, вторая специальность, название
{Математическая и~квантовая химия}
\newcommand{\thesisDegree}             % Диссертация, ученая степень
{кандидата химических наук}
\newcommand{\thesisDegreeShort}        % Диссертация, ученая степень, краткая запись
{канд. хим. наук}
\newcommand{\thesisCity}               % Диссертация, город написания диссертации
{Москва}
\newcommand{\thesisYear}               % Диссертация, год написания диссертации
{2020}
\newcommand{\thesisOrganization}       % Диссертация, организация
{%Федеральное государственное образовательное учреждение высшего образования \tqt{
  Московский государственный университет им.~М.\:В.~Ломоносова. Химический факультет. Кафедра медицинской химии и тонкого органического синтеза.
  %\tqt{АББРЕВИАТУРА}} 
}
\newcommand{\thesisOrganizationShort}  % Диссертация, краткое название организации для доклада
{\todo{НазУчДисРаб}}

\newcommand{\thesisInOrganization}     % Диссертация, организация в предложном падеже: Работа выполнена в ...
{\todo{учреждении с~длинным длинным длинным длинным названием, в~котором
    выполнялась данная диссертационная работа}}

%% \newcommand{\supervisorDead}{}           % Рисовать рамку вокруг фамилии
\newcommand{\supervisorFio}              % Научный руководитель, ФИО
{Палюлин Владимир Александрович}
\newcommand{\supervisorRegalia}          % Научный руководитель, регалии
{к.~х.~н., в.~н.~c.}
\newcommand{\supervisorFioShort}         % Научный руководитель, ФИО
{В.\:А.~Палюлин}
\newcommand{\supervisorRegaliaShort}     % Научный руководитель, регалии
{к.~х.~н., в.~н.~c.}

%% \newcommand{\supervisorTwoDead}{}        % Рисовать рамку вокруг фамилии
%% \newcommand{\supervisorTwoFio}           % Второй научный руководитель, ФИО
%% {\todo{Фамилия Имя Отчество}}
%% \newcommand{\supervisorTwoRegalia}       % Второй научный руководитель, регалии
%% {\todo{уч. степень, уч. звание}}
%% \newcommand{\supervisorTwoFioShort}      % Второй научный руководитель, ФИО
%% {\todo{И.\,О.~Фамилия}}
%% \newcommand{\supervisorTwoRegaliaShort}  % Второй научный руководитель, регалии
%% {\todo{уч.~ст.,~уч.~зв.}}

\newcommand{\opponentOneFio}           % Оппонент 1, ФИО
{\todo{Фамилия Имя Отчество}}
\newcommand{\opponentOneRegalia}       % Оппонент 1, регалии
{\todo{доктор физико-математических наук, профессор}}
\newcommand{\opponentOneJobPlace}      % Оппонент 1, место работы
{\todo{Не очень длинное название для места работы}}
\newcommand{\opponentOneJobPost}       % Оппонент 1, должность
{\todo{старший научный сотрудник}}

\newcommand{\opponentTwoFio}           % Оппонент 2, ФИО
{\todo{Фамилия Имя Отчество}}
\newcommand{\opponentTwoRegalia}       % Оппонент 2, регалии
{\todo{кандидат физико-математических наук}}
\newcommand{\opponentTwoJobPlace}      % Оппонент 2, место работы
{\todo{Основное место работы c длинным длинным длинным длинным названием}}
\newcommand{\opponentTwoJobPost}       % Оппонент 2, должность
{\todo{старший научный сотрудник}}

%% \newcommand{\opponentThreeFio}         % Оппонент 3, ФИО
%% {\todo{Фамилия Имя Отчество}}
%% \newcommand{\opponentThreeRegalia}     % Оппонент 3, регалии
%% {\todo{кандидат физико-математических наук}}
%% \newcommand{\opponentThreeJobPlace}    % Оппонент 3, место работы
%% {\todo{Основное место работы c длинным длинным длинным длинным названием}}
%% \newcommand{\opponentThreeJobPost}     % Оппонент 3, должность
%% {\todo{старший научный сотрудник}}

\newcommand{\leadingOrganizationTitle} % Ведущая организация, дополнительные строки. Удалить, чтобы не отображать в автореферате
{\todo{Федеральное государственное бюджетное образовательное учреждение высшего
    профессионального образования с~длинным длинным длинным длинным названием}}

\newcommand{\defenseDate}              % Защита, дата
{\todo{DD mmmmmmmm YYYY~г.~в~XX часов}}
\newcommand{\defenseCouncilNumber}     % Защита, номер диссертационного совета
{\todo{Д\,123.456.78}}
\newcommand{\defenseCouncilTitle}      % Защита, учреждение диссертационного совета
{\todo{Название учреждения}}
\newcommand{\defenseCouncilAddress}    % Защита, адрес учреждение диссертационного совета
{\todo{Адрес}}
\newcommand{\defenseCouncilPhone}      % Телефон для справок
{\todo{+7~(0000)~00-00-00}}

\newcommand{\defenseSecretaryFio}      % Секретарь диссертационного совета, ФИО
{\todo{Фамилия Имя Отчество}}
\newcommand{\defenseSecretaryRegalia}  % Секретарь диссертационного совета, регалии
{\todo{д-р~физ.-мат. наук}}            % Для сокращений есть ГОСТы, например: ГОСТ Р 7.0.12-2011 + http://base.garant.ru/179724/#block_30000

\newcommand{\synopsisLibrary}          % Автореферат, название библиотеки
{\todo{Название библиотеки}}
\newcommand{\synopsisDate}             % Автореферат, дата рассылки
{\todo{DD mmmmmmmm YYYY года}}

% To avoid conflict with beamer class use \providecommand
\providecommand{\keywords}%            % Ключевые слова для метаданных PDF диссертации и автореферата
{}             % Основные сведения
\input{common/fonts}            % Определение шрифтов (частичное)
\input{common/styles}           % Стили общие для диссертации и автореферата
\input{Dissertation/disstyles}  % Стили для диссертации
\input{Dissertation/userstyles} % Стили для специфических пользовательских задач

%%% Библиография. Выбор движка для реализации %%%
% Здесь только проверка установленного ключа. Сама настройка выбора движка
% размещена в common/setup.tex
\ifnumequal{\value{bibliosel}}{0}{%
    \input{biblio/predefined}   % Встроенная реализация с загрузкой файла через движок bibtex8
}{
    %%% Реализация библиографии пакетами biblatex и biblatex-gost с использованием движка biber %%%

\usepackage{csquotes} % biblatex рекомендует его подключать. Пакет для оформления сложных блоков цитирования.
%%% Загрузка пакета с основными настройками %%%
\makeatletter
\ifnumequal{\value{draft}}{0}{% Чистовик
\usepackage[%
backend=biber,% движок
bibencoding=utf8,% кодировка bib файла
sorting=none,% настройка сортировки списка литературы
style=gost-numeric,% стиль цитирования и библиографии (по ГОСТ)
language=autobib,% получение языка из babel/polyglossia, default: autobib % если ставить autocite или auto, то цитаты в тексте с указанием страницы, получат указание страницы на языке оригинала
autolang=other,% многоязычная библиография
clearlang=true,% внутренний сброс поля language, если он совпадает с языком из babel/polyglossia
defernumbers=true,% нумерация проставляется после двух компиляций, зато позволяет выцеплять библиографию по ключевым словам и нумеровать не из большего списка
sortcites=true,% сортировать номера затекстовых ссылок при цитировании (если в квадратных скобках несколько ссылок, то отображаться будут отсортированно, а не абы как)
doi=false,% Показывать или нет ссылки на DOI
isbn=false,% Показывать или нет ISBN, ISSN, ISRN
]{biblatex}[2016/09/17]
\ltx@iffilelater{biblatex-gost.def}{2017/05/03}%
{\toggletrue{bbx:gostbibliography}%
\renewcommand*{\revsdnamepunct}{\addcomma}}{}
}{%Черновик
\usepackage[%
backend=biber,% движок
bibencoding=utf8,% кодировка bib файла
sorting=none,% настройка сортировки списка литературы
% defernumbers=true, % откомментируйте, если требуется правильная нумерация ссылок на литературу в режиме черновика. Замедляет сборку
]{biblatex}[2016/09/17]%
}
\makeatother

\ifxetexorluatex
\else
% Исправление случая неподдержки знака номера в pdflatex
    \DefineBibliographyStrings{russian}{number={\textnumero}}
\fi

\ifsynopsis
\ifnumgreater{\value{usefootcite}}{0}{
    \ExecuteBibliographyOptions{autocite=footnote}
    \newbibmacro*{cite:full}{%
        \printtext[bibhypertarget]{%
            \usedriver{%
                \DeclareNameAlias{sortname}{default}%
            }{%
                \thefield{entrytype}%
            }%
        }%
        \usebibmacro{shorthandintro}%
    }
    \DeclareCiteCommand{\smartcite}[\mkbibfootnote]{%
        \usebibmacro{prenote}%
    }{%
        \usebibmacro{citeindex}%
        \usebibmacro{cite:full}%
    }{%
        \multicitedelim%
    }{%
        \usebibmacro{postnote}%
    }
}{}
\fi

%%% Подключение файлов bib %%%
\addbibresource[label=bl-external]{biblio/external.bib}
\addbibresource[label=bl-author]{biblio/author.bib}

%http://tex.stackexchange.com/a/141831/79756
%There is a way to automatically map the language field to the langid field. The following lines in the preamble should be enough to do that.
%This command will copy the language field into the langid field and will then delete the contents of the language field. The language field will only be deleted if it was successfully copied into the langid field.
\DeclareSourcemap{ %модификация bib файла перед тем, как им займётся biblatex
    \maps{
        \map{% перекидываем значения полей language в поля langid, которыми пользуется biblatex
            \step[fieldsource=language, fieldset=langid, origfieldval, final]
            \step[fieldset=language, null]
        }
        \map{% перекидываем значения полей numpages в поля pagetotal, которыми пользуется biblatex
            \step[fieldsource=numpages, fieldset=pagetotal, origfieldval, final]
            \step[fieldset=numpages, null]
        }
        \map{% перекидываем значения полей pagestotal в поля pagetotal, которыми пользуется biblatex
            \step[fieldsource=pagestotal, fieldset=pagetotal, origfieldval, final]
            \step[fieldset=pagestotal, null]
        }
        \map[overwrite]{% перекидываем значения полей shortjournal, если они есть, в поля journal, которыми пользуется biblatex
            \step[fieldsource=shortjournal, final]
            \step[fieldset=journal, origfieldval]
            \step[fieldset=shortjournal, null]
        }
        \map[overwrite]{% перекидываем значения полей shortbooktitle, если они есть, в поля booktitle, которыми пользуется biblatex
            \step[fieldsource=shortbooktitle, final]
            \step[fieldset=booktitle, origfieldval]
            \step[fieldset=shortbooktitle, null]
        }
        \map{% если в поле medium написано "Электронный ресурс", то устанавливаем поле media, которым пользуется biblatex, в значение eresource.
            \step[fieldsource=medium,
            match=\regexp{Электронный\s+ресурс},
            final]
            \step[fieldset=media, fieldvalue=eresource]
            \step[fieldset=medium, null]
        }
        \map{% использование media=text по умолчанию
            \step[fieldset=media, fieldvalue=text]
        }
        \map[overwrite]{% стираем значения всех полей issn
            \step[fieldset=issn, null]
        }
        \map[overwrite]{% стираем значения всех полей abstract, поскольку ими не пользуемся, а там бывают "неприятные" латеху символы
            \step[fieldsource=abstract]
            \step[fieldset=abstract,null]
        }
        \map[overwrite]{ % переделка формата записи даты
            \step[fieldsource=urldate,
            match=\regexp{([0-9]{2})\.([0-9]{2})\.([0-9]{4})},
            replace={$3-$2-$1$4}, % $4 вставлен исключительно ради нормальной работы программ подсветки синтаксиса, которые некорректно обрабатывают $ в таких конструкциях
            final]
        }
        \map[overwrite]{ % стираем ключевые слова
            \step[fieldsource=keywords]
            \step[fieldset=keywords,null]
        }
        % реализация foreach различается для biblatex v3.12 и v3.13.
        % Для версии v3.13 эта конструкция заменяет последующие 5 структур map
        % \map[overwrite,foreach={authorvak,authorscopus,authorwos,authorconf,authorother}]{ % записываем информацию о типе публикации в ключевые слова
        %     \step[fieldsource=$MAPLOOP,final=true]
        %     \step[fieldset=keywords,fieldvalue={,biblio$MAPLOOP},append=true]
        % }
        \map[overwrite]{ % записываем информацию о типе публикации в ключевые слова
            \step[fieldsource=authorvak,final=true]
            \step[fieldset=keywords,fieldvalue={,biblioauthorvak},append=true]
        }
        \map[overwrite]{ % записываем информацию о типе публикации в ключевые слова
            \step[fieldsource=authorscopus,final=true]
            \step[fieldset=keywords,fieldvalue={,biblioauthorscopus},append=true]
        }
        \map[overwrite]{ % записываем информацию о типе публикации в ключевые слова
            \step[fieldsource=authorwos,final=true]
            \step[fieldset=keywords,fieldvalue={,biblioauthorwos},append=true]
        }
        \map[overwrite]{ % записываем информацию о типе публикации в ключевые слова
            \step[fieldsource=authorconf,final=true]
            \step[fieldset=keywords,fieldvalue={,biblioauthorconf},append=true]
        }
        \map[overwrite]{ % записываем информацию о типе публикации в ключевые слова
            \step[fieldsource=authorother,final=true]
            \step[fieldset=keywords,fieldvalue={,biblioauthorother},append=true]
        }
        \map[overwrite]{ % добавляем ключевые слова, чтобы различать источники
            \perdatasource{biblio/external.bib}
            \step[fieldset=keywords, fieldvalue={,biblioexternal},append=true]
        }
        \map[overwrite]{ % добавляем ключевые слова, чтобы различать источники
            \perdatasource{biblio/author.bib}
            \step[fieldset=keywords, fieldvalue={,biblioauthor},append=true]
        }
        \map[overwrite]{ % добавляем ключевые слова, чтобы различать источники
            \step[fieldset=keywords, fieldvalue={,bibliofull},append=true]
        }
%        \map[overwrite]{% стираем значения всех полей series
%            \step[fieldset=series, null]
%        }
        \map[overwrite]{% перекидываем значения полей howpublished в поля organization для типа online
            \step[typesource=online, typetarget=online, final]
            \step[fieldsource=howpublished, fieldset=organization, origfieldval]
            \step[fieldset=howpublished, null]
        }
        % Так отключаем [Электронный ресурс]
%        \map[overwrite]{% стираем значения всех полей media=eresource
%            \step[fieldsource=media,
%            match={eresource},
%            final]
%            \step[fieldset=media, null]
%        }
    }
}

\ifsynopsis
\else
\DeclareSourcemap{ %модификация bib файла перед тем, как им займётся biblatex
    \maps{
        \map[overwrite]{% стираем значения всех полей addendum
            \perdatasource{biblio/author.bib}
            \step[fieldset=addendum, null] %чтобы избавиться от информации об объёме авторских статей, в отличие от автореферата
        }
    }
}
\fi

\defbibfilter{vakscopuswos}{%
    keyword=biblioauthorvak or keyword=biblioauthorscopus or keyword=biblioauthorwos
}

\defbibfilter{scopuswos}{%
    keyword=biblioauthorscopus or keyword=biblioauthorwos
}

%%% Убираем неразрывные пробелы перед двоеточием и точкой с запятой %%%
%\makeatletter
%\ifnumequal{\value{draft}}{0}{% Чистовик
%    \renewcommand*{\addcolondelim}{%
%      \begingroup%
%      \def\abx@colon{%
%        \ifdim\lastkern>\z@\unkern\fi%
%        \abx@puncthook{:}\space}%
%      \addcolon%
%      \endgroup}
%
%    \renewcommand*{\addsemicolondelim}{%
%      \begingroup%
%      \def\abx@semicolon{%
%        \ifdim\lastkern>\z@\unkern\fi%
%        \abx@puncthook{;}\space}%
%      \addsemicolon%
%      \endgroup}
%}{}
%\makeatother

%%% Правка записей типа thesis, чтобы дважды не писался автор
%\ifnumequal{\value{draft}}{0}{% Чистовик
%\DeclareBibliographyDriver{thesis}{%
%  \usebibmacro{bibindex}%
%  \usebibmacro{begentry}%
%  \usebibmacro{heading}%
%  \newunit
%  \usebibmacro{author}%
%  \setunit*{\labelnamepunct}%
%  \usebibmacro{thesistitle}%
%  \setunit{\respdelim}%
%  %\printnames[last-first:full]{author}%Вот эту строчку нужно убрать, чтобы автор диссертации не дублировался
%  \newunit\newblock
%  \printlist[semicolondelim]{specdata}%
%  \newunit
%  \usebibmacro{institution+location+date}%
%  \newunit\newblock
%  \usebibmacro{chapter+pages}%
%  \newunit
%  \printfield{pagetotal}%
%  \newunit\newblock
%  \usebibmacro{doi+eprint+url+note}%
%  \newunit\newblock
%  \usebibmacro{addendum+pubstate}%
%  \setunit{\bibpagerefpunct}\newblock
%  \usebibmacro{pageref}%
%  \newunit\newblock
%  \usebibmacro{related:init}%
%  \usebibmacro{related}%
%  \usebibmacro{finentry}}
%}{}

%\newbibmacro{string+doi}[1]{% новая макрокоманда на простановку ссылки на doi
%    \iffieldundef{doi}{#1}{\href{http://dx.doi.org/\thefield{doi}}{#1}}}

%\ifnumequal{\value{draft}}{0}{% Чистовик
%\renewcommand*{\mkgostheading}[1]{\usebibmacro{string+doi}{#1}} % ссылка на doi с авторов. стоящих впереди записи
%\renewcommand*{\mkgostheading}[1]{#1} % только лишь убираем курсив с авторов
%}{}
%\DeclareFieldFormat{title}{\usebibmacro{string+doi}{#1}} % ссылка на doi с названия работы
%\DeclareFieldFormat{journaltitle}{\usebibmacro{string+doi}{#1}} % ссылка на doi с названия журнала
%%% Тире как разделитель в библиографии традиционной руской длины:
\renewcommand*{\newblockpunct}{\addperiod\addnbspace\cyrdash\space\bibsentence}
%%% Убрать тире из разделителей элементов в библиографии:
%\renewcommand*{\newblockpunct}{%
%    \addperiod\space\bibsentence}%block punct.,\bibsentence is for vol,etc.

%%% Возвращаем запись «Режим доступа» %%%
%\DefineBibliographyStrings{english}{%
%    urlfrom = {Mode of access}
%}
%\DeclareFieldFormat{url}{\bibstring{urlfrom}\addcolon\space\url{#1}}

%%% В списке литературы обозначение одной буквой диапазона страниц англоязычного источника %%%
\DefineBibliographyStrings{english}{%
    pages = {p\adddot} %заглавность буквы затем по месту определяется работой самого biblatex
}

%%% В ссылке на источник в основном тексте с указанием конкретной страницы обозначение одной большой буквой %%%
%\DefineBibliographyStrings{russian}{%
%    page = {C\adddot}
%}

%%% Исправление длины тире в диапазонах %%%
% \cyrdash --- тире «русской» длины, \textendash --- en-dash
\DefineBibliographyExtras{russian}{%
  \protected\def\bibrangedash{%
    \cyrdash\penalty\value{abbrvpenalty}}% almost unbreakable dash
  \protected\def\bibdaterangesep{\bibrangedash}%тире для дат
}
\DefineBibliographyExtras{english}{%
  \protected\def\bibrangedash{%
    \cyrdash\penalty\value{abbrvpenalty}}% almost unbreakable dash
  \protected\def\bibdaterangesep{\bibrangedash}%тире для дат
}

%Set higher penalty for breaking in number, dates and pages ranges
\setcounter{abbrvpenalty}{10000} % default is \hyphenpenalty which is 12

%Set higher penalty for breaking in names
\setcounter{highnamepenalty}{10000} % If you prefer the traditional BibTeX behavior (no linebreaks at highnamepenalty breakpoints), set it to ‘infinite’ (10 000 or higher).
\setcounter{lownamepenalty}{10000}

%%% Set low penalties for breaks at uppercase letters and lowercase letters
%\setcounter{biburllcpenalty}{500} %управляет разрывами ссылок после маленьких букв RTFM biburllcpenalty
%\setcounter{biburlucpenalty}{3000} %управляет разрывами ссылок после больших букв, RTFM biburlucpenalty

%%% Список литературы с красной строки (без висячего отступа) %%%
%\defbibenvironment{bibliography} % переопределяем окружение библиографии из gost-numeric.bbx пакета biblatex-gost
%  {\list
%     {\printtext[labelnumberwidth]{%
%       \printfield{prefixnumber}%
%       \printfield{labelnumber}}}
%     {%
%      \setlength{\labelwidth}{\labelnumberwidth}%
%      \setlength{\leftmargin}{0pt}% default is \labelwidth
%      \setlength{\labelsep}{\widthof{\ }}% Управляет длиной отступа после точки % default is \biblabelsep
%      \setlength{\itemsep}{\bibitemsep}% Управление дополнительным вертикальным разрывом между записями. \bibitemsep по умолчанию соответствует \itemsep списков в документе.
%      \setlength{\itemindent}{\bibhang}% Пользуемся тем, что \bibhang по умолчанию принимает значение \parindent (абзацного отступа), который переназначен в styles.tex
%      \addtolength{\itemindent}{\labelwidth}% Сдвигаем правее на величину номера с точкой
%      \addtolength{\itemindent}{\labelsep}% Сдвигаем ещё правее на отступ после точки
%      \setlength{\parsep}{\bibparsep}%
%     }%
%      \renewcommand*{\makelabel}[1]{\hss##1}%
%  }
%  {\endlist}
%  {\item}

%%% Макросы автоматического подсчёта количества авторских публикаций.
% Печатают невидимую (пустую) библиографию, считая количество источников.
% http://tex.stackexchange.com/a/66851/79756
%
\makeatletter
        \newtotcounter{citenum}
        \defbibenvironment{counter}
            {\setcounter{citenum}{0}\renewcommand{\blx@driver}[1]{}} % begin code: убирает весь выводимый текст
            {} % end code
            {\stepcounter{citenum}} % item code: cчитает "печатаемые в библиографию" источники

        \newtotcounter{citeauthorvak}
        \defbibenvironment{countauthorvak}
            {\setcounter{citeauthorvak}{0}\renewcommand{\blx@driver}[1]{}}
            {}
            {\stepcounter{citeauthorvak}}

        \newtotcounter{citeauthorscopus}
        \defbibenvironment{countauthorscopus}
                {\setcounter{citeauthorscopus}{0}\renewcommand{\blx@driver}[1]{}}
                {}
                {\stepcounter{citeauthorscopus}}

        \newtotcounter{citeauthorwos}
        \defbibenvironment{countauthorwos}
                {\setcounter{citeauthorwos}{0}\renewcommand{\blx@driver}[1]{}}
                {}
                {\stepcounter{citeauthorwos}}

        \newtotcounter{citeauthorother}
        \defbibenvironment{countauthorother}
                {\setcounter{citeauthorother}{0}\renewcommand{\blx@driver}[1]{}}
                {}
                {\stepcounter{citeauthorother}}

        \newtotcounter{citeauthorconf}
        \defbibenvironment{countauthorconf}
                {\setcounter{citeauthorconf}{0}\renewcommand{\blx@driver}[1]{}}
                {}
                {\stepcounter{citeauthorconf}}

        \newtotcounter{citeauthor}
        \defbibenvironment{countauthor}
                {\setcounter{citeauthor}{0}\renewcommand{\blx@driver}[1]{}}
                {}
                {\stepcounter{citeauthor}}

        \newtotcounter{citeauthorvakscopuswos}
        \defbibenvironment{countauthorvakscopuswos}
                {\setcounter{citeauthorvakscopuswos}{0}\renewcommand{\blx@driver}[1]{}}
                {}
                {\stepcounter{citeauthorvakscopuswos}}

        \newtotcounter{citeauthorscopuswos}
        \defbibenvironment{countauthorscopuswos}
                {\setcounter{citeauthorscopuswos}{0}\renewcommand{\blx@driver}[1]{}}
                {}
                {\stepcounter{citeauthorscopuswos}}

        \newtotcounter{citeexternal}
        \defbibenvironment{countexternal}
                {\setcounter{citeexternal}{0}\renewcommand{\blx@driver}[1]{}}
                {}
                {\stepcounter{citeexternal}}
\makeatother

\defbibheading{nobibheading}{} % пустой заголовок, для подсчёта публикаций с помощью невидимой библиографии
\defbibheading{pubgroup}{\section*{#1}} % обычный стиль, заголовок-секция
\defbibheading{pubsubgroup}{\noindent\textbf{#1}} % для подразделов "по типу источника"

%%%Сортировка списка литературы Русский-Английский (предварительно удалить dissertation.bbl) (начало)
%%%Источник: https://github.com/odomanov/biblatex-gost/wiki/%D0%9A%D0%B0%D0%BA-%D1%81%D0%B4%D0%B5%D0%BB%D0%B0%D1%82%D1%8C,-%D1%87%D1%82%D0%BE%D0%B1%D1%8B-%D1%80%D1%83%D1%81%D1%81%D0%BA%D0%BE%D1%8F%D0%B7%D1%8B%D1%87%D0%BD%D1%8B%D0%B5-%D0%B8%D1%81%D1%82%D0%BE%D1%87%D0%BD%D0%B8%D0%BA%D0%B8-%D0%BF%D1%80%D0%B5%D0%B4%D1%88%D0%B5%D1%81%D1%82%D0%B2%D0%BE%D0%B2%D0%B0%D0%BB%D0%B8-%D0%BE%D1%81%D1%82%D0%B0%D0%BB%D1%8C%D0%BD%D1%8B%D0%BC
%\DeclareSourcemap{
%	\maps[datatype=bibtex]{
%		\map{
%			\step[fieldset=langid, fieldvalue={tempruorder}]
%		}
%		\map[overwrite]{
%			\step[fieldsource=langid, match=russian, final]
%			\step[fieldsource=presort, 
%			match=\regexp{(.+)}, 
%			replace=\regexp{aa$1}]
%		}
%		\map{
%			\step[fieldsource=langid, match=russian, final]
%			\step[fieldset=presort, fieldvalue={az}]
%		}
%		\map[overwrite]{
%			\step[fieldsource=langid, notmatch=russian, final]
%			\step[fieldsource=presort, 
%			match=\regexp{(.+)}, 
%			replace=\regexp{za$1}]
%		}
%		\map{
%			\step[fieldsource=langid, notmatch=russian, final]
%			\step[fieldset=presort, fieldvalue={zz}]
%		}
%		\map{
%			\step[fieldsource=langid, match={tempruorder}, final]
%			\step[fieldset=langid, null]
%		}
%	}
%}
%Сортировка списка литературы (конец)

%%% Создание команд для вывода списка литературы %%%
\newcommand*{\insertbibliofull}{
    \printbibliography[keyword=bibliofull,section=0,title=\bibtitlefull]
    \ifnumequal{\value{draft}}{0}{
      \printbibliography[heading=nobibheading,env=counter,keyword=bibliofull,section=0]
    }{}
}
\newcommand*{\insertbiblioauthor}{
    \printbibliography[heading=pubgroup, section=0, keyword=biblioauthor, title=\bibtitleauthor]
}
\newcommand*{\insertbiblioauthorimportant}{
    \printbibliography[heading=pubgroup, section=2, keyword=biblioauthor, title=\bibtitleauthorimportant]
}

% Вариант вывода печатных работ автора, с группировкой по типу источника.
% Порядок команд `\printbibliography` должен соответствовать порядку в файле common/characteristic.tex
\newcommand*{\insertbiblioauthorgrouped}{
    \section*{\bibtitleauthor}
    \ifsynopsis
    \printbibliography[heading=pubsubgroup, section=0, keyword=biblioauthorvak,    title=\bibtitleauthorvak,resetnumbers=true] % Работы автора из списка ВАК (сброс нумерации)
    \else
    \printbibliography[heading=pubsubgroup, section=0, keyword=biblioauthorvak,    title=\bibtitleauthorvak,resetnumbers=false] % Работы автора из списка ВАК (сквозная нумерация)
    \fi
    \printbibliography[heading=pubsubgroup, section=0, keyword=biblioauthorwos,    title=\bibtitleauthorwos,resetnumbers=false]% Работы автора, индексируемые Web of Science
    \printbibliography[heading=pubsubgroup, section=0, keyword=biblioauthorscopus, title=\bibtitleauthorscopus,resetnumbers=false]% Работы автора, индексируемые Scopus
    \printbibliography[heading=pubsubgroup, section=0, keyword=biblioauthorconf,   title=\bibtitleauthorconf,resetnumbers=false]% Тезисы конференций
    \printbibliography[heading=pubsubgroup, section=0, keyword=biblioauthorother,  title=\bibtitleauthorother,resetnumbers=false]% Прочие работы автора
}

\newcommand*{\insertbiblioexternal}{
    \printbibliography[heading=pubgroup,    section=0, keyword=biblioexternal,     title=\bibtitlefull]
}
     % Реализация пакетом biblatex через движок biber
}

% Вывести информацию о выбранных опциях в лог сборки
\typeout{Selected options:}
\typeout{Draft mode: \arabic{draft}}
\typeout{Font: \arabic{fontfamily}}
\typeout{AltFont: \arabic{usealtfont}}
\typeout{Bibliography backend: \arabic{bibliosel}}
\typeout{Precompile images: \arabic{imgprecompile}}
% Вывести информацию о версиях используемых библиотек в лог сборки
\listfiles

%%% Управление компиляцией отдельных частей диссертации %%%
% Необходимо сначала иметь полностью скомпилированный документ, чтобы все
% промежуточные файлы были в наличии
% Затем, для вывода отдельных частей можно воспользоваться командой \includeonly
% Ниже примеры использования команды:
%
%\includeonly{Dissertation/part2}
%\includeonly{Dissertation/contents,Dissertation/appendix,Dissertation/conclusion}
%
% Если все команды закомментированы, то документ будет выведен в PDF файл полностью

%%% Данные в виде таблиц:
\pgfplotstableread{Results/CCCC-Scan.dat}{\ScanButane}
\pgfplotstableread{Results/CCOC-Scan.dat}{\ScanMeOEt}
\pgfplotstableread{Results/CCSC-Scan.dat}{\ScanMeSEt}
\pgfplotstableread{Results/CCQC-Scan.dat}{\ScanMeCOEt}

\begin{document}

\input{common/renames}                 % Переопределение именований

%%% Структура диссертации (ГОСТ Р 7.0.11-2011, 4)
% Титульный лист (ГОСТ Р 7.0.11-2001, 5.1)
\thispagestyle{empty}
\begin{center}
\thesisOrganization
\end{center}
%
\vspace{0pt plus4fill} %число перед fill = кратность относительно некоторого расстояния fill, кусками которого заполнены пустые места
\IfFileExists{images/logo.pdf}{
  \begin{minipage}[b]{0.5\linewidth}
    \begin{flushleft}
      \includegraphics[height=3.5cm]{logo}
    \end{flushleft}
  \end{minipage}%
  \begin{minipage}[b]{0.5\linewidth}
    \begin{flushright}
      На правах рукописи\\
%      \textsl {УДК \thesisUdk}
    \end{flushright}
  \end{minipage}
}{
\begin{flushright}
На правах рукописи

%\textsl {УДК \thesisUdk}
\end{flushright}
}
%
\vspace{0pt plus5fill} %число перед fill = кратность относительно некоторого расстояния fill, кусками которого заполнены пустые места
\begin{center}
{\large \thesisAuthor}
\end{center}
%
\vspace{0pt plus1fill} %число перед fill = кратность относительно некоторого расстояния fill, кусками которого заполнены пустые места
\begin{center}
\textbf {\large %\MakeUppercase
\thesisTitle}

\vspace{0pt plus2fill} %число перед fill = кратность относительно некоторого расстояния fill, кусками которого заполнены пустые места
{%\small
Специальность \thesisSpecialtyNumber\ "---

<<\thesisSpecialtyTitle>>
}

\ifdefined\thesisSpecialtyTwoNumber
{%\small
Специальность \thesisSpecialtyTwoNumber\ "---

<<\thesisSpecialtyTwoTitle>>
}
\fi

\vspace{0pt plus2fill} %число перед fill = кратность относительно некоторого расстояния fill, кусками которого заполнены пустые места
Диссертация на соискание учёной степени

\thesisDegree
\end{center}
%
\vspace{0pt plus4fill} %число перед fill = кратность относительно некоторого расстояния fill, кусками которого заполнены пустые места
\begin{flushright}
\ifdefined\supervisorTwoFio
Научные руководители:

\supervisorRegalia

\ifdefined\supervisorDead
\framebox{\supervisorFio}
\else
\supervisorFio
\fi

\supervisorTwoRegalia

\ifdefined\supervisorTwoDead
\framebox{\supervisorTwoFio}
\else
\supervisorFio
\fi
\else
Научный руководитель:

\supervisorRegalia

\ifdefined\supervisorDead
\framebox{\supervisorFio}
\else
\supervisorFio
\fi
\fi

\end{flushright}
%
\vspace{0pt plus4fill} %число перед fill = кратность относительно некоторого расстояния fill, кусками которого заполнены пустые места
{\centering\thesisCity\ "--- \thesisYear\par}
           % Титульный лист
\include{Dissertation/contents}        % Оглавление
\ifnumequal{\value{contnumfig}}{1}{}{\counterwithout{figure}{chapter}}
\ifnumequal{\value{contnumtab}}{1}{}{\counterwithout{table}{chapter}}
\include{Dissertation/introduction}    % Введение
\ifnumequal{\value{contnumfig}}{1}{\counterwithout{figure}{chapter}
}{\counterwithin{figure}{chapter}}
\ifnumequal{\value{contnumtab}}{1}{\counterwithout{table}{chapter}
}{\counterwithin{table}{chapter}}
\chapter{Конформационное поведение производных и гетероаналогов бицикло[3.3.1]нонана (обзор литературы)}

\epigraph{Quia igitur huius mundi machina sphaerica est, dicendum est in primis quid sit sphaera?}{De sphaera Lincolniensem Secundum tractatus, cap. I}

Под термином \tqt{\emph{циклическая система}} здесь и в дальнейшем мы подразумеваем циклы и полициклические каркасы в общем, рассматривая прежде всего их остовы, сосотоящие обычно из тяжёлых (неводородных) атомов.

\begin{figure}
  \centering
  %\vspace{\bigskipamount}
  \vspace{\bigskipamount}
  Двухкоординированнные~(тип~I):

  \vspace{\medskipamount}
  \begin{tabular}{cc|cc|}
    \chemfig{>[:+30]B(-H)-[:+150]} & 
    \chemfig{>[:+30]B(-R)-[:+150]} &
    \chemfig{>[:+30]C(-[:+55]R)(-[:-55]R)-[:+150]} &
    \chemfig{>[:+30]C(=X)-[:+150]} \\
  \end{tabular}
  ~
  \begin{tabular}{cc|c|}
    \chemfig{>[:+30]N(-H)-[:+150]} & 
    \chemfig{>[:+30]N(-R)-[:+150]} & 
    \chemfig{>[:+30]O-[:+150]} \\
  \end{tabular}
  ~
  \begin{tabular}{c|ccc|}
    \chemfig{>[:+30]Si(-[:+55]R)(-[:-55]R)-[:+150]} &
    \chemfig{>[:+30]P(-H)-[:+150]} & \chemfig{>[:+30]P(-R)-[:+150]} & \chemfig{>[:+30]P(-[:+55]R)(=[:-55]O)-[:+150]} \\
  \end{tabular}
  ~
  \begin{tabular}{ccc|}
    \chemfig{>[:+30]S-[:+150]} & \chemfig{>[:+30]S(=O)-[:+150]} &
    \chemfig{>[:+30]S(=[:+55]O)(=[:-55]O)-[:+150]} \\
  \end{tabular}
  
  \vspace{\bigskipamount}
  Трёхкоординированные~(тип~II):

  \vspace{\medskipamount}
  \begin{tabular}{c|c|cc|c|cc}
    \chemfig{B(<[:-120])(<:[:-60])-[:+90]} & \chemfig{C(<[:-120])(<:[:-60])(-[:+150]R)-[:+30]} &   \chemfig{N(<[:-120])(<:[:-60])-[:+30]} & \chemfig{N^+(<[:-120])(<:[:-60])(-[:+150]R)-[:+30]} & \chemfig{Si(<[:-120])(<:[:-60])(-[:+150]R)-[:+30]} &  \chemfig{P(<[:-120])(<:[:-60])-[:+30]} & \chemfig{P(<[:-120])(<:[:-60])(=[:+150]O)-[:+30]}
  \end{tabular}
  \vspace{\medskipamount}
  \caption{\label{fig:Atom:Types}Стереохимическая типизация атомов в циклических системах}
\end{figure}

Мы используем широкое понятие о \tqt{замещении} в циклических системах; то есть, если дополнительно не оговаривается иное, мы имеем в виду под замещением не только введение или изменение заместителя (\emph{замещение} в~узком смысле слова), но и \emph{гетероаналогию} или \emph{гетероаналогичное замещение} атома \ce{C} соответствующего циклического углеводорода на~гетероатом того же стереохимического типа~(рис.~\ref{fig:Atom:Types}).

Конформационный анализ молекулярных структур часто включает их разбиение на~подструктуры с дальнейшим выражением конформации структуры в~целом как суперпозиции конформаций подструктур. Система бицикло[3.3.1]нонана~\cmpd{Bicycle331} допускает несколько вариантов подобных разбиений~(рис.~\ref{fig:Decomposition:331}):

\begin{enumerate}
\item\label{item:331:Decomposition:6:6} Два 1,3-конденсированных шестичленных цикла (здесь и~в~дальнейшем циклы будем обозначать вместе с~направлением обхода): $\text{\CycleFirst} = \left(5\rightarrow 9\rightarrow 1\rightarrow 2\rightarrow 3\rightarrow 4\right)$ и~$\text{\CycleSecond} = \left(1\rightarrow 9\rightarrow 5\rightarrow 6\rightarrow 7\rightarrow 8\right)$;
\item\label{item:331:Decomposition:8:1} Восьмичленный цикл~$\text{\CycleThird} = \left(8\rightarrow 7\rightarrow 6\rightarrow 5\rightarrow 4\rightarrow 3\rightarrow 2\rightarrow 1\right)$ с одноатомным мостиком $-9-$ между положениями 1 и 5;
\item\label{item:331:Decomposition:2x5:3} Два пятичленных \emph{псевдоцикла} (группы из атомов, трактуемых как цикл, но не обязательно соединённых связями в соответствующем порядке) $\text{\CyclePseudoFirst}=(1\rightarrow 2\rightarrow 3\rightarrow 4\rightarrow 5)$ и $\text{\CyclePseudoSecond}=(5\rightarrow 6\rightarrow 7\rightarrow 8\rightarrow 1)$, примыкающих к трёхчленному псевдоциклу $(1\rightarrow 9\rightarrow 5)$;
\item\label{item:331:Decomposition:2x2:2} Три связанных трёхатомных фрагмента: $2-3-4$, $8-7-6$ (т.\,н. \tqt{крылья}) и $1-9-5$;
\end{enumerate}

\begin{figure}
\centerfloat{
\begin{tabular}{|c|c|c|c|}
%\toprule
\multicolumn{4}{c}{\chemfig{[:-30,1.25]9*6(-1(-[:180]2-[:120]3-[:+60]4-[:0]\phantom{5}?)-8-7-6-5-)}} \\
\multicolumn{4}{c}{\cmpd{Bicycle331} } \\
\midrule
\ref{item:331:Decomposition:6:6} & \ref{item:331:Decomposition:8:1} &
\ref{item:331:Decomposition:2x5:3} & \ref{item:331:Decomposition:2x2:2} \\
\chemfig{[:-30]3?-[:-60]2-[:0]1-[:+120]9-[:+60]5-[:+180]4?}\chemfig{[:-30]9*6(-1-8-7-6-5-)} &
\chemfig{[:-30]9*6(-[,,,,dash pattern=on 1pt off 1pt]1(-[:180]2-[:120]3-[:+60]4-[:0]\phantom{5}?)-8-7-6-5-[,,,,dash pattern=on 1pt off 1pt])} & 
\chemfig{[:-30]9*6(-1?[a](-[:180]2-[:120]3-[:+60]4-[:0]\phantom{5}?[a,1,{dash pattern=on 1pt off 1pt}])-8-7-6-5-)} &
\chemfig{[:-30]9*6(-1(-[:180,,,,dash pattern=on 1pt off 1pt]2-[:120]3-[:+60]4-[:0,,,,dash pattern=on 1pt off 1pt]\phantom{5}?)-[,,,,dash pattern=on 1pt off 1pt]8-7-6-[,,,,dash pattern=on 1pt off 1pt]5-)} \\
$\leftarrow$\qquad$\longleftarrow$ & $\longrightarrow$ & $\leftarrow$\quad$\rightarrow$ & \\ 
\bottomrule
\end{tabular}
\vspace{\bigskipamount}
}
\caption{\label{fig:Decomposition:331} Номенклатура положений и структурные декомпозиции системы бицикло[3.3.1]нонана. Для циклов {\CycleFirst}--{\CyclePseudoSecond} указано направление обхода }
\end{figure}

Конформация бициклического скелета~\cmpd{Bicycle331} традиционно определяется согласно разбиению~\ref{item:331:Decomposition:6:6}, т.\,е. из двух составляющих его конформаций шестичленных циклов. Поэтому нам кажется закономерным вначале кратко сформулировать хорошо известные основные положения конформационного анализа шестичленных циклов.

\begin{figure}
\centerfloat{
\begin{tabular}{c|cc|c}
\toprule\cmpd{Bicycle331:XYZ}: & 
\multicolumn{3}{c}{\chemfig{[:-30]Z*6(-(-[:180]-[:120]X-[:+60]-[:0]?)--Y---)}} \\
\midrule & \ConfName{К} & \ConfName{Т} & \ConfName{В} \\
\midrule\CycleFirst: & 
\chemfig{X?<[:-60]-[:+20,,,,line width=\boldbondwidth]>[:-20]Z-[:+120]-[:-160]?} & 
\chemfig{X?-[:-150,0.75]<[:-30,0.75]-[:+30,1.5,,,line width=\boldbondwidth]Z>[:-30,0.75]-[:-150,0.75]?} &
\chemfig{X(-[:-35]?)<[:-60]-[:+0,1.5,,,line width=\boldbondwidth]>[:+60]Z-[:-145]?}  \\
\chemfig{[:-30]X*6(---Z---)} &
&
\chemfig{?-[:-150,0.75]X<[:-30,0.75]-[:+30,1.5,,,line width=\boldbondwidth]>[:-30,0.75]Z-[:-150,0.75]?} &
\chemfig{(-[:-35]?)<[:-60]X-[:+0,1.5,,,line width=\boldbondwidth]>[:+60]-[:-145]Z?} \\
\midrule\CycleSecond: & 
\chemfig{Y?<[:-60]-[:+20,,,,line width=\boldbondwidth]>[:-20]Z-[:+120]-[:-160]?}&
\chemfig{Y?-[:-150,0.75]<[:-30,0.75]-[:+30,1.5,,,line width=\boldbondwidth]Z>[:-30,0.75]-[:-150,0.75]?} &
\chemfig{Y(-[:-35]?)<[:-60]-[:+0,1.5,,,line width=\boldbondwidth]>[:+60]Z-[:-145]?}  \\
\chemfig{[:-30]Y*6(---Z---)} & 
&
\chemfig{?-[:-150,0.75]Y<[:-30,0.75]-[:+30,1.5,,,line width=\boldbondwidth]>[:-30,0.75]Z-[:-150,0.75]?} &
\chemfig{(-[:-35]?)<[:-60]Y-[:+0,1.5,,,line width=\boldbondwidth]>[:+60]-[:-145]Z?} \\
\midrule\cmpd{Cyclohexane}: &
\ChemPicture{?<[:-60]-[:+20,,,,line width=\boldbondwidth]>[:-20]-[:+120]-[:-160]?} & 
\ChemPicture{?-[:-150,0.75]<[:-30,0.75]-[:+30,1.5,,,line width=\boldbondwidth]>[:-30,0.75]-[:-150,0.75]?} &
\ChemPicture{(-[:-35]?)<[:-60]-[:+0,1.5,,,line width=\boldbondwidth]>[:+60]-[:-145]?} \\
\ce{X}=\ce{Y}=\ce{CH2} & ($\SymGroup{D}{3d}$) & ($\SymGroup{D}{2}$) & ($\SymGroup{C}{2v}$) \\
\bottomrule
\end{tabular}
\vspace{\medskipamount}
}
\caption{\label{fig:Conformations:Six}Базисные конформации шестичленных циклов для 1,4-дизамещённых аналогов циклогексана~\cmpd{Cyclohexane}}
\end{figure}

Для пространственных форм насыщенных шестичленных циклов выделяются три класса так называемых \emph{базисных} конформаций: \tqt{кресло}~(\ConfName{К}), \tqt{твист}~(\ConfName{T}) и \tqt{ванна}~(\ConfName{В})~(рис.~\ref{fig:Conformations:Six}); остальные формы являются промежуточными и могут быть выражены как \emph{выпуклые} комбинации этих трёх. Оптимальной является структура \ConfName{К}; \ConfName{В} обычно является переходным состоянием процесса взаимопревращения двух конформеров со~структурой \ConfName{T}; соответствующий процесс называется \emph{пседовращением} шестичленного цикла. Причиной такой нестабильности \tqt{ванны} признано сильное повышение энергии структур с ординарными связями в заслонённой конформации (так называемое \emph{питцеровское напряжение}). 

\begin{figure}
  \centerfloat{
    \begin{tabular}{cc}
      \multicolumn{2}{c}{\ChemPicture{?[a]<[:-30,1.5]-[:+30,,,,line width=\boldbondwidth](>[:+120]-[:-120](-[:-150]?[a])(-[:-30]-[:+30,1.25]?[b]))-[:-+30,,,,line width=\boldbondwidth]?[b,{<}]}} \\
      \multicolumn{2}{c}{\BB{}~(\TT{})} \\
      \midrule\ChemPicture{?[a]<[:-30,1.25]-[:+30,,,,line width=\boldbondwidth](>[:+120]-[:-120](-[:-150]?[a])(-[:-30]-[:-60]?[b]))-[:-+30,,,,line width=\boldbondwidth]?[b,{<}]} & 
      \ChemPicture{?[a]<[:+60]-[:+30,,,,line width=\boldbondwidth](>[:+120]-[:-120](-[:-150]?[a])(-[:-30]-[:+30,1.25]?[b]))-[:-+30,,,,line width=\boldbondwidth]?[b,{<}]} \\
      \multicolumn{2}{c}{\BC{} \(\equiv\) \CB{}} \\
      \midrule\multicolumn{2}{c}{\ChemPicture{?[a]<[:+60]-[:+30,,,,line width=\boldbondwidth](>[:+120]-[:-120](-[:-150]?[a])(-[:-30]-[:-60]?[b]))-[:-+30,,,,line width=\boldbondwidth]?[b,{<}]} } \\
      \multicolumn{2}{c}{\CC{}} \\
    \end{tabular}
  }
  \vspace{\medskipamount}
  \caption{\label{fig:System331:Conf}Основные конформации бицикло[3.3.1]нонана~\cmpd{Bicycle331}}
\end{figure}

Для насыщенных производных~\cmpd{Bicycle331} выделяются конформации \tqt{двойное кресло}~(\CC{}), \tqt{ванна-кресло}~(\BC{}), \tqt{кресло-ванна} (\CB{}), \tqt{кресло-ванна}~(\CB{}) и~\tqt{двойной ванны} (\BB{})~(рис.~\ref{fig:System331:379XYZ:Conf}). В случае структурной эквивалентности \tqt{крыльев} (например, для исходного углеводорода~\cmpd{Bicycle331}) наблюдается \emph{конформационное вырождение} "--- формы \CB{} и~\BC{} становятся структурно и~энергетически неразличимыми~(рис.~\ref{fig:System331:Conf}).

\begin{figure}
\centerfloat{
\begin{tabular}{c|c}
\multicolumn{2}{c}{
\ChemPicture{X?[a]<[:-30,1.5]-[:+30,,,,line width=\boldbondwidth](>[:+120]Z-[:-120](-[:-150]?[a])(-[:-30]-[:+30,1.25]Y?[b]))-[:-+30,,,,line width=\boldbondwidth]?[b,{<}]} 
}
\\
\multicolumn{2}{c}{\BB{}~(\TT{})} \\
\midrule
\ChemPicture{X?[a]<[:-30,1.25]-[:+30,,,,line width=\boldbondwidth](>[:+120]Z-[:-120](-[:-150]?[a])(-[:-30]-[:-60]Y?[b]))-[:-+30,,,,line width=\boldbondwidth]?[b,{<}]
}
& 
\ChemPicture{X?[a]<[:+60]-[:+30,,,,line width=\boldbondwidth](>[:+120]Z-[:-120](-[:-150]?[a])(-[:-30]-[:+30,1.25]Y?[b]))-[:-+30,,,,line width=\boldbondwidth]?[b,{<}]}
\\
\BC{} & \CB{}
\\
\midrule
\multicolumn{2}{c}{ %
\ChemPicture{X?[a]<[:+60]-[:+30,,,,line width=\boldbondwidth](>[:+120]Z-[:-120](-[:-150]?[a])(-[:-30]-[:-60]Y?[b]))-[:-+30,,,,line width=\boldbondwidth]?[b,{<}]} % 
}
\\
\multicolumn{2}{c}{\CC{}}
\\
\end{tabular}
}
\vspace{\medskipamount}
\caption{\label{fig:System331:379XYZ:Conf}Основные конформации 3X,~7Y,~9Z-аналогов бицикло[3.3.1]нонана}
\end{figure}

Форма \BB{} у~производных бицикло[3.3.1]нонана из-за питцеровских напряжений встречается очень редко (характерна для сильно затруднённых и напряжённых молекул), и в этих случаях обычно наблюдается экспериментально в виде формы «двойного твиста» (\TT{}).

Торсионные напряжения также являются фактором устойчивости и для формы \CB{}, содержащей подструктуру \tqt{ванны}, в которой, однако, заслонённые конформации связей незначительно отклоняются от идеальной формы \ConfName{В} с эндоциклическими двугранными углами $\tau_{9123}\simeq 0$ и $\tau_{9543}\simeq 0$. Часто эти отклонения симметричны, т. е., $\tau_{9123} = - \tau_{9543}$. Всё это является, в частности, показателем стереохимической (конформационной) жёсткости формы \tqt{кресла} насыщенного шестичленного цикла. Такая жёсткость образуется из сочетания низкой энергии и высоких собственных частот колебания структуры. Жёсткие конформации одной из подструктур молекулы могут стабилизировать менее энергетически выгодные формы других как термодинамически, так и кинетически.

\begin{figure}
\centerfloat{
}
\caption{\label{fig:Interactions:37}3,7-взаимодействия в молекулах аналогов бицикло[3.3.1]нонана}
\end{figure}

Конформация \CC{} состоит из кресловидных шестичленных циклов, почти свободных от питцеровского напряжения. Основным фактором устойчивости для этой конформации оказывается невалентное взаимодействие между \tqt{крыльями} бициклической системы. Особенно значима в~этом смысле роль положений 3 и 7~(рис.~\ref{fig:Interactions:37}). Дисперсионное или электростатическое отталкивание между этими положениями или \emph{эндо}-заместителями в~них приводит к~повышению энергии и~дестабилизации \CC{} относительно других форм, тогда как притяжение стабилизирует соответствующую структуру.

\begin{figure}
\centerfloat{}
\caption{\label{fig:Interactions:2468}2,4~(6,8)-взаимодействия в молекулах аналогов бицикло[3.3.1]нонана}
\end{figure}

Дополнительным фактором, влияющим на относительные энергии конформеров \BC{}/\CB{} и~\CC{} являются взаимодействия положений 2 и~4 (6 и~8) внутри \tqt{крыльев}~(рис.~\ref{fig:Interactions:2468}).). Они частично сходны с 1,3-диаксиальными взаимодействиями в соответствующих конформерах 1,3-\emph{цис}-дизамещённых аналогов циклогексана. Отталкивание приводит к уплощению циклов и сближению двух упомянутых форм.

Максимальная симметрия бициклического скелета~\cmpd{Bicycle331}, порождённая из~симметрической группы $\SymGroup{S}{9}$ как группа автоморфизмов соответствующего молекулярного графа, порождается системой орбит (эквивалентных вершин) вида $\left\langle(2\,4\,6\,8)(1\,5)(3\,7)(9)\right\rangle\subset\SymGroup{S}{9}$. В эвклидовом трёхмерном (\tqt{мировом}) пространстве $\SymGroup{E}{3}\left(\AGroup{R}\right)\simeq\AGroup{R}^3$ эта группа (изоморфно реализуется) представлена (?)  в~виде точечной группы симметрии \(\SymGroup{C}{2v}\) для~\CC{} и~\BB{}. Для других конформаций симметрия снижается до~подгрупп \(\SymGroup{C}{s}\) у~\BC{}/\CB{} и~$\SymGroup{C}{2}$ у~\TT{}.

Конформационное поведение производных бицикло[3.3.1]нонана в общем случае определяется взаимопревращением конформеров, в которых бициклический скелет приобретает либо конформации \CC{}, либо \BC{}/\CB{}.


% !TeX spellcheck = ru_RU
% !TeX encoding = UTF-8
\chapter{Исследование напряжений в производных бицикло[3.3.1]нонана}\label{ch:Strain:331}

Ранние работы по анализу напряжений в связи со структурой алмаза и адамантана\dots~\cite{Mohr:1918}

\paragraph{Реакции разделения связей} Наименьшую погрешность обеспечивают \emph{гипергомодесмотические} реакции, сохраняющие не только состав, как гомодесмические... но и число связей различных типов. В частности, для 3,7,9-замещённых аналогов бицикло[3.3.1]нонана выводится уравнение
\begin{center}
  \chemfig{[:-30]Z*6(-(-[:180]-[:120]X-[:+60]-[:0]?)--Y---)}
  + \chemfig{H_3C-[:-30]X-[:+30]CH_3} + \chemfig{H_3C-[:-30]Y-[:+30]CH_3} + \chemfig{H_3C-[:-30]Z-[:+30]CH_3} + 4\(\cdot\) \chemfig{H_3C-[:-30]-[:+30]CH_3}
  + 4\(\cdot\)\chemfig{H_3C-[:+30](-[:+90]CH_3)-[:-30]CH_3}
  \(\longrightarrow\)
  
  \(\longrightarrow\) 
  2\(\cdot\)\chemfig{H_3C-[:+30]X-[:-30]-[:+30]CH_3} + 2\(\cdot\)\chemfig{H_3C-[:+30]Y-[:-30]-[:+30]CH_3} + 4\(\cdot\)\chemfig{H_3C-[:+30](-[:+90]CH_3)-[:-30]-[:+30]CH_3} + 
  2\(\cdot\)\chemfig{H_3C-[:+30](-[:+90]CH_3)-[:-30]Z-[:+30]CH_3}
\end{center}

Для 1,3,5,7-аналогов исследуемой бициклической системы
\begin{center}
  \chemfig{[:-30]*6(-A(-[:180]-[:120]X-[:+60]-[:0]\phantom{Z}?)--Y--Z-)}
  + \chemfig{H_3C-[:-30]X-[:+30]CH_3} + \chemfig{H_3C-[:-30]Y-[:+30]CH_3} + 5\(\cdot\) \chemfig{H_3C-[:-30]-[:+30]CH_3}
  + 2\(\cdot\)\chemfig{H_3C-[:+30]A(-[:+90]CH_3)-[:-30]CH_3}
  + 2\(\cdot\)\chemfig{H_3C-[:+30]Z(-[:+90]CH_3)-[:-30]CH_3}
  \(\longrightarrow\)
  
  \(\longrightarrow\) 
  2\(\cdot\)\chemfig{H_3C-[:+30]X-[:-30]-[:+30]CH_3} + 2\(\cdot\)\chemfig{H_3C-[:+30]Y-[:-30]-[:+30]CH_3} + 3\(\cdot\)\chemfig{H_3C-[:+30]A(-[:+90]CH_3)-[:-30]-[:+30]CH_3} + 3\(\cdot\)\chemfig{H_3C-[:+30]Z(-[:+90]CH_3)-[:-30]-[:+30]CH_3}
\end{center}

Энергетический эффект этих реакций, взятый с обратным знаком, оценивает общее напряжение молекулы относительно ациклических форм, считающихся свободными от напряжения. 

В частности, для 3,7-аналогов 1-азабицикло[3.3.1]нонана такое разделение связей описывается уравнением
\begin{center}
  \chemfig{[:-30]*6(-N(-[:180]-[:120]X-[:+60]-[:0]?)--Y---)} +
  + \chemfig{H_3C-[:-30]X-[:+30]CH_3} + \chemfig{H_3C-[:-30]Y-[:+30]CH_3}
  + 5\(\cdot\)\chemfig{H_3C-[:-30]-[:+30]CH_3}
  + 2\(\cdot\)\chemfig{H_3C-[:-30](-[:-90]CH_3)-[:+30]CH_3}
  + 2\(\cdot\)\chemfig{H_3C-[:-30]N(-[:-90]CH_3)-[:+30]CH_3}
  \(\longrightarrow\)
  
  \(\longrightarrow\) 
  2\(\cdot\)\chemfig{H_3C-[:+30]X-[:-30]-[:+30]CH_3} + 2\(\cdot\)\chemfig{H_3C-[:+30]Y-[:-30]-[:+30]CH_3} + 3\(\cdot\)\chemfig{H_3C-[:+30](-[:+90]CH_3)-[:-30]-[:+30]CH_3} + 
  3\(\cdot\)\chemfig{H_3C-[:+30]N(-[:+90]CH_3)-[:-30]-[:+30]CH_3}
\end{center}

Для аналогов 1,5-диаазабицикло[3.3.1]нонана
\begin{center}
  \chemfig{[:-30]*6(-N(-[:180]-[:120]X-[:+60]-[:0]\phantom{N}?)--Y--N-)}
  + \chemfig{H_3C-[:-30]X-[:+30]CH_3} + \chemfig{H_3C-[:-30]Y-[:+30]CH_3} + 5\(\cdot\) \chemfig{H_3C-[:-30]-[:+30]CH_3}
  + 2\(\cdot\)\chemfig{H_3C-[:+30](-[:+90]CH_3)-[:-30]CH_3}
  + 2\(\cdot\)\chemfig{H_3C-[:+30]N(-[:+90]CH_3)-[:-30]CH_3}
  \(\longrightarrow\)
  
  \(\longrightarrow\) 
  2\(\cdot\)\chemfig{H_3C-[:+30]X-[:-30]-[:+30]CH_3} + 2\(\cdot\)\chemfig{H_3C-[:+30]Y-[:-30]-[:+30]CH_3} + 
  6\(\cdot\)\chemfig{H_3C-[:+30]N(-[:+90]CH_3)-[:-30]-[:+30]CH_3}
\end{center}

Для более подробного рассмотрения отдельных взаимодействий необходимо комбинировать эти уравнения с другими.

Для расчётов энергий индивидуальных реагентов применяются неэмпирические (\emph{ab~initio}) методы. Они включают на первом этапе оптимизацию геометрии молекул. Далее расчёт собственных частот колебаний для конформеров. 

Часто несколько методов расчёта комбинируются в виде так называемых \emph{аддитивных схем}\dots Полная энергия $\Energytot$ и энергия нулевых колебаний $\EnergyZPE$

\begin{equation}
E_\varpi = \underbrace{{\overbrace{E_0^\infty + E_c^\infty}^{E_\infty}} + E_k}_{\sim\Energytot} + \EnergyZPE
\end{equation}

Ниже в таблице~\ref{tab:Reagents:Opt} на с.~\pageref{tab:Reagents:Opt}--\pageref{tab:Reagents:Opt:Ends} приводятся результаты оптимизации молекул-\tqt{реагентов} в приближении RI-MP2~\cite{MP:1934,RI:MP2:1997} с использованием расширенного корреляционно-согласованного трёхэкспоненциального базисного набора Даннинга (aug-cc-pVTZ).~\cite{Dunning:1989,Peterson:1993,Woon:1993} Если группа симметрии оптимальной структуры не указана явно, она является тривиальной~(\SymGroup{C}{1}). Все полученные в ходе расчёта собственные частоты молекул действительны.

%\caption{\label{tab:Desmotic:Reagents}Результаты расчёта реагентов для реакций разделения связей}

\begin{longtabu} to \textwidth {rr|SS|X[l]}
\multicolumn{5}{p{\textwidth}}{Таблица\label{tab:Reagents:Opt}~\arabic{table}. Строение молекул-\tqt{реагентов} для реакций разделения связей по данным неэмпирических расчётов} \\ %
\toprule\multicolumn{2}{c|}{Структура / симметрия}  & \(E_{2}\),~\si{\hartree} & \(E_{ZPE}\),~\si{\hartree} & Геометрия \endfirsthead
\multicolumn{5}{p{\textwidth}}{\emph{продолжение таблицы}~\arabic{table}} \\
\toprule\multicolumn{2}{c|}{Структура / симметрия} & \(E_{2}\),~\si{\hartree} & \(E_{ZPE}\),~\si{\hartree} & Геометрия \endhead
  \midrule\multicolumn{5}{l}{\emph{продолжение на следующей странице}} \endfoot
  \bottomrule\multicolumn{5}{l}{} \endlastfoot
\midrule\chemfig{H_3C-[:-30]-[:+30]CH_3} / \(\SymGroup{C}{2v}\) &\cmpd{Propane} & \num{-118.859638} & \num{0.104589} &
\ce{C-C}~\SI{1.523}{\angstrom}; $\angle$\ce{CCC}~\ang{112.3} \\
\midrule \chemfig{H_3C-[:-30]O-[:+30]CH_3} / \(\SymGroup{C}{2v}\) &\cmpd{Me2O} & \num{-154.737526} & \num{0.080693} & 
\ce{C-O}~\SI{1.411}{\angstrom};
$\angle$\ce{COC}~\ang{110.5} \\
\midrule \chemfig{H_3C-[:-30]S-[:+30]CH_3} / \(\SymGroup{C}{2v}\) & \cmpd{Me2S} & \num{-477.351949} & \num{0.076641} & 
\ce{C-S}~\SI{1.804}{\angstrom}; 
$\angle$\ce{CSC}~\ang{97.9} \\
%\midrule\cmpd{Dimethylamine} & \chemfig{H_3C-[:-30]N(-[:-90]H)-[:+30]CH_3} & \(\SymGroup{C}{s}\) & \\
\midrule \chemfig{H_3C-[:-30](-[:-90]CH_3)-[:+30]CH_3} / \(\SymGroup{C}{3v}\) & \cmpd{Isobutane} & \num{-158.087267} & \num{0.132478} & \ce{C-C}~\SI{1.524}{\angstrom}; $\angle$\ce{CCC}~\ang{110.7}\\
\midrule\chemfig{H_3C-[:-30]C(=[:-90]O)-[:+30]CH_3} / \(\SymGroup{C}{2v}\) & \cmpd{Acetone} &  \num{-192.789711} & \num{0.084256} & \ce{C=O}~\SI{1.220}{\angstrom}; 
\ce{C-C}~\SI{1.508}{\angstrom};
$\angle$\ce{OCC}~\ang{122.1};
$\angle$\ce{CCC}~\ang{115.8}\\
\midrule \chemfig{H_3C-[:-30]N(-[:-90]CH_3)-[:+30]CH_3} / \(\SymGroup{C}{3v}\) & \cmpd{Trimethylamine} & \num{-174.100139} & \num{0.121385} & 
\ce{C-N}~\SI{1.451}{\angstrom};
$\angle$\ce{CNC}~\ang{110.2} \\
\midrule\chemfig{H_3C-[:-30]-[:+30]-[:-30]CH_3} / \(\SymGroup{C}{2h}\) &\cmpd{Butane} &  \num{-158.084303} & \num{0.133309} & 
\ce{H_3C-CH_2}~\SI{1.524}{\angstrom}; 
\ce{H_2C-CH_2}~\SI{1.523}{\angstrom}; 
$\angle$\ce{CCC}~\ang{112.3} \\
\midrule\chemfig{H_3C-[:-30]O-[:+30]-[:-30]CH_3} / \(\SymGroup{C}{s}\) & \cmpd{MeOEt} &  \num{-193.966394} & \num{0.108895} & 
\ce{H3C-O}~\SI{1.412}{\angstrom};
\ce{H_2C-O}~\SI{1.417}{\angstrom};
\ce{H2C-CH_3}~\SI{1.509}{\angstrom};
$\angle$\ce{COC}~\ang{111.2};
$\angle$\ce{OCC}~\ang{108.1} \\
\midrule\chemfig{H_3C-[:-30]S-[:+30]-[:-30]CH_3} / \(\SymGroup{C}{s}\) & \cmpd{MeSEt} &  \num{-516.576661} & \num{0.105402} &
\ce{H_3C-S}~\SI{1.805}{\angstrom}; 
\ce{H_2C-S}~\SI{1.809}{\angstrom}; 
\ce{H_2C-CH_3}~\SI{1.519}{\angstrom}; 
$\angle$\ce{CSC}~\ang{98.7}; 
$\angle$\ce{SCC}~\ang{109.3} \\
\midrule\chemfig{H_3C-[:-30]C(=[:-90]O)-[:+30]-[:-30]CH_3} / \(\SymGroup{C}{s}\) & \cmpd{Butanone} & \num{-232.014552} & \num{0.113700} & \ce{C=O}~\SI{1.220}{\angstrom};
\ce{H_3C-CO}~\SI{1.509}{\angstrom}; 
\ce{H_2C-CO}~\SI{1.513}{\angstrom}; 
\ce{H_3C-CH_2}~\SI{1.517}{\angstrom};
$\angle\ce{OCC(H_3)}~\ang{121.9}$;
$\angle\ce{OCC(H_2)}~\ang{122.1}$; 
$\angle\ce{CC(O)C}~\ang{116.0}$; 
$\angle\ce{C(O)CC}~\ang{113.4}$;
%$\tau_{\ce{OCCC}}\approx\ang{0}$ 
\\
\midrule\chemfig{H_3C-[:-30](-[:-90]CH_3)-[:+30]-[:-30]CH_3}  & \cmpd{Isopentane} & \num{-197.311142} & \num{0.161457} & \\
\midrule\chemfig{H_3C-[:-30](-[:-90]CH_3)-[:+30]O-[:-30]CH_3} & \cmpd{Me2CHOMe} & \num{-233.195233} & \num{0.137450} &  \\
\midrule\chemfig{H_3C-[:-30](-[:-90]CH_3)-[:+30]S-[:-30]CH_3} & \cmpd{Me2CHSMe} & & & \\
\midrule \chemfig{H_3C-[:-30]N(-[:-90]CH_3)-[:+30]-[:-30]CH_3} & \cmpd{Me2NEt} & & & \\
%\midrule\cmpd{Isopentanone} &\chemfig{H_3C-[:+30](-[:+90]CH_3)-[:-30]C(=[:-90]O)-[:+30]CH_3} & & & \\
\end{longtabu}\label{tab:Reagents:Opt:Ends}

За исключением простейших структур \cmpd{Propane}--\cmpd{Acetone}, молекулы реагентов обладают конформационной подвижностью. Инверсия атома азота триметиламина~\cmpd{Trimethylamine} взаимодействует со внутренним вращением метильных групп. 
%Внутреннее вращение \emph{н}-бутана~\cmpd{Butane} и его аналогов\dots

\pgfplotstablecreatecol[%
create col/assign/.code={%
\getthisrow{Etot}\entryval%
%\ifnum\pgfplotstablerow>0
\pgfmathsetmacro{\entryval}{min(\entryval,\pgfmathaccuma)}%
%\fi
\edef\pgfmathaccuma{\entryval}
\pgfkeyslet{/pgfplots/table/create col/next content}\entryval
}
]{Emin}\ScanButane %
\pgfplotstablecreatecol[%
create col/assign/.code={%
  \getthisrow{Etot}\entryval%
  \getthisrow{Emin}\refval%
  %\ifnum\pgfplotstablerow>0
  \pgfmathsetmacro{\entryval}{\entryval - \refval}%
  \pgfmathsetmacro{\entryval}{\AUtoKJM * \entryval}%
  %\fi
  \edef\pgfmathaccuma{\entryval}
  \pgfkeyslet{/pgfplots/table/create col/next content}\entryval
}
]{Erel}\ScanButane %
\pgfplotstablesave[columns={Tau,Etot,Erel}]{\ScanButane}{Results/ScanButane.save.dat}

\pgfplotstablecreatecol[%
create col/assign/.code={%
  \getthisrow{Etot}\entry%
  \ifnum\pgfplotstablerow>0
  \pgfmathsetmacro{\entry}{min(\entry,\pgfmathaccuma)}%
  \fi
  \edef\pgfmathaccuma{\entry}
  \pgfkeyslet{/pgfplots/table/create col/next content}\entry
}]{Emin}\ScanMeOEt %
\pgfplotstablecreatecol[%
create col/assign/.code={%
  \getthisrow{Etot}\entryval%
  \getthisrow{Emin}\refval%
  %\ifnum\pgfplotstablerow>0
  \pgfmathsetmacro{\entryval}{\entryval - \refval}%
  \pgfmathsetmacro{\entryval}{\AUtoKJM * \entryval}%
  %\fi
  \edef\pgfmathaccuma{\entryval}
  \pgfkeyslet{/pgfplots/table/create col/next content}\entryval
}
]{Erel}\ScanMeOEt %
\pgfplotstablesave[columns={Tau,Etot,Erel}]{\ScanMeOEt}{Results/ScanMeOEt.save.dat}

\pgfplotstablecreatecol[%
create col/assign/.code={%
  \getthisrow{Etot}\entry%
  \ifnum\pgfplotstablerow>0
  \pgfmathsetmacro{\entry}{min(\entry,\pgfmathaccuma)}%
  \fi
  \edef\pgfmathaccuma{\entry}
  \pgfkeyslet{/pgfplots/table/create col/next content}\entry
}]{Emin}\ScanMeSEt %
\pgfplotstablecreatecol[%
create col/assign/.code={%
  \getthisrow{Etot}\entryval%
  \getthisrow{Emin}\refval%
  %\ifnum\pgfplotstablerow>0
  \pgfmathsetmacro{\entryval}{\entryval - \refval}%
  \pgfmathsetmacro{\entryval}{\AUtoKJM * \entryval}%
  %\fi
  \edef\pgfmathaccuma{\entryval}
  \pgfkeyslet{/pgfplots/table/create col/next content}\entryval
}
]{Erel}\ScanMeSEt %
\pgfplotstablesave[columns={Tau,Etot,Erel}]{\ScanMeSEt}{Results/ScanMeSEt.save.dat}

\pgfplotstablecreatecol[%
create col/assign/.code={%
  \getthisrow{Etot}\entry%
  \ifnum\pgfplotstablerow>0
  \pgfmathsetmacro{\entry}{min(\entry,\pgfmathaccuma)}%
  \fi
  \edef\pgfmathaccuma{\entry}
  \pgfkeyslet{/pgfplots/table/create col/next content}\entry
}]{Emin}\ScanMeCOEt %
\pgfplotstablecreatecol[%
create col/assign/.code={%
  \getthisrow{Etot}\entryval%
  \getthisrow{Emin}\refval%
  %\ifnum\pgfplotstablerow>0
  \pgfmathsetmacro{\entryval}{\entryval - \refval}%
  \pgfmathsetmacro{\entryval}{\AUtoKJM * \entryval}%
  %\fi
  \edef\pgfmathaccuma{\entryval}
  \pgfkeyslet{/pgfplots/table/create col/next content}\entryval
}
]{Erel}\ScanMeCOEt %
\pgfplotstablesave[columns={Tau,Etot,Erel}]{\ScanMeCOEt}{Results/ScanMeCOEt.save.dat}

\pgfplotstablecreatecol[%
create col/assign/.code={%
  \getthisrow{Etot}\entry%
  \ifnum\pgfplotstablerow>0
  \pgfmathsetmacro{\entry}{min(\entry,\pgfmathaccuma)}%
  \fi
  \edef\pgfmathaccuma{\entry}
  \pgfkeyslet{/pgfplots/table/create col/next content}\entry
}]{Emin}\ScanIsoPentane %
\pgfplotstablecreatecol[%
create col/assign/.code={%
  \getthisrow{Etot}\entryval%
  \getthisrow{Emin}\refval%
  %\ifnum\pgfplotstablerow>0
  \pgfmathsetmacro{\entryval}{\entryval - \refval}%
  \pgfmathsetmacro{\entryval}{\AUtoKJM * \entryval}%
  %\fi
  \edef\pgfmathaccuma{\entryval}
  \pgfkeyslet{/pgfplots/table/create col/next content}\entryval
}
]{Erel}\ScanIsoPentane %
\pgfplotstablesave[columns={Tau,Etot,Erel}]{\ScanIsoPentane}{Results/ScanIsoPentane.save.dat}

\pgfplotstablecreatecol[%
create col/assign/.code={%
  \getthisrow{Etot}\entry%
  \ifnum\pgfplotstablerow>0
  \pgfmathsetmacro{\entry}{min(\entry,\pgfmathaccuma)}%
  \fi
  \edef\pgfmathaccuma{\entry}
  \pgfkeyslet{/pgfplots/table/create col/next content}\entry
}]{Emin}\ScanIPrOMe %
\pgfplotstablecreatecol[%
create col/assign/.code={%
  \getthisrow{Etot}\entryval%
  \getthisrow{Emin}\refval%
  %\ifnum\pgfplotstablerow>0
  \pgfmathsetmacro{\entryval}{\entryval - \refval}%
  \pgfmathsetmacro{\entryval}{\AUtoKJM * \entryval}%
  %\fi
  \edef\pgfmathaccuma{\entryval}
  \pgfkeyslet{/pgfplots/table/create col/next content}\entryval
}
]{Erel}\ScanIPrOMe %
\pgfplotstablesave[columns={Tau,Etot,Erel}]
{\ScanIPrOMe}{Results/ScanIPrOMe.save.dat}
%
\pgfplotstablecreatecol[%
create col/assign/.code={%
  \getthisrow{Etot}\entry%
  \ifnum\pgfplotstablerow>0
  \pgfmathsetmacro{\entry}{min(\entry,\pgfmathaccuma)}%
  \fi
  \edef\pgfmathaccuma{\entry}
  \pgfkeyslet{/pgfplots/table/create col/next content}\entry
}]{Emin}\ScanIPrSMe %
\pgfplotstablecreatecol[%
create col/assign/.code={%
  \getthisrow{Etot}\entryval%
  \getthisrow{Emin}\refval%
  %\ifnum\pgfplotstablerow>0
  \pgfmathsetmacro{\entryval}{\entryval - \refval}%
  \pgfmathsetmacro{\entryval}{\AUtoKJM * \entryval}%
  %\fi
  \edef\pgfmathaccuma{\entryval}
  \pgfkeyslet{/pgfplots/table/create col/next content}\entryval
}
]{Erel}\ScanIPrSMe %
\pgfplotstablesave[columns={Tau,Etot,Erel}]
{\ScanIPrSMe}{Results/ScanIPrSMe.save.dat}
%
\pgfplotstablecreatecol[%
create col/assign/.code={%
  \getthisrow{Etot}\entry%
  \ifnum\pgfplotstablerow>0
  \pgfmathsetmacro{\entry}{min(\entry,\pgfmathaccuma)}%
  \fi
  \edef\pgfmathaccuma{\entry}
  \pgfkeyslet{/pgfplots/table/create col/next content}\entry
}]{Emin}\ScanMeMeNEt %
\pgfplotstablecreatecol[%
create col/assign/.code={%
  \getthisrow{Etot}\entryval%
  \getthisrow{Emin}\refval%
  %\ifnum\pgfplotstablerow>0
  \pgfmathsetmacro{\entryval}{\entryval - \refval}%
  \pgfmathsetmacro{\entryval}{\AUtoKJM * \entryval}%
  %\fi
  \edef\pgfmathaccuma{\entryval}
  \pgfkeyslet{/pgfplots/table/create col/next content}\entryval
}
]{Erel}\ScanMeMeNEt %
\pgfplotstablesave[columns={Tau,Etot,Erel}]
{\ScanMeMeNEt}{Results/ScanMeMeNEt.save.dat}

\begin{figure}
\caption{Потенциалы внутреннего вращения молекул \cmpd{Butane}--\cmpd{Butanone}\label{fig:Internal:Rotation:Butanes}}

\centerfloat{
\begin{tikzpicture}
  \begin{axis} [%
    %title={\chemfig{H_3C-[:+30]\circ-[:-30]\circ-[:+30]CH_3}},
    %legend pos=outer north east,
    legend style={draw=none},
    xlabel={\chemfig{[,0.75]H_3C-[:+30]\circ-[:-30]\circ-[:+30]CH_3} / $\tau$,~\textdegree},
    xmin=-180, xmax=180,
    xtick={-180,-120,...,180},
    ylabel={$\Delta\Func{E}$,~\si{\kilo\joule\per\mole}},
    ymin=-1, ymax=32]
    \addplot [black,thick,smooth,domain=-180:180] table [x=Tau,y=Erel] 
    {Results/ScanButane.save.dat};\addlegendentry{\chemfig{[,0.675]-[:+30]-[:-30]-[:+30]}~\cmpd{Butane}}
    \addplot [red,thick,smooth,domain=-180:180] table [x=Tau,y=Erel]
    {Results/ScanMeOEt.save.dat};\addlegendentry{\chemfig{[,0.675]-[:+30]-[:-30]O-[:+30]}~\cmpd{MeOEt}}
    \addplot [brown,thick,smooth,domain=-180:180] table [x=Tau,y=Erel] 
    {Results/ScanMeSEt.save.dat};\addlegendentry{\chemfig{[,0.675]-[:+30]-[:-30]S-[:+30]}~\cmpd{MeSEt}}
    \addplot [green,thick,smooth,domain=-180:180] table [x=Tau,y=Erel] 
    {Results/ScanMeCOEt.save.dat}; \addlegendentry{\chemfig{[,0.675]-[:+30]-[:-30]C(=[:-90]O)-[:+30]}~\cmpd{Butanone}}
  \end{axis}
\end{tikzpicture}
}
\end{figure}

\begin{figure}
\caption{Потенциалы внутреннего вращения молекул \cmpd{Isopentane}--\cmpd{Me2NEt}\label{fig:Internal:Rotation:Isopentanes}}

  \centerfloat{
\begin{tikzpicture}
  \begin{axis} [%
    legend style={draw=none},
    xlabel={\chemfig{[,0.75]H_3C-[:+30]\bullet(-[:+90]CH_3)-[:-30]\circ-[:+30]CH_3}/$\tau$,~\textdegree},
    xmin=-180, xmax=180,
    xtick={-180,-120,...,180},
    ylabel={$\Delta\Func{E}$,~\si{\kilo\joule\per\mole}},
    ymin=-1, ymax=32]
    \addplot [black,thick,smooth,domain=-180:180] table [x=Tau,y=Erel] 
    {Results/ScanIsopentane.save.dat};\addlegendentry{\chemfig{[,0.675]-[:+30](-[:+90])-[:-30]-[:+30]}~\cmpd{Isopentane}}
    \addplot [red,thick,smooth,domain=-180:180] table [x=Tau,y=Erel] 
    {Results/ScanIPrOMe.save.dat};\addlegendentry{\chemfig{[,0.675]-[:+30](-[:+90])-[:-30]O-[:+30]}~\cmpd{Me2CHOMe}}
    \addplot [brown,thick,smooth,domain=-180:180] table [x=Tau,y=Erel] 
    {Results/ScanIPrSMe.save.dat};\addlegendentry{\chemfig{[,0.675]-[:+30](-[:+90])-[:-30]S-[:+30]}~\cmpd{Me2CHSMe}}
    \addplot [blue,thick,smooth,domain=-180:180] table [x=Tau,y=Erel] 
    {Results/ScanMeMeNEt.save.dat};\addlegendentry{\chemfig{[,0.75]-[:+30]N(-[:+90])-[:-30]-[:+30]}~\cmpd{Me2NEt}}
\end{axis}
\end{tikzpicture}%
}
\end{figure}

\pgfplotstableset{
%columns/TauAll/.style = {column name = $\tau$,~\textdegree},
%columns/EButane/.style = {column name = \cmpd{Butane}},0
%columns/EMeOEt/.style = {column name = $\Delta\Func{E}$~(\cmpd{MeOEt})},
}

\pgfplotstablenew[
create on use/Tau/.style = {create col/copy column from table={\ScanButane}{Tau}},
create on use/EButane/.style = {create col/copy column from table={\ScanButane}{Erel}},
create on use/EMeOEt/.style = {create col/copy column from table={\ScanMeOEt}{Erel}},
create on use/EMeSEt/.style = {create col/copy column from table={\ScanMeSEt}{Erel}},
create on use/EMeCOEt/.style = {create col/copy column from table={\ScanMeCOEt}{Erel}},
create on use/EIsoPentane/.style = {create col/copy column from table={\ScanIsoPentane}{Erel}},
create on use/EIPrOMe/.style = {create col/copy column from table={\ScanIPrOMe}{Erel}},
create on use/EIPrSMe/.style = {create col/copy column from table={\ScanIPrSMe}{Erel}},
create on use/EDMeNEt/.style = {create col/copy column from table={\ScanMeMeNEt}{Erel}},
columns={Tau,EButane,EMeOEt,EMeSEt,EMeCOEt,EIsoPentane,EIPrOMe,EIPrSMe,EDMeNEt}]{25}\ScanReagents

%\begin{table}
%\label{tab:Rot:Energy:Reagents} 
%\caption{Внутреннее вращение молекул~\cmpd{Butane} -- \cmpd{Me2NEt},~\si{\kilo\joule\per\mole} }
%\centering
%\pgfplotstabletypeset[column type=r,fixed, fixed zerofill, precision=1,
%columns/EButane/.style = {column name = \cmpd{Butane}},
%columns/EMeOEt/.style = {column name = \cmpd{MeOEt}},
%columns/EMeSEt/.style = {column name = \cmpd{MeSEt}},
%columns/EMeCOEt/.style = {column name = \cmpd{Butanone}},
%columns/EIsoPentane/.style = {column name = \cmpd{Isopentane}},
%columns/EIPrOMe/.style = {column name = \cmpd{Me2CHOMe}},
%columns/EIPrSMe/.style = {column name = \cmpd{Me2CHSMe}},
%columns/EDMeNEt/.style = {column name = \cmpd{Me2NEt}},
%]\ScanReagents
%\end{table}

Глобально оптимальными, судя по~рис.~\ref{fig:Internal:Rotation:Butanes}--\ref{fig:Internal:Rotation:Isopentanes}, являются \emph{транс}- или \emph{транс,~гош}-конформеры с $\tau\approx+\ang{60}$ и $\tau\approx\ang{180}$. В дальнейшем во всех случаях для расчётов энергетических эффектов реакций разделения связей были использованы оптимизированные значения энергии именно для таких структур.


\section{Напряжения свободных шестичленных циклов}

Реакция разделения связей для 1,3- и 1,4-дизамещённых аналогов формально одинаковы:

\begin{center}
  \begin{equation*} \underbrace{\text{\chemfig{[:-30]X*6(---Y---)}} \qquad \text{\chemfig{[:-30]*6(-X----Y-)}}}_{\text{как бы это всё поизяЧнее сделать?}} \end{equation*}
+ \chemfig{H_3C-[:-30]X-[:+30]CH_3} + \chemfig{H_3C-[:-30]Y-[:+30]CH_3} + 4\(\cdot\) \chemfig{H_3C-[:-30]-[:+30]CH_3} \(\longrightarrow\)
  
  \(\longrightarrow\) 
    2\(\cdot\)\chemfig{H_3C-[:+30]X-[:-30]-[:+30]CH_3} +     2\(\cdot\)\chemfig{H_3C-[:+30]-[:-30]-[:+30]CH_3} + 2\(\cdot\)\chemfig{H_3C-[:+30]Y-[:-30]-[:+30]CH_3}
\end{center}

Реакции \emph{позиционного диспропорционирования} дизамещённых аналогов циклогексана

\begin{center}
  \chemfig{*6(-X---Y--)} + \chemfig{*6(------)} \(\longrightarrow\) \chemfig{*6(-X-----)} + \chemfig{*6(----Y--)}
\end{center}
  
B таблице~\ref{tab:Cycle:Six:Opt} на с.~\pageref{tab:Cycle:Six:Opt}--\pageref{tab:Cycle:Six:Opt:Ends} приводятся результаты оптимизации молекул производных циклогексана, включая эндоциклические торсионные углы $\tau$.

\begin{center}
  \chemfig{*6(------)} \quad\chemfig{*6(-X-----)} \quad\chemfig{*6(---Y--X-)} \quad\chemfig{*6(-X---Y--)} \quad\chemfig{*6(-Z--Y--X-)}
\end{center}

\begin{longtabu} to \textwidth {rc|SS|X[l]}
\multicolumn{5}{p{\textwidth}}{Таблица\label{tab:Cycle:Six:Opt}~\arabic{table}. Строение молекул циклогексана и его аналогов по данным неэмпирических расчётов} \\ %  
\toprule\multicolumn{2}{c|}{Структура}  & \(E_{2}\),~\si{\hartree} & \(E_{ZPE}\),~\si{\hartree} & Геометрия \endfirsthead
  \multicolumn{5}{p{\textwidth}}{\emph{продолжение таблицы}~\thetable:} \\
  \toprule\multicolumn{2}{c|}{Структура} & \(E_{2}\),~\si{\hartree} & \(E_{ZPE}\),~\si{\hartree} & Геометрия\endhead
  \midrule\multicolumn{5}{l}{\emph{продолжение на следующей странице}}\endfoot
  \bottomrule\multicolumn{5}{l}{}\endlastfoot
  \midrule\cmpd{Cyclohexane} & \chemfig{?-[:+160]-[:-120]<[:+20]-[:-20,,,,line width=\boldbondwidth]>[:+60]?} / \SymGroup{D}{3d} &  \num{-235.344605} & \num{0.171358}& \ce{C-C}~\SI{1.526}{\angstrom}; $\angle\ce{CCC}~\ang{110.7}$; $\tau=\pm\ang{56.9}$ 
  \\ \bottomrule \multicolumn{5}{l}{Монозамещённые (\emph{пентаметиленовые}) производные~(\SymGroup{C}{s}) } \\
\toprule\cmpd{Pyran} & \chemfig{?-[:+160]-[:-120]O<[:+20]-[:-20,,,,line width=\boldbondwidth]>[:+60]?} &  \num{-271.229641} & \num{0.147833} & $\angle\ce{COC}~\ang{110.9}$; $\angle\ce{OCC}~\ang{111.9}$; $\angle\ce{CC_3C}~\ang{109.9}$; $\angle\ce{CC_4C}~\ang{109.1}$; \\
\midrule\cmpd{Thiopyran} & \chemfig{?-[:+160]-[:-120]S<[:+20]-[:-20,,,,line width=\boldbondwidth]>[:+60]?} & \num{-593.841537}  & \num{0.143971} & \chemfig{[-30,1.5]*6(@{6}-@{1}S-@{2}-@{3}-@{4}-@{5}-)} \\
\midrule\cmpd{Cyclohexanone} & \chemfig{?-[:+160]-[:-120]C(=[:-135,0.875]O)<[:+20]-[:-20,,,,line width=\boldbondwidth]>[:+60]?} & & & \ce{C-C} \\
\midrule\cmpd{MeCyclohexaneEq} & \chemfig{?-[:+160]-[:-120](-[:+155]H_3C)(-[:-90,0.75]H)<[:+20]-[:-20,,,,line width=\boldbondwidth]>[:+60]?} & & & \ce{C-C} \\
\midrule\cmpd{NMePiperidineEq} & \chemfig{?-[:+160]-[:-120]N(-[:+155]H_3C)<[:+20]-[:-20,,,,line width=\boldbondwidth]>[:+60]?} & & & \ce{C-C} \\ \bottomrule
\multicolumn{5}{l}{Симметричные 1,3-дизамещённые производные~(\ce{X}=\ce{Y},~\SymGroup{C}{s}) } \\ 
\toprule\cmpd{Dioxane13} & \chemfig{?-[:+160]-[:-120]O<[:+20]-[:-20,,,,line width=\boldbondwidth]O>[:+60]?} & & & \ce{C-C} \\ 
\midrule\cmpd{Dithiane13} & \chemfig{?-[:+160]-[:-120]S<[:+20]-[:-20,,,,line width=\boldbondwidth]S>[:+60]?} & & & \ce{C-C} \\
\midrule\cmpd{N2Me2Diazine13} & \chemfig{?-[:+160]-[:-120]N(-[:+160]H_3C)<[:+20]-[:-20,,,,line width=\boldbondwidth]N(-[:+20]CH_3)>[:+60]?} & & & \ce{C-C} \\
\bottomrule
\multicolumn{5}{l}{Симметричные 1,4-дизамещённые производные~(\ce{X}=\ce{Y},~\SymGroup{C}{2h}) } \\
\toprule\cmpd{Dioxane14} & \chemfig{?-[:+160]-[:-120]O<[:+20]-[:-20,,,,line width=\boldbondwidth]>[:+60]O?} & & & \ce{C-C} \\
\midrule\cmpd{Dithiane14} & \chemfig{?-[:+160]-[:-120]S<[:+20]-[:-20,,,,line width=\boldbondwidth]>[:+60]S?} & & & \ce{C-C} \\
\midrule\cmpd{Cyclohexandione14} & \chemfig{?-[:-120]C(=[:-135,0.875]O)<[:+20]-[:-20,,,,line width=\boldbondwidth]>[:+60]C(=[:+45,0.875]O)-[:-160]?} & & & \ce{C-C} \\
\midrule\cmpd{Me2CyclohexaneEE} & \chemfig{?-[:-120](-[:+155]H_3C)(-[:-90,0.75]H)<[:+20]-[:-20,,,,line width=\boldbondwidth]>[:+60](-[:+90,0.75]H)(-[:-25]CH_3)-[:-160]?} & & & \ce{C-C} \\
\midrule\cmpd{N2Me2PiperazineEE} & \chemfig{?-[:-120]N(-[:+155]H_3C)<[:+20]-[:-20,,,,line width=\boldbondwidth]>[:+60]N(-[:-25]CH_3)-[:-160]?} & & & \ce{C-C} \\
\bottomrule\multicolumn{5}{l}{Несимметричные 1,4-дизамещённые производные~(\ce{X}$\ne$\ce{Y},~\SymGroup{C}{s}) } \\
\toprule\cmpd{Oxathiane4} & \chemfig{?-[:+160]-[:-120]S<[:+20]-[:-20,,,,line width=\boldbondwidth]>[:+60]O?} & & & \ce{C-C} \\
\midrule \cmpd{Pyranone4} & \chemfig{?-[:+160]-[:-120]C(=[:-135,0.875]O)<[:+20]-[:-20,,,,line width=\boldbondwidth]>[:+60]O?} & & & \ce{C-C} \\
\midrule\cmpd{Me4PyraneEq} & \chemfig{?-[:+160]-[:-120](-[:+155]H_3C)(-[:-90,0.75]H)<[:+20]-[:-20,,,,line width=\boldbondwidth]>[:+60]O?} & & & \ce{C-C} \\
\midrule\cmpd{NMeMorpholineEq} & \chemfig{?-[:+160]-[:-120]N(-[:+155]H_3C)<[:+20]-[:-20,,,,line width=\boldbondwidth]>[:+60]O?} & & & \ce{C-C} \\
\midrule \cmpd{Thiopyranone4} & \chemfig{?-[:+160]-[:-120]C(=[:-135,0.875]O)<[:+20]-[:-20,,,,line width=\boldbondwidth]>[:+60]S?} && & \ce{C-C} \\
\midrule\cmpd{Me4ThiopyraneEq} & \chemfig{?-[:+160]-[:-120](-[:+155]H_3C)(-[:-90,0.75]H)<[:+20]-[:-20,,,,line width=\boldbondwidth]>[:+60]S?}  & & & \ce{C-C} \\
\midrule\cmpd{NMeThiomorpholineEq} & \chemfig{?-[:+160]-[:-120]N(-[:+155]H_3C)<[:+20]-[:-20,,,,line width=\boldbondwidth]>[:+60]S?} & & & \ce{C-C} \\
\midrule\cmpd{Me4CyclohexanoneEq} & \chemfig{?-[:-120]C(=[:-135,0.875]O)<[:+20]-[:-20,,,,line width=\boldbondwidth]>[:+60](-[:+90,0.75]H)(-[:-25]CH_3)-[:-160]?} & & & \ce{C-C} \\
\midrule\cmpd{NMePiperidone4Eq} & \chemfig{?-[:-120]C(=[:-135,0.875]O)<[:+20]-[:-20,,,,line width=\boldbondwidth]>[:+60]N(-[:-25]CH_3)-[:-160]?} & & & \ce{C-C} \\
\midrule\cmpd{4MePiperidineEE} & \chemfig{?-[:-120]N(-[:+155]H_3C)<[:+20]-[:-20,,,,line width=\boldbondwidth]>[:+60](-[:+90,0.75]H)(-[:-25]CH_3)-[:-160]?} & & & \ce{C-C} \\
\bottomrule\multicolumn{5}{l}{Симметричные 1,3,5-тризамещённые производные~(\ce{X}=\ce{Y}=\ce{Z},~\SymGroup{C}{3v}) } \\
\toprule\cmpd{STrioxane} & \chemfig{O?-[:+160]-[:-120]O<[:+20]-[:-20,,,,line width=\boldbondwidth]O>[:+60]?} & \num{} & & \ce{C-O} \\
\midrule\cmpd{STrithiane} & \chemfig{S?-[:+160]-[:-120]S<[:+20]-[:-20,,,,line width=\boldbondwidth]S>[:+60]?} & & & \ce{C-S} \\
\bottomrule\multicolumn{5}{l}{Несимметричные 1,3,5-тризамещённые производные~(\ce{X}=\ce{Y},~\SymGroup{C}{s}) } \\
\toprule\cmpd{Thiodioxane35} & \chemfig{S?-[:+160]-[:-120]O<[:+20]-[:-20,,,,line width=\boldbondwidth]O>[:+60]?} & & & \ce{C-C} \\ 
\midrule\cmpd{Oxa5dithiane13} & \chemfig{O?-[:+160]-[:-120]S<[:+20]-[:-20,,,,line width=\boldbondwidth]S>[:+60]?} & & & \ce{C-C} \\
\midrule\cmpd{N2Me2Oxa5diazine13} & \chemfig{O?-[:+160]-[:-120]N(-[:+160]H_3C)<[:+20]-[:-20,,,,line width=\boldbondwidth]N(-[:+20]CH_3)>[:+60]?} & & & \ce{C-C} \\
\midrule\cmpd{N2Me2Thia5diazine13} & \chemfig{S?-[:+160]-[:-120]N(-[:+160]H_3C)<[:+20]-[:-20,,,,line width=\boldbondwidth]N(-[:+20]CH_3)>[:+60]?} & & & \ce{C-C} \\
\bottomrule
\end{longtabu}\label{tab:Cycle:Six:Opt:Ends}

Длины связей~(\AA) и эндоциклические двугранные углы

\vspace{\medskipamount}
\chemfig{[,1.75]*6(@{6}-@{1}-@{2}-@{3}-@{4}-@{5}-)} %
\namebond{3}{4}{\small\ang{56.9}}  %
\namebond{1}{2}{\small\num{1.526}} %
\quad
\chemfig{[,1.75]*6(@{6}-@{1}O-@{2}-@{3}-@{4}-@{5}-)} %
\namebond{1}{2}{\small\ang{60.5}}  %
\namebond{2}{3}{\small\ang{-57.4}} %
\namebond{3}{4}{\small\ang{53.4}}  %
\namebond{5}{4}{\small\num{1.526}} % 
\namebond{6}{5}{\small\num{1.520}} % 
\namebond{6}{1}{\small\num{1.422}} %
\quad 
\chemfig{[,1.75]*6(@{6}-@{1}S-@{2}-@{3}-@{4}-@{5}-)}
\namebond{1}{2}{\small\ang{56.3}} 
\namebond{2}{3}{\small\ang{-62.1}}
\namebond{3}{4}{\small\ang{59.6}}
\namebond{5}{4}{\small\num{1.525}}
\namebond{6}{5}{\small\num{1.521}}
\namebond{6}{1}{\small\num{1.814}}
\quad

%\begin{longtabu} to \textwidth {X} \end{longtabu}

\vspace{\medskipamount}
Энергии напряжения... $\EStrain$


%% методические главы
\chapter{Оформление различных элементов}\label{ch:ch1}

\section{Форматирование текста}\label{sec:ch1/sec1}

Мы можем сделать \textbf{жирный текст} и \textit{курсив}.

\section{Ссылки}\label{sec:ch1/sec2}

Сошлёмся на библиографию.
Одна ссылка: \cite[с.~54]{Sokolov}\cite[с.~36]{Gaidaenko}.
Две ссылки: \cite{Sokolov,Gaidaenko}.
Ссылка на собственные работы: \cite{vakbib1, confbib2}.
Много ссылок: %\cite[с.~54]{Lermontov,Management,Borozda} % такой «фокус»
%вызывает biblatex warning относительно опции sortcites, потому что неясно, к
%какому источнику относится уточнение о страницах, а bibtex об этой проблеме
%даже не предупреждает
\cite{Lermontov, Management, Borozda, Marketing, Constitution, FamilyCode,
Gost.7.0.53, Razumovski, Lagkueva, Pokrovski, Methodology, Nasirova, Berestova,
Kriger}%
\ifnumequal{\value{bibliosel}}{0}{% Примеры для bibtex8
    \cite{Sirotko, Lukina, Encyclopedia}%
}{% Примеры для biblatex через движок biber
    \cite{Sirotko2, Lukina2, Encyclopedia2}%
}%
.
И~ещё немного ссылок:~\cite{Article,Book,Booklet,Conference,Inbook,Incollection,Manual,Mastersthesis,
Misc,Phdthesis,Proceedings,Techreport,Unpublished}
% Следует обратить внимание, что пробел после запятой внутри \cite{}
% обрабатывается ожидаемо, а пробел перед запятой, может вызывать проблемы при
% обработке ссылок.
\cite{medvedev2006jelektronnye, CEAT:CEAT581, doi:10.1080/01932691.2010.513279,
Gosele1999161,Li2007StressAnalysis, Shoji199895, test:eisner-sample,
test:eisner-sample-shorted, AB_patent_Pomerantz_1968, iofis_patent1960}
\ifnumequal{\value{bibliosel}}{0}{% Примеры для bibtex8
}{% Примеры для biblatex через движок biber
    \cite{patent2h, patent3h, patent2}%
}%
.

\ifnumequal{\value{bibliosel}}{0}{% Примеры для bibtex8
Попытка реализовать несколько ссылок на конкретные страницы
для \texttt{bibtex} реализации библиографии:
[\citenum{Sokolov}, с.~54; \citenum{Gaidaenko}, с.~36].
}{% Примеры для biblatex через движок biber
Несколько источников (мультицитата):
% Тут специально написано по-разному тире, для демонстрации, что
% применение специальных тире в настоящий момент в biblatex приводит к непоказу
% "с.".
\cites[vii--x, 5, 7]{Sokolov}[v"--~x, 25, 526]{Gaidaenko}[vii--x, 5, 7]{Techreport},
работает только в \texttt{biblatex} реализации библиографии.
}%

Ссылки на собственные работы:~\cite{vakbib1, confbib1}

Сошлёмся на приложения: Приложение~\ref{app:A}, Приложение~\ref{app:B2}.

Сошлёмся на формулу: формула~\eqref{eq:equation1}.

Сошлёмся на изображение: рисунок~\ref{fig:knuth}.

Стандартной практикой является добавление к ссылкам префикса, характеризующего тип элемента.
Это не является строгим требованием, но~позволяет лучше ориентироваться в документах большого размера.
Например, для ссылок на рисунки используется префикс \textit{fig},
для ссылки на~таблицу "--- \textit{tab}.

В таблице~\ref{tab:tab_pref} приложения~\ref{app:B4} приведён список рекомендуемых
к использованию стандартных префиксов.

\section{Формулы}\label{sec:ch1/sec3}

Благодаря пакету \textit{icomma}, \LaTeX~одинаково хорошо воспринимает
в~качестве десятичного разделителя и запятую (\(3,1415\)), и точку (\(3.1415\)).

\subsection{Ненумерованные одиночные формулы}\label{subsec:ch1/sec3/sub1}

Вот так может выглядеть формула, которую необходимо вставить в~строку
по~тексту: \(x \approx \sin x\) при \(x \to 0\).

А вот так выглядит ненумерованная отдельностоящая формула c подстрочными
и надстрочными индексами:
\[
(x_1+x_2)^2 = x_1^2 + 2 x_1 x_2 + x_2^2
\]

Формула с неопределенным интегралом:
\[
\int f(\alpha+x)=\sum\beta
\]

При использовании дробей формулы могут получаться очень высокие:
\[
  \frac{1}{\sqrt{2}+
  \displaystyle\frac{1}{\sqrt{2}+
  \displaystyle\frac{1}{\sqrt{2}+\cdots}}}
\]

В формулах можно использовать греческие буквы:
%Все \original... команды заранее, ради этого примера, определены в Dissertation\userstyles.tex
\[
\alpha\beta\gamma\delta\originalepsilon\epsilon\zeta\eta\theta%
\vartheta\iota\kappa\varkappa\lambda\mu\nu\xi\pi\varpi\rho\varrho%
\sigma\varsigma\tau\upsilon\originalphi\phi\chi\psi\omega\Gamma\Delta%
\Theta\Lambda\Xi\Pi\Sigma\Upsilon\Phi\Psi\Omega
\]
\[%https://texfaq.org/FAQ-boldgreek
\boldsymbol{\alpha\beta\gamma\delta\originalepsilon\epsilon\zeta\eta%
\theta\vartheta\iota\kappa\varkappa\lambda\mu\nu\xi\pi\varpi\rho%
\varrho\sigma\varsigma\tau\upsilon\originalphi\phi\chi\psi\omega\Gamma%
\Delta\Theta\Lambda\Xi\Pi\Sigma\Upsilon\Phi\Psi\Omega}
\]

Для добавления формул можно использовать пары \verb+$+\dots\verb+$+ и \verb+$$+\dots\verb+$$+,
но~они считаются устаревшими.
Лучше использовать их функциональные аналоги \verb+\(+\dots\verb+\)+ и \verb+\[+\dots\verb+\]+.

\subsection{Ненумерованные многострочные формулы}\label{subsec:ch1/sec3/sub2}

Вот так можно написать две формулы, не нумеруя их, чтобы знаки <<равно>> были
строго друг под другом:
\begin{align}
  f_W & =  \min \left( 1, \max \left( 0, \frac{W_{soil} / W_{max}}{W_{crit}} \right)  \right), \nonumber \\
  f_T & =  \min \left( 1, \max \left( 0, \frac{T_s / T_{melt}}{T_{crit}} \right)  \right), \nonumber
\end{align}

Выровнять систему ещё и по переменной \( x \) можно, используя окружение
\verb|alignedat| из пакета \verb|amsmath|. Вот так:
\[
    |x| = \left\{
    \begin{alignedat}{2}
        &&x, \quad &\text{eсли } x\geqslant 0 \\
        &-&x, \quad & \text{eсли } x<0
    \end{alignedat}
    \right.
\]
Здесь первый амперсанд (в исходном \LaTeX\ описании формулы) означает
выравнивание по~левому краю, второй "--- по~\( x \), а~третий "--- по~слову
<<если>>. Команда \verb|\quad| делает большой горизонтальный пробел.

Ещё вариант:
\[
    |x|=
    \begin{cases}
    \phantom{-}x, \text{если } x \geqslant 0 \\
    -x, \text{если } x<0
    \end{cases}
\]

Кроме того, для  нумерованных формул \verb|alignedat| делает вертикальное
выравнивание номера формулы по центру формулы. Например, выравнивание
компонент вектора:
\begin{equation}
\label{eq:2p3}
\begin{alignedat}{2}
{\mathbf{N}}_{o1n}^{(j)} = \,{\sin} \phi\,n\!\left(n+1\right)
         {\sin}\theta\,
         \pi_n\!\left({\cos} \theta\right)
         \frac{
               z_n^{(j)}\!\left( \rho \right)
              }{\rho}\,
           &{\boldsymbol{\hat{\mathrm e}}}_{r}\,+   \\
+\,
{\sin} \phi\,
         \tau_n\!\left({\cos} \theta\right)
         \frac{
            \left[\rho z_n^{(j)}\!\left( \rho \right)\right]^{\prime}
              }{\rho}\,
            &{\boldsymbol{\hat{\mathrm e}}}_{\theta}\,+   \\
+\,
{\cos} \phi\,
         \pi_n\!\left({\cos} \theta\right)
         \frac{
            \left[\rho z_n^{(j)}\!\left( \rho \right)\right]^{\prime}
              }{\rho}\,
            &{\boldsymbol{\hat{\mathrm e}}}_{\phi}\:.
\end{alignedat}
\end{equation}

Ещё об отступах. Иногда для лучшей <<читаемости>> формул полезно
немного исправить стандартные интервалы \LaTeX\ с учётом логической
структуры самой формулы. Например в формуле~\ref{eq:2p3} добавлен
небольшой отступ \verb+\,+ между основными сомножителями, ниже
результат применения всех вариантов отступа:
\begin{align*}
\backslash! &\quad f(x) = x^2\! +3x\! +2 \\
  \mbox{по-умолчанию} &\quad f(x) = x^2+3x+2 \\
\backslash, &\quad f(x) = x^2\, +3x\, +2 \\
\backslash{:} &\quad f(x) = x^2\: +3x\: +2 \\
\backslash; &\quad f(x) = x^2\; +3x\; +2 \\
\backslash \mbox{space} &\quad f(x) = x^2\ +3x\ +2 \\
\backslash \mbox{quad} &\quad f(x) = x^2\quad +3x\quad +2 \\
\backslash \mbox{qquad} &\quad f(x) = x^2\qquad +3x\qquad +2
\end{align*}

Можно использовать разные математические алфавиты:
\begin{align}
\mathcal{ABCDEFGHIJKLMNOPQRSTUVWXYZ} \nonumber \\
\mathfrak{ABCDEFGHIJKLMNOPQRSTUVWXYZ} \nonumber \\
\mathbb{ABCDEFGHIJKLMNOPQRSTUVWXYZ} \nonumber
\end{align}

Посмотрим на систему уравнений на примере аттрактора Лоренца:

\[
\left\{
  \begin{array}{rl}
    \dot x = & \sigma (y-x) \\
    \dot y = & x (r - z) - y \\
    \dot z = & xy - bz
  \end{array}
\right.
\]

А для вёрстки матриц удобно использовать многоточия:
\[
\left(
  \begin{array}{ccc}
    a_{11} & \ldots & a_{1n} \\
    \vdots & \ddots & \vdots \\
    a_{n1} & \ldots & a_{nn} \\
  \end{array}
\right)
\]

\subsection{Нумерованные формулы}\label{subsec:ch1/sec3/sub3}

А вот так пишется нумерованная формула:
\begin{equation}
  \label{eq:equation1}
  e = \lim_{n \to \infty} \left( 1+\frac{1}{n} \right) ^n
\end{equation}

Нумерованных формул может быть несколько:
\begin{equation}
  \label{eq:equation2}
  \lim_{n \to \infty} \sum_{k=1}^n \frac{1}{k^2} = \frac{\pi^2}{6}
\end{equation}

Впоследствии на формулы~\eqref{eq:equation1} и~\eqref{eq:equation2} можно ссылаться.

Сделать так, чтобы номер формулы стоял напротив средней строки, можно,
используя окружение \verb|multlined| (пакет \verb|mathtools|) вместо
\verb|multline| внутри окружения \verb|equation|. Вот так:
\begin{equation} % \tag{S} % tag - вписывает свой текст
  \label{eq:equation3}
    \begin{multlined}
        1+ 2+3+4+5+6+7+\dots + \\
        + 50+51+52+53+54+55+56+57 + \dots + \\
        + 96+97+98+99+100=5050
    \end{multlined}
\end{equation}

Используя команду \verb|\eqrefs|, можно
красиво ссылаться сразу на несколько формул
\eqrefs{eq:equation1, eq:equation3, eq:equation2}, даже перепутав
порядок ссылок \verb|\eqrefs{eq1, eq3, eq2}|.
Аналогично, для ссылок на несколько рисунков, таблиц и т.д. \refs{sec:ch1/sec1, sec:ch1/sec2, sec:ch1/sec3} можно использовать команду \verb|\refs|.
Обе эти команды определены в файле \verb|common/packages.tex|.

Уравнения~\eqrefs{eq:subeq_1,eq:subeq_2} демонстрируют возможности
окружения \verb|\subequations|.
\begin{subequations}
    \label{eq:subeq_1}
    \begin{gather}
        y = x^2 + 1 \label{eq:subeq_1-1} \\
        y = 2 x^2 - x + 1 \label{eq:subeq_1-2}
    \end{gather}
\end{subequations}
Ссылки на отдельные уравнения~\eqrefs{eq:subeq_1-1,eq:subeq_1-2,eq:subeq_2-1}.
\begin{subequations}
    \label{eq:subeq_2}
    \begin{align}
        y &= x^3 + x^2 + x + 1 \label{eq:subeq_2-1} \\
        y &= x^2
    \end{align}
\end{subequations}

\subsection{Форматирование чисел и размерностей величин}\label{sec:units}

Числа форматируются при помощи команды \verb|\num|:
\num{5,3};
\num{2,3e8};
\num{12345,67890};
\num{2,6 d4};
\num{1+-2i};
\num{.3e45};
\num[exponent-base=2]{5 e64};
\num[exponent-base=2,exponent-to-prefix]{5 e64};
\num{1.654 x 2.34 x 3.430}
\num{1 2 x 3 / 4}.
Для написания последовательности чисел можно использовать команды \verb|\numlist| и \verb|\numrange|:
\numlist{10;30;50;70}; \numrange{10}{30}.
Значения углов можно форматировать при помощи команды \verb|\ang|:
\ang{2.67};
\ang{30,3};
\ang{-1;;};
\ang{;-2;};
\ang{;;-3};
\ang{300;10;1}.

Обратите внимание, что ГОСТ запрещает использование знака <<->> для обозначения отрицательных чисел
за исключением формул, таблиц и рисунков.
Вместо него следует использовать слово <<минус>>.

Размерности можно записывать при помощи команд \verb|\si| и \verb|\SI|:
\si{\farad\squared\lumen\candela};
\si{\joule\per\mole\per\kelvin};
\si[per-mode = symbol-or-fraction]{\joule\per\mole\per\kelvin};
\si{\metre\per\second\squared};
\SI{0.10(5)}{\neper};
\SI{1.2-3i e5}{\joule\per\mole\per\kelvin};
\SIlist{1;2;3;4}{\tesla};
\SIrange{50}{100}{\volt}.
Список единиц измерений приведён в таблицах~\refs{tab:unit:base,
tab:unit:derived,tab:unit:accepted,tab:unit:physical,tab:unit:other}.
Приставки единиц приведены в~таблице~\ref{tab:unit:prefix}.

С дополнительными опциями форматирования можно ознакомиться в~описании пакета \texttt{siunitx};
изменить или добавить единицы измерений можно в~файле \texttt{siunitx.cfg}.

\begin{table}
    \centering
    \captionstyle{\centering}
    \caption{Основные величины СИ}\label{tab:unit:base}
    \begin{tabular}{llc}
        \toprule
        Название  & Команда                & Символ         \\
        \midrule
        Ампер     & \verb|\ampere| & \si{\ampere}   \\
        Кандела   & \verb|\candela| & \si{\candela}  \\
        Кельвин   & \verb|\kelvin| & \si{\kelvin}   \\
        Килограмм & \verb|\kilogram| & \si{\kilogram} \\
        Метр      & \verb|\metre| & \si{\metre}    \\
        Моль      & \verb|\mole| & \si{\mole}     \\
        Секунда   & \verb|\second| & \si{\second}   \\
        \bottomrule
    \end{tabular}
\end{table}

\begin{table}
  \small
  \centering
  \begin{threeparttable}% выравнивание подписи по границам таблицы
    \caption{Производные единицы СИ}\label{tab:unit:derived}
    \begin{tabular}{llc|llc}
        \toprule
        Название       & Команда                 & Символ              & Название & Команда & Символ \\
        \midrule
        Беккерель      & \verb|\becquerel|  & \si{\becquerel}     &
        Ньютон         & \verb|\newton|  & \si{\newton}                                      \\
        Градус Цельсия & \verb|\degreeCelsius| & \si{\degreeCelsius} &
        Ом             & \verb|\ohm| & \si{\ohm}                                         \\
        Кулон          & \verb|\coulomb| & \si{\coulomb}       &
        Паскаль        & \verb|\pascal| & \si{\pascal}                                      \\
        Фарад          & \verb|\farad| & \si{\farad}         &
        Радиан         & \verb|\radian| & \si{\radian}                                      \\
        Грей           & \verb|\gray| & \si{\gray}          &
        Сименс         & \verb|\siemens| & \si{\siemens}                                     \\
        Герц           & \verb|\hertz| & \si{\hertz}         &
        Зиверт         & \verb|\sievert| & \si{\sievert}                                     \\
        Генри          & \verb|\henry| & \si{\henry}         &
        Стерадиан      & \verb|\steradian| & \si{\steradian}                                   \\
        Джоуль         & \verb|\joule| & \si{\joule}         &
        Тесла          & \verb|\tesla| & \si{\tesla}                                       \\
        Катал          & \verb|\katal| & \si{\katal}         &
        Вольт          & \verb|\volt| & \si{\volt}                                        \\
        Люмен          & \verb|\lumen| & \si{\lumen}         &
        Ватт           & \verb|\watt| & \si{\watt}                                        \\
        Люкс           & \verb|\lux| & \si{\lux}           &
        Вебер          & \verb|\weber| & \si{\weber}                                       \\
        \bottomrule
    \end{tabular}
  \end{threeparttable}
\end{table}

\begin{table}
  \centering
  \begin{threeparttable}% выравнивание подписи по границам таблицы
    \caption{Внесистемные единицы}\label{tab:unit:accepted}

    \begin{tabular}{llc}
        \toprule
        Название        & Команда                 & Символ          \\
        \midrule
        День            & \verb|\day| & \si{\day}       \\
        Градус          & \verb|\degree| & \si{\degree}    \\
        Гектар          & \verb|\hectare| & \si{\hectare}   \\
        Час             & \verb|\hour| & \si{\hour}      \\
        Литр            & \verb|\litre| & \si{\litre}     \\
        Угловая минута  & \verb|\arcminute| & \si{\arcminute} \\
        Угловая секунда & \verb|\arcsecond| & \si{\arcsecond} \\ %
        Минута          & \verb|\minute| & \si{\minute}    \\
        Тонна           & \verb|\tonne| & \si{\tonne}     \\
        \bottomrule
    \end{tabular}
  \end{threeparttable}
\end{table}

\begin{table}
    \centering
    \captionstyle{\centering}
    \caption{Внесистемные единицы, получаемые из эксперимента}\label{tab:unit:physical}
    \begin{tabular}{llc}
        \toprule
        Название                & Команда                 & Символ                 \\
        \midrule
        Астрономическая единица & \verb|\astronomicalunit| & \si{\astronomicalunit} \\
        Атомная единица массы   & \verb|\atomicmassunit| & \si{\atomicmassunit}   \\
        Боровский радиус        & \verb|\bohr| & \si{\bohr}             \\
        Скорость света          & \verb|\clight| & \si{\clight}           \\
        Дальтон                 & \verb|\dalton| & \si{\dalton}           \\
        Масса электрона         & \verb|\electronmass| & \si{\electronmass}     \\
        Электрон(-)вольт          & \verb|\electronvolt| & \si{\electronvolt}     \\
        Элементарный заряд      & \verb|\elementarycharge| & \si{\elementarycharge} \\
        Энергия Хартри          & \verb|\hartree| & \si{\hartree}          \\
        Постоянная Планка       & \verb|\planckbar| & \si{\planckbar}        \\
        \bottomrule
    \end{tabular}
\end{table}

\begin{table}
  \centering
  \begin{threeparttable}% выравнивание подписи по границам таблицы
    \caption{Другие внесистемные единицы}\label{tab:unit:other}
    \begin{tabular}{llc}
        \toprule
        Название                  & Команда                 & Символ             \\
        \midrule
        Ангстрем                  & \verb|\angstrom| & \si{\angstrom}     \\
        Бар                       & \verb|\bar| & \si{\bar}          \\
        Барн                      & \verb|\barn| & \si{\barn}         \\
        Бел                       & \verb|\bel| & \si{\bel}          \\
        Децибел                   & \verb|\decibel| & \si{\decibel}      \\
        Узел                      & \verb|\knot| & \si{\knot}         \\
        Миллиметр ртутного столба & \verb|\mmHg| & \si{\mmHg}         \\
        Морская миля              & \verb|\nauticalmile| & \si{\nauticalmile} \\
        Непер                     & \verb|\neper| & \si{\neper}        \\
        \bottomrule
    \end{tabular}
  \end{threeparttable}
\end{table}

\begin{table}
  \small
  \centering
  \begin{threeparttable}% выравнивание подписи по границам таблицы
    \caption{Приставки СИ}\label{tab:unit:prefix}
    \begin{tabular}{llcc|llcc}
        \toprule
        Приставка & Команда                 & Символ      & Степень &
        Приставка & Команда                 & Символ      & Степень   \\
        \midrule
        Иокто     & \verb|\yocto| & \si{\yocto} & -24     &
        Дека      & \verb|\deca| & \si{\deca}  & 1         \\
        Зепто     & \verb|\zepto| & \si{\zepto} & -21     &
        Гекто     & \verb|\hecto| & \si{\hecto} & 2         \\
        Атто      & \verb|\atto| & \si{\atto}  & -18     &
        Кило      & \verb|\kilo| & \si{\kilo}  & 3         \\
        Фемто     & \verb|\femto| & \si{\femto} & -15     &
        Мега      & \verb|\mega| & \si{\mega}  & 6         \\
        Пико      & \verb|\pico| & \si{\pico}  & -12     &
        Гига      & \verb|\giga| & \si{\giga}  & 9         \\
        Нано      & \verb|\nano| & \si{\nano}  & -9      &
        Терра     & \verb|\tera| & \si{\tera}  & 12        \\
        Микро     & \verb|\micro| & \si{\micro} & -6      &
        Пета      & \verb|\peta| & \si{\peta}  & 15        \\
        Милли     & \verb|\milli| & \si{\milli} & -3      &
        Экса      & \verb|\exa| & \si{\exa}   & 18        \\
        Санти     & \verb|\centi| & \si{\centi} & -2      &
        Зетта     & \verb|\zetta| & \si{\zetta} & 21        \\
        Деци      & \verb|\deci| & \si{\deci}  & -1      &
        Иотта     & \verb|\yotta| & \si{\yotta} & 24        \\
        \bottomrule
    \end{tabular}
  \end{threeparttable}
\end{table}

\subsection{Заголовки с формулами: \texorpdfstring{\(a^2 + b^2 = c^2\)}{%
a\texttwosuperior\ + b\texttwosuperior\ = c\texttwosuperior},
\texorpdfstring{\(\left\vert\textrm{{Im}}\Sigma\left(
\protect\varepsilon\right)\right\vert\approx const\)}{|ImΣ (ε)| ≈ const},
\texorpdfstring{\(\sigma_{xx}^{(1)}\)}{σ\_\{xx\}\textasciicircum\{(1)\}}
}\label{subsec:with_math}

Пакет \texttt{hyperref} берёт текст для закладок в pdf-файле из~аргументов
команд типа \verb|\section|, которые могут содержать математические формулы,
а~также изменения цвета текста или шрифта, которые не отображаются в~закладках.
Чтобы использование формул в заголовках не вызывало в~логе компиляции появление
предупреждений типа <<\texttt{Token not allowed in~a~PDF string
(Unicode):(hyperref) removing...}>>, следует использовать конструкцию
\verb|\texorpdfstring{}{}|, где в~первых фигурных скобках указывается
формула, а~во~вторых "--- запись формулы для закладок.
           % Глава 1
\chapter{Длинное название главы, в которой мы смотрим на~примеры того, как будут верстаться изображения и~списки}\label{ch:ch2}

\section{Одиночное изображение}\label{sec:ch2/sec1}

\begin{figure}[ht]
  \centerfloat{
    \includegraphics[scale=0.27]{latex}
  }
  \caption{TeX.}\label{fig:latex}
\end{figure}

Для выравнивания изображения по-центру используется команда \verb+\centerfloat+, которая является во
многом улучшенной версией встроенной команды \verb+\centering+.

\section{Длинное название параграфа, в котором мы узнаём как сделать две картинки с~общим номером и названием}\label{sec:ch2/sect2}

А это две картинки под общим номером и названием:
\begin{figure}[ht]
  \begin{minipage}[b][][b]{0.49\linewidth}\centering
    \includegraphics[width=0.5\linewidth]{knuth1} \\ а)
  \end{minipage}
  \hfill
  \begin{minipage}[b][][b]{0.49\linewidth}\centering
    \includegraphics[width=0.5\linewidth]{knuth2} \\ б)
  \end{minipage}
  \caption{Очень длинная подпись к изображению,
      на котором представлены две фотографии Дональда Кнута}
  \label{fig:knuth}
\end{figure}

Те~же~две картинки под~общим номером и~названием,
но с автоматизированной нумерацией подрисунков:
\begin{figure}[ht]
    \centerfloat{
        \hfill
        \subbottom[List-of-Figures entry][Первый подрисунок\label{fig:knuth_2-1}]{%
            \includegraphics[width=0.25\linewidth]{knuth1}}
        \hfill
        \subbottom[\label{fig:knuth_2-2}]{%
            \includegraphics[width=0.25\linewidth]{knuth2}}
        \hfill
        \subbottom[Третий подрисунок]{%
            \includegraphics[width=0.3\linewidth]{example-image-c}}
        \hfill
    }
    \legend{Подрисуночный текст, описывающий обозначения, например. Согласно
    ГОСТ 2.105, пункт 4.3.1, располагается перед наименованием рисунка.}
    \caption[Этот текст попадает в названия рисунков в списке рисунков]{Очень
    длинная подпись к второму изображению, на~котором представлены две
    фотографии Дональда Кнута}\label{fig:knuth_2}
\end{figure}

На рисунке~\ref{fig:knuth_2-1} показан Дональд Кнут без головного убора.
На рисунке~\ref{fig:knuth_2}\subcaptionref*{fig:knuth_2-2}
показан Дональд Кнут в головном уборе.

Возможно вставлять векторные картинки, рассчитываемые \LaTeX\ <<на~лету>>
с~их~предварительной компиляцией. Надписи в таких рисунках будут выполнены
тем же~шрифтом, который указан для документа в целом.
На~рисунке~\ref{fig:tikz_example} на~странице~\pageref{fig:tikz_example}
представлен пример схемы, рассчитываемой пакетом \verb|tikz| <<на~лету>>.
Для ускорения компиляции, подобные рисунки могут быть <<кешированы>>, что
определяется настройками в~\verb|common/setup.tex|.
Причём имя предкомпилированного
файла и~папка расположения таких файлов могут быть отдельно заданы,
что удобно, если не~для подготовки диссертации,
то~для подготовки научных публикаций.
\begin{figure}[ht]
    \centerfloat{
        \ifdefmacro{\tikzsetnextfilename}{\tikzsetnextfilename{tikz_example_compiled}}{}% присваиваемое предкомпилированному pdf имя файла (не обязательно)
        \input{Dissertation/images/tikz_scheme.tikz}

    }
    \legend{}
    \caption[Пример \texttt{tikz} схемы]{Пример рисунка, рассчитываемого
        \texttt{tikz}, который может быть предкомпилирован}\label{fig:tikz_example}
\end{figure}

Множество программ имеют либо встроенную возможность экспортировать векторную
графику кодом \verb|tikz|, либо соответствующий пакет расширения.
Например, в GeoGebra есть встроенный экспорт,
для Inkscape есть пакет svg2tikz,
для Python есть пакет matplotlib2tikz,
для R есть пакет tikzdevice.

\section{Пример вёрстки списков}\label{sec:ch2/sec3}

\noindent Нумерованный список:
\begin{enumerate}
  \item Первый пункт.
  \item Второй пункт.
  \item Третий пункт.
\end{enumerate}

\noindent Маркированный список:
\begin{itemize}
  \item Первый пункт.
  \item Второй пункт.
  \item Третий пункт.
\end{itemize}

\noindent Вложенные списки:
\begin{itemize}
  \item Имеется маркированный список.
  \begin{enumerate}
    \item В нём лежит нумерованный список,
    \item в котором
    \begin{itemize}
      \item лежит ещё один маркированный список.
    \end{itemize}
  \end{enumerate}
\end{itemize}

\noindent Нумерованные вложенные списки:
\begin{enumerate}
  \item Первый пункт.
  \item Второй пункт.
  \item Вообще, по ГОСТ 2.105 первый уровень нумерации
  (при необходимости ссылки в тексте документа на одно из перечислений)
  идёт буквами русского или латинского алфавитов,
  а второй "--- цифрами со~скобками.
  Здесь отходим от ГОСТ.
    \begin{enumerate}
      \item в нём лежит нумерованный список,
      \item в котором
        \begin{enumerate}
          \item ещё один нумерованный список,
          \item третий уровень нумерации не нормирован ГОСТ 2.105;
          \item обращаем внимание на строчность букв,
          \item в этом списке
          \begin{itemize}
            \item лежит ещё один маркированный список.
          \end{itemize}
        \end{enumerate}

    \end{enumerate}

  \item Четвёртый пункт.
\end{enumerate}

\section{Традиции русского набора}

Много полезных советов приведено в материале
\tqt{\href{https://kostyrka.ru/main/en/typesetting-and-typography-crash-course-by-kostyrka/}{Краткий курс благородного набора}} (автор А.\:В.~Костырка).
Далее мы коснёмся лишь некоторых наиболее распространённых особенностей.

\subsection{Пробелы}

В~русском наборе принято:
\begin{itemize}
    \item единицы измерения, знак процента отделять пробелами от~числа:
        10~кВт, 15~\% (согласно ГОСТ 8.417, раздел 8);
    \item \(\tg 20\text{\textdegree}\), но: 20~{\textdegree}C
        (согласно ГОСТ 8.417, раздел 8);
    \item знак номера, параграфа отделять от~числа: №~5, \S~8;
    \item стандартные сокращения: т.\:е., и~т.\:д., и~т.\:п.;
    \item неразрывные пробелы в~предложениях.
\end{itemize}

\subsection{Математические знаки и символы}

Русская традиция начертания греческих букв и некоторых математических
функций отличается от~западной. Это исправляется серией
\verb|\renewcommand|.
\begin{itemize}
%Все \original... команды заранее, ради этого примера, определены в Dissertation\userstyles.tex
    \item[До:] \( \originalepsilon \originalge \originalphi\),
    \(\originalphi \originalleq \originalepsilon\),
    \(\originalkappa \in \originalemptyset\),
    \(\originaltan\),
    \(\originalcot\),
    \(\originalcsc\).
    \item[После:] \( \epsilon \ge \phi\),
    \(\phi \leq \epsilon\),
    \(\kappa \in \emptyset\),
    \(\tan\),
    \(\cot\),
    \(\csc\).
\end{itemize}

Кроме того, принято набирать греческие буквы вертикальными, что
решается подключением пакета \verb|upgreek| (см. закомментированный
блок в~\verb|userpackages.tex|) и~аналогичным переопределением в
преамбуле (см.~закомментированный блок в~\verb|userstyles.tex|). В
этом шаблоне такие переопределения уже включены.

Знаки математических операций принято переносить. Пример переноса
в~формуле~\eqref{eq:equation3}.

\subsection{Кавычки}
В английском языке приняты одинарные и двойные кавычки в~виде ‘...’ и~“...”.
В России приняты французские («...») и~немецкие („...“) кавычки (они называются
«ёлочки» и~«лапки», соответственно). ,,Лапки`` обычно используются внутри
<<ёлочек>>, например, <<... наш гордый ,,Варяг``...>>.

Французкие левые и правые кавычки набираются
как лигатуры \verb|<<| и~\verb|>>|, а~немецкие левые
и правые кавычки набираются как лигатуры \verb|,,| и~\verb|‘‘| (\verb|``|).

Вместо лигатур или команд с~активным символом "\ можно использовать команды
\verb|\glqq| и \verb|\grqq| для набора немецких кавычек и команды \verb|\flqq|
и~\verb|\frqq| для набора французских кавычек. Они определены в пакете
\verb|babel|.

\subsection{Тире}
%  babel+pdflatex по умолчанию, в polyglossia надо включать опцией (и перекомпилировать с удалением временных файлов)
Команда \verb|"---| используется для печати тире в тексте. Оно несколько короче
английского длинного тире. Кроме того, команда задаёт небольшую жёсткую отбивку
от слова, стоящего перед тире. При этом, само тире не~отрывается от~слова.
После тире следует такая же отбивка от текста, как и~перед тире. При наборе
текста между словом и командой, за которым она следует, должен стоять пробел.

В составных словах, таких, как <<Закон Менделеева"--~Клапейрона>>, для печати
тире надо использовать команду \verb|"--~|. Она ставит более короткое,
по~сравнению с~английским, тире и позволяет делать переносы во втором слове.
При~наборе текста команда \verb|"--~| не отделяется пробелом от слова,
за~которым она следует (\verb|Менделеева"--~|). Следующее за командой слово
может быть  отделено от~неё пробелом или перенесено на другую строку.

Если прямая речь начинается с~абзаца, то перед началом её печатается тире
командой \verb|"--*|. Она печатает русское тире и жёсткую отбивку нужной
величины перед текстом.

\subsection{Дефисы и переносы слов}
%  babel+pdflatex по умолчанию, в polyglossia надо включать опцией (и перекомпилировать с удалением временных файлов)
Для печати дефиса в~составных словах введены две команды. Команда~\verb|"~|
печатает дефис и~запрещает делать переносы в~самих словах, а~команда \verb|"=|
печатает дефис, оставляя \TeX ’у право делать переносы в~самих словах.

В отличие от команды \verb|\-|, команда \verb|"-| задаёт место в~слове, где
можно делать перенос, не~запрещая переносы и~в~других местах слова.

Команда \verb|""| задаёт место в~слове, где можно делать перенос, причём дефис
при~переносе в~этом месте не~ставится.

Команда \verb|",| вставляет небольшой пробел после инициалов с~правом переноса
в~фамилии.

\section{Текст из панграмм и формул}

Любя, съешь щипцы, "--- вздохнёт мэр, "--- кайф жгуч. Шеф взъярён тчк щипцы
с~эхом гудбай Жюль. Эй, жлоб! Где туз? Прячь юных съёмщиц в~шкаф. Экс-граф?
Плюш изъят. Бьём чуждый цен хвощ! Эх, чужак! Общий съём цен шляп (юфть) "---
вдрызг! Любя, съешь щипцы, "--- вздохнёт мэр, "--- кайф жгуч. Шеф взъярён тчк
щипцы с~эхом гудбай Жюль. Эй, жлоб! Где туз? Прячь юных съёмщиц в~шкаф.
Экс-граф? Плюш изъят. Бьём чуждый цен хвощ! Эх, чужак! Общий съём цен шляп
(юфть) "--- вдрызг! Любя, съешь щипцы, "--- вздохнёт мэр, "--- кайф жгуч. Шеф
взъярён тчк щипцы с~эхом гудбай Жюль. Эй, жлоб! Где туз? Прячь юных съёмщиц
в~шкаф. Экс-граф? Плюш изъят. Бьём чуждый цен хвощ! Эх, чужак! Общий съём цен
шляп (юфть) "--- вдрызг! Любя, съешь щипцы, "--- вздохнёт мэр, "--- кайф жгуч.
Шеф взъярён тчк щипцы с~эхом гудбай Жюль. Эй, жлоб! Где туз? Прячь юных съёмщиц
в~шкаф. Экс-граф? Плюш изъят. Бьём чуждый цен хвощ! Эх, чужак! Общий съём цен
шляп (юфть) "--- вдрызг! Любя, съешь щипцы, "--- вздохнёт мэр, "--- кайф жгуч.
Шеф взъярён тчк щипцы с~эхом гудбай Жюль. Эй, жлоб! Где туз? Прячь юных съёмщиц
в~шкаф. Экс-граф? Плюш изъят. Бьём чуждый цен хвощ! Эх, чужак! Общий съём цен
шляп (юфть) "--- вдрызг! Любя, съешь щипцы, "--- вздохнёт мэр, "--- кайф жгуч.
Шеф взъярён тчк щипцы с~эхом гудбай Жюль. Эй, жлоб! Где туз? Прячь юных съёмщиц
в~шкаф. Экс-граф? Плюш изъят. Бьём чуждый цен хвощ! Эх, чужак! Общий съём цен
шляп (юфть) "--- вдрызг! Любя, съешь щипцы, "--- вздохнёт мэр, "--- кайф жгуч.
Шеф взъярён тчк щипцы с~эхом гудбай Жюль. Эй, жлоб! Где туз? Прячь юных съёмщиц
в~шкаф. Экс-граф? Плюш изъят. Бьём чуждый цен хвощ! Эх, чужак! Общий съём цен
шляп (юфть) "--- вдрызг! Любя, съешь щипцы, "--- вздохнёт мэр, "--- кайф жгуч.
Шеф взъярён тчк щипцы с~эхом гудбай Жюль. Эй, жлоб! Где туз? Прячь юных съёмщиц
в~шкаф. Экс-граф? Плюш изъят. Бьём чуждый цен хвощ! Эх, чужак! Общий съём цен
шляп (юфть) "--- вдрызг! Любя, съешь щипцы, "--- вздохнёт мэр, "--- кайф жгуч.
Шеф взъярён тчк щипцы с~эхом гудбай Жюль. Эй, жлоб! Где туз? Прячь юных съёмщиц
в~шкаф. Экс-граф? Плюш изъят. Бьём чуждый цен хвощ! Эх, чужак! Общий съём цен
шляп (юфть) "--- вдрызг! Любя, съешь щипцы, "--- вздохнёт мэр, "--- кайф жгуч.
Шеф взъярён тчк щипцы с~эхом гудбай Жюль. Эй, жлоб! Где туз? Прячь юных съёмщиц
в~шкаф. Экс-граф? Плюш изъят. Бьём чуждый цен хвощ! Эх, чужак! Общий съём цен
шляп (юфть) "--- вдрызг! Любя, съешь щипцы, "--- вздохнёт мэр, "--- кайф жгуч.
Шеф взъярён тчк щипцы с~эхом гудбай Жюль. Эй, жлоб! Где туз? Прячь юных съёмщиц
в~шкаф. Экс-граф? Плюш изъят. Бьём чуждый цен хвощ! Эх, чужак! Общий съём цен
шляп (юфть) "--- вдрызг!Любя, съешь щипцы, "--- вздохнёт мэр, "--- кайф жгуч.
Шеф взъярён тчк щипцы с~эхом гудбай Жюль. Эй, жлоб! Где туз? Прячь юных съёмщиц
в~шкаф. Экс-граф? Плюш изъят. Бьём чуждый цен хвощ! Эх, чужак! Общий съём цен

Ку кхоро адолэжкэнс волуптариа хаж, вим граэко ыкчпэтында ты. Граэкы жэмпэр
льюкяльиюч квуй ку, аэквюы продыжщэт хаж нэ. Вим ку магна пырикульа, но квюандо
пожйдонёюм про. Квуй ат рыквюы ёнэрмйщ. Выро аккузата вим нэ.
\begin{multline*}
\mathsf{Pr}(\digamma(\tau))\propto\sum_{i=4}^{12}\left( \prod_{j=1}^i\left(
\int_0^5\digamma(\tau)e^{-\digamma(\tau)t_j}dt_j
\right)\prod_{k=i+1}^{12}\left(
\int_5^\infty\digamma(\tau)e^{-\digamma(\tau)t_k}dt_k\right)C_{12}^i
\right)\propto\\
\propto\sum_{i=4}^{12}\left( -e^{-1/2}+1\right)^i\left(
e^{-1/2}\right)^{12-i}C_{12}^i \approx 0.7605,\quad
\forall\tau\neq\overline{\tau}
\end{multline*}
Квуй ыёюз омниюм йн. Экз алёквюам кончюлату квуй, ты альяквюам ёнвидюнт пэр.
Зыд нэ коммодо пробатуж. Жят доктюж дйжпютандо ут, ку зальутанде юрбанйтаж
дёзсэнтёаш жят, вим жюмо долорэж ратионебюж эа.

Ад ентэгры корпора жплэндидэ хаж. Эжт ат факэтэ дычэрунт пэржыкюти. Нэ нам
доминг пэрчёус. Ку квюо ёужто эррэм зючкёпит. Про хабэо альбюкиюс нэ.
\[
        \begin{pmatrix}
                a_{11} & a_{12} & a_{13} \\
                a_{21} & a_{22} & a_{23}
        \end{pmatrix}
\]

\[
        \begin{vmatrix}
                a_{11} & a_{12} & a_{13} \\
                a_{21} & a_{22} & a_{23}
        \end{vmatrix}
\]

\[
        \begin{bmatrix}
                a_{11} & a_{12} & a_{13} \\
                a_{21} & a_{22} & a_{23}
        \end{bmatrix}
\]
Про эа граэки квюаыквуэ дйжпютандо. Ыт вэл тебиквюэ дэфянятйоныс, нам жолюм
квюандо мандамюч эа. Эож пауло лаудым инкедыринт нэ, пэрпэтюа форынчйбюж пэр
эю. Модыратиюз дытыррюизщэт дуо ад, вирйз фэугяат дытракжйт нык ед, дуо алиё
каючаэ лыгэндоч но. Эа мольлиз юрбанйтаж зигнёфэрумквюы эжт.

Про мандамюч кончэтытюр ед. Трётанё прёнкипыз зигнёфэрумквюы вяш ан. Ат хёз
эквюедым щуавятатэ. Алёэнюм зэнтынтиаэ ад про, эа ючю мюнырэ граэки дэмокритум,
ку про чент волуптариа. Ыльит дыкоры аляквюид еюж ыт. Ку рыбюм мюндй ютенам
дуо.
\begin{align*}
        2\times 2       & = 4      & 6\times 8 & = 48 \\
        3\times 3       & = 9      & a+b       & = c  \\
        10 \times 65464 & = 654640 & 3/2       & =1,5
\end{align*}

\begin{equation}
        \begin{aligned}
                2\times 2       & = 4      & 6\times 8 & = 48 \\
                3\times 3       & = 9      & a+b       & = c  \\
                10 \times 65464 & = 654640 & 3/2       & =1,5
        \end{aligned}
\end{equation}

Пэр йн тальэ пожтэа, мыа ед попюльо дэбетиз жкрибэнтур. Йн квуй аппэтырэ
мэнандря, зыд аляквюид хабымуч корпора йн. Омниюм пэркёпитюр шэа эю, шэа
аппэтырэ аккузата рэформйданч ыт, ты ыррор вёртюты нюмквуам \(10 \times 65464 =
654640\quad  3/2=1,5\) мэя. Ипзум эуежмод \(a+b = c\) мальюизчыт ад дуо. Ад
фэюгаят пытынтёюм адвыржаряюм вяш. Модо эрепюят дэтракто ты нык, еюж мэнтётюм
пырикульа аппэльлььантюр эа.

Мэль ты дэлььынётё такематыш. Зэнтынтиаэ конклььюжионэмквуэ ан мэя. Вёжи лебыр
квюаыквуэ квуй нэ, дуо зймюл дэлььиката ку. Ыам ку алиё путынт.

%Большая фигурная скобка только справа
\[\left. %ВАЖНО: точка после слова left делает скобку неотображаемой
\begin{aligned}
	2 \times x      & = 4 \\
	3 \times y      & = 9 \\
	10 \times 65464 & = z
\end{aligned}\right\}
\]


Конвынёры витюпырата но нам, тебиквюэ мэнтётюм позтюлант ед про. Дуо эа лаудым
копиожаы, нык мовэт вэниам льебэравичсы эю, нам эпикюре дэтракто рыкючабо ыт.
Вэрйтюж аккюжамюз ты шэа, дэбетиз форынчйбюж жкряпшэрит ыт прё. Ан еюж тымпор
рыфэррэнтур, ючю дольор котёдиэквюэ йн. Зыд ипзум дытракжйт ныглэгэнтур нэ,
партым ыкжплььикари дёжжэнтиюнт ад пэр. Мэль ты кытэрож молыжтйаы, нам но ыррор
жкрипта аппарэат.

\[ \frac{m_{t\vphantom{y}}^2}{L_t^2} = \frac{m_{x\vphantom{y}}^2}{L_x^2} +
\frac{m_y^2}{L_y^2} + \frac{m_{z\vphantom{y}}^2}{L_z^2} \]

Вэре льаборэж тебиквюэ хаж ут. Ан пауло торквюатоз хаж, нэ пробо фэугяат
такематыш шэа. Мэльёуз пэртинакёа юлламкорпэр прё ад, но мыа рыквюы конкыптам.
Хёз квюот пэртинакёа эи, ельлюд трактатоз пэр ад. Зыд ед анёмал льаборэж
номинави, жят ад конгуы льабятюр. Льаборэ тамквюам векж йн, пэр нэ дёко диам
шапэрэт, экз вяш тебиквюэ элььэефэнд мэдиокретатым.

Нэ про натюм фюйзчыт квюальизквюэ, аэквюы жкаывола мэль ку. Ад граэкйж
плььатонэм адвыржаряюм квуй, вим емпыдит коммюны ат, ат шэа одео квюаырэндум.
Вёртюты ажжынтиор эффикеэнди эож нэ, доминг лаборамюз эи ыам. Чэнзэрет
мныжаркхюм экз эож, ыльит тамквюам факильизиж нык эи. Квуй ан элыктрам
тинкидюнт ентырпрытаряш. Йн янвыняры трактатоз зэнтынтиаэ зыд. Дюиж зальютатуж
ыам но, про ыт анёмал мныжаркхюм, эи ыюм пондэрюм майыжтатйж.
           % Глава 2
\chapter{Вёрстка таблиц}\label{ch:ch3}

\section{Таблица обыкновенная}\label{sec:ch3/sect1}

Так размещается таблица:

\begin{table} [htbp]
    \caption{Название таблицы}\label{tab:Ts0Sib}%

  \centering
  \begin{threeparttable}% выравнивание подписи по границам таблицы
    \begin{tabular}{| p{3cm} || p{3cm} | p{3cm} | p{4cm}l |}
    \hline
    \hline
    Месяц   & \centering \(T_{min}\), К & \centering \(T_{max}\), К &\centering  \((T_{max} - T_{min})\), К & \\
    \hline
    Декабрь &\centering  253.575   &\centering  257.778    &\centering      4.203  &   \\
    Январь  &\centering  262.431   &\centering  263.214    &\centering      0.783  &   \\
    Февраль &\centering  261.184   &\centering  260.381    &\centering     \(-\)0.803  &   \\
    \hline
    \hline
    \end{tabular}
  \end{threeparttable}
\end{table}

\begin{table} [htbp]% Пример записи таблицы с номером, но без отображаемого наименования
  \centering
  \begin{threeparttable}% выравнивание подписи по границам таблицы
    \captiondelim{}% должен стоять до самого пустого caption
    \caption{}%
    \label{tab:test1}%
    \begin{SingleSpace}
      \begin{tabular}{| c | c | c | c |}
	  \hline
	  Оконная функция & \({2N}\)& \({4N}\)& \({8N}\)\\ \hline
	  Прямоугольное   & 8.72  & 8.77  & 8.77  \\ \hline
	  Ханна           & 7.96  & 7.93  & 7.93  \\ \hline
	  Хэмминга        & 8.72  & 8.77  & 8.77  \\ \hline
	  Блэкмана        & 8.72  & 8.77  & 8.77  \\ \hline
      \end{tabular}%
    \end{SingleSpace}
  \end{threeparttable}
\end{table}

Таблица~\ref{tab:test2} "--- пример таблицы, оформленной в~классическом книжном
варианте или~очень близко к~нему. \mbox{ГОСТу} по~сути не~противоречит. Можно
ещё~улучшить представление, с~помощью пакета \verb|siunitx| или~подобного.

\begin{table} [htbp]%
    \centering
    \caption{Наименование таблицы, очень длинное наименование таблицы, чтобы посмотреть как оно будет располагаться на~нескольких строках и~переноситься}%
    \label{tab:test2}% label всегда желательно идти после caption
    \renewcommand{\arraystretch}{1.5}%% Увеличение расстояния между рядами, для улучшения восприятия.
    \begin{SingleSpace}
        \begin{tabular}{@{}@{\extracolsep{20pt}}llll@{}} %Вертикальные полосы не используются принципиально, как и лишние горизонтальные (допускается по ГОСТ 2.105 пункт 4.4.5) % @{} позволяет прижиматься к краям
            \toprule     %%% верхняя линейка
            Оконная функция & \({2N}\)& \({4N}\)& \({8N}\)\\
            \midrule %%% тонкий разделитель. Отделяет названия столбцов. Обязателен по ГОСТ 2.105 пункт 4.4.5
            Прямоугольное   & 8.72  & 8.77  & 8.77  \\
            Ханна           & 7.96  & 7.93  & 7.93  \\
            Хэмминга        & 8.72  & 8.77  & 8.77  \\
            Блэкмана        & 8.72  & 8.77  & 8.77  \\
            \bottomrule %%% нижняя линейка
        \end{tabular}%
    \end{SingleSpace}
\end{table}

\section{Таблица с многострочными ячейками и примечанием}

В таблице~\ref{tab:makecell} приведён пример использования команды
\verb+\multicolumn+ для объединения горизонтальных ячеек таблицы,
и команд пакета \textit{makecell} для добавления разрыва строки внутри ячеек.

\begin{table} [htbp]
  \centering
  \begin{threeparttable}% выравнивание подписи по границам таблицы
    \caption{Пример использования функций пакета \textit{makecell}}%
    \label{tab:makecell}%
    \begin{tabular}{| c | c | c | c |}
	    \hline
	    Колонка 1                                    & Колонка 2        & \thead{Название колонки 3, \\ не помещающееся в одну строку} & Колонка 4 \\ \hline
	    \multicolumn{4}{|c|}{Выравнивание по центру}                                                                                               \\ \hline
	    \multicolumn{2}{|r|}{\makecell{Выравнивание к \\ правому краю}} & \multicolumn{2}{l|}{Выравнивание к левому краю}                         \\ \hline
	    \makecell{В этой ячейке \\ много информации} & 8.72             & 8.55                                                         & 8.44      \\ \cline{3-4}
	    А в этой мало                                & 8.22             & \multicolumn{2}{c|}{5}                                                  \\ \hline
    \end{tabular}%
  \end{threeparttable}
\end{table}

Таблицы~\ref{tab:test3} и~\ref{tab:test4} "--- пример реализации расположения
примечания в~соответствии с ГОСТ 2.105. Каждый вариант со своими достоинствами
и~недостатками. Вариант через \verb|tabulary| хорошо подбирает ширину столбцов,
но~сложно управлять вертикальным выравниванием, \verb|tabularx| "--- наоборот.
\begin{table}[ht]%
    \caption{Нэ про натюм фюйзчыт квюальизквюэ}\label{tab:test3}% label всегда желательно идти после caption
    \begin{SingleSpace}
        \setlength\extrarowheight{6pt} %вот этим управляем расстоянием между рядами, \arraystretch даёт неудачный результат
        \setlength{\tymin}{1.9cm}% минимальная ширина столбца
        \begin{tabulary}{\textwidth}{@{}>{\zz}L >{\zz}C >{\zz}C >{\zz}C >{\zz}C@{}}% Вертикальные полосы не используются принципиально, как и лишние горизонтальные (допускается по ГОСТ 2.105 пункт 4.4.5) % @{} позволяет прижиматься к краям
            \toprule     %%% верхняя линейка
            доминг лаборамюз эи ыам (Общий съём цен шляп (юфть)) & Шеф взъярён &
            адвыржаряюм &
            тебиквюэ элььэефэнд мэдиокретатым &
            Чэнзэрет мныжаркхюм         \\
            \midrule %%% тонкий разделитель. Отделяет названия столбцов. Обязателен по ГОСТ 2.105 пункт 4.4.5
            Эй, жлоб! Где туз? Прячь юных съёмщиц в~шкаф Плюш изъят. Бьём чуждый цен хвощ! &
            \({\approx}\) &
            \({\approx}\) &
            \({\approx}\) &
            \( + \) \\
            Эх, чужак! Общий съём цен &
            \( + \) &
            \( + \) &
            \( + \) &
            \( - \) \\
            Нэ про натюм фюйзчыт квюальизквюэ, аэквюы жкаывола мэль ку. Ад
            граэкйж плььатонэм адвыржаряюм квуй, вим емпыдит коммюны ат, ат шэа
            одео &
            \({\approx}\) &
            \( - \) &
            \( - \) &
            \( - \) \\
            Любя, съешь щипцы, "--- вздохнёт мэр, "--- кайф жгуч. &
            \( - \) &
            \( + \) &
            \( + \) &
            \({\approx}\) \\
            Нэ про натюм фюйзчыт квюальизквюэ, аэквюы жкаывола мэль ку. Ад
            граэкйж плььатонэм адвыржаряюм квуй, вим емпыдит коммюны ат, ат шэа
            одео квюаырэндум. Вёртюты ажжынтиор эффикеэнди эож нэ. &
            \( + \) &
            \( - \) &
            \({\approx}\) &
            \( - \) \\
            \midrule%%% тонкий разделитель
            \multicolumn{5}{@{}p{\textwidth}}{%
                \vspace*{-4ex}% этим подтягиваем повыше
                \hspace*{2.5em}% абзацный отступ - требование ГОСТ 2.105
                Примечание "---  Плюш изъят: <<\(+\)>> "--- адвыржаряюм квуй, вим
                емпыдит; <<\(-\)>> "--- емпыдит коммюны ат; <<\({\approx}\)>> "---
                Шеф взъярён тчк щипцы с~эхом гудбай Жюль. Эй, жлоб! Где туз?
                Прячь юных съёмщиц в~шкаф. Экс-граф?
            }
            \\
            \bottomrule %%% нижняя линейка
        \end{tabulary}%
    \end{SingleSpace}
\end{table}

Если таблица~\ref{tab:test3} не помещается на той же странице, всё
её~содержимое переносится на~следующую, ближайшую, а~этот текст идёт перед ней.
\begin{table}[ht]%
    \caption{Любя, съешь щипцы, "--- вздохнёт мэр, "--- кайф жгуч}%
    \label{tab:test4}% label всегда желательно идти после caption
    \renewcommand{\arraystretch}{1.6}%% Увеличение расстояния между рядами, для улучшения восприятия.
    \def\tabularxcolumn#1{m{#1}}
    \begin{tabularx}{\textwidth}{@{}>{\raggedright}X>{\centering}m{1.9cm} >{\centering}m{1.9cm} >{\centering}m{1.9cm} >{\centering\arraybackslash}m{1.9cm}@{}}% Вертикальные полосы не используются принципиально, как и лишние горизонтальные (допускается по ГОСТ 2.105 пункт 4.4.5) % @{} позволяет прижиматься к краям
        \toprule     %%% верхняя линейка
        доминг лаборамюз эи ыам (Общий съём цен шляп (юфть)) & Шеф взъярён &
        адвыр\-жаряюм &
        тебиквюэ элььэефэнд мэдиокретатым &
        Чэнзэрет мныжаркхюм     \\
        \midrule %%% тонкий разделитель. Отделяет названия столбцов. Обязателен по ГОСТ 2.105 пункт 4.4.5
        Эй, жлоб! Где туз? Прячь юных съёмщиц в~шкаф Плюш изъят.
        Бьём чуждый цен хвощ! &
        \({\approx}\) &
        \({\approx}\) &
        \({\approx}\) &
        \( + \) \\
        Эх, чужак! Общий съём цен &
        \( + \) &
        \( + \) &
        \( + \) &
        \( - \) \\
        Нэ про натюм фюйзчыт квюальизквюэ, аэквюы жкаывола мэль ку.
        Ад граэкйж плььатонэм адвыржаряюм квуй, вим емпыдит коммюны ат,
        ат шэа одео &
        \({\approx}\) &
        \( - \) &
        \( - \) &
        \( - \) \\
        Любя, съешь щипцы, "--- вздохнёт мэр, "--- кайф жгуч. &
        \( - \) &
        \( + \) &
        \( + \) &
        \({\approx}\) \\
        Нэ про натюм фюйзчыт квюальизквюэ, аэквюы жкаывола мэль ку. Ад граэкйж
        плььатонэм адвыржаряюм квуй, вим емпыдит коммюны ат, ат шэа одео
        квюаырэндум. Вёртюты ажжынтиор эффикеэнди эож нэ. &
        \( + \) &
        \( - \) &
        \({\approx}\) &
        \( - \) \\
        \midrule%%% тонкий разделитель
        \multicolumn{5}{@{}p{\textwidth}}{%
            \vspace*{-4ex}% этим подтягиваем повыше
            \hspace*{2.5em}% абзацный отступ - требование ГОСТ 2.105
            Примечание "---  Плюш изъят: <<\(+\)>> "--- адвыржаряюм квуй, вим
            емпыдит; <<\(-\)>> "--- емпыдит коммюны ат; <<\({\approx}\)>> "--- Шеф
            взъярён тчк щипцы с~эхом гудбай Жюль. Эй, жлоб! Где туз? Прячь юных
            съёмщиц в~шкаф. Экс-граф?
        }
        \\
        \bottomrule %%% нижняя линейка
    \end{tabularx}%
\end{table}

\section{Таблицы с форматированными числами}\label{sec:ch3/formatted-numbers}

В таблицах~\refs{tab:S:parse,tab:S:align} представлены примеры использования опции
форматирования чисел \texttt{S}, предоставляемой пакетом \texttt{siunitx}.

\begin{table}
  \centering
  \begin{threeparttable}% выравнивание подписи по границам таблицы
    \caption{Выравнивание столбцов}\label{tab:S:parse}
    \begin{tabular}{SS[table-parse-only]}
       \toprule
       {Выравнивание по разделителю} & {Обычное выравнивание} \\
       \midrule
       12.345                        & 12.345                 \\
       6,78                          & 6,78                   \\
       -88.8(9)                      & -88.8(9)               \\
       4.5e3                         & 4.5e3                  \\
       \bottomrule
    \end{tabular}
  \end{threeparttable}
\end{table}

\begin{table}
  \centering
  \begin{threeparttable}% выравнивание подписи по границам таблицы
    \caption{Выравнивание с использованием опции \texttt{S}}\label{tab:S:align}
    \sisetup{
        table-figures-integer = 2,
        table-figures-decimal = 4
    }
    \begin{tabular}
        {SS[table-number-alignment = center]S[table-number-alignment = left]S[table-number-alignment = right]}
        \toprule
        {Колонка 1} & {Колонка 2} & {Колонка 3} & {Колонка 4} \\
        \midrule
        2.3456      & 2.3456      & 2.3456      & 2.3456      \\
        34.2345     & 34.2345     & 34.2345     & 34.2345     \\
        56.7835     & 56.7835     & 56.7835     & 56.7835     \\
        90.473      & 90.473      & 90.473      & 90.473      \\
        \bottomrule
    \end{tabular}
  \end{threeparttable}
\end{table}

\section{Параграф "--- два}\label{sec:ch3/sect2}

Некоторый текст.

\section{Параграф с подпараграфами}\label{sec:ch3/sect3}

\subsection{Подпараграф "--- один}\label{subsec:ch3/sect3/sub1}

Некоторый текст.

\subsection{Подпараграф "--- два}\label{subsec:ch3/sect3/sub2}

Некоторый текст.

\clearpage           % Глава 3
\chapter*{Заключение}                       % Заголовок
\addcontentsline{toc}{chapter}{Заключение}  % Добавляем его в оглавление

%% Согласно ГОСТ Р 7.0.11-2011:
%% 5.3.3 В заключении диссертации излагают итоги выполненного исследования, рекомендации, перспективы дальнейшей разработки темы.
%% 9.2.3 В заключении автореферата диссертации излагают итоги данного исследования, рекомендации и перспективы дальнейшей разработки темы.
%% Поэтому имеет смысл сделать эту часть общей и загрузить из одного файла в автореферат и в диссертацию:

Основные результаты работы заключаются в следующем.
%% Согласно ГОСТ Р 7.0.11-2011:
%% 5.3.3 В заключении диссертации излагают итоги выполненного исследования, рекомендации, перспективы дальнейшей разработки темы.
%% 9.2.3 В заключении автореферата диссертации излагают итоги данного исследования, рекомендации и перспективы дальнейшей разработки темы.
\begin{enumerate}
  \item Неэмпирическими расчётными методами было исследовано конформационное поведение бицикло$[3.3.1]$нонана, его замещённых производных и насыщенных гетероциклических аналогов.
  \item Для анализа закономерностей конформационных эффектов замещения аппарат гипергомодесмотических реакций, разработанный Уилером для каркасных и полиненасыщенных углеводородов, был распространён на насыщенные гетероциклы и гетеробициклические соединения.
  \item На основе анализа \ldots
  \item Численные исследования показали, что \ldots
  \item Математическое моделирование показало \ldots
  \item Для выполнения поставленных задач был создан \ldots
\end{enumerate}

И какая-нибудь заключающая фраза.

Последний параграф может включать благодарности.  В заключение автор выражает благодарность и большую признательность научному руководителю Иванову~И.\:И. за~поддержку, помощь, обсуждение результатов и~научное руководство. Также автор благодарит Сидорова~А.\:А. и~Петрова~Б.\:Б. за~помощь в~работе с~образцами, Рабиновича~В.\:В. за предоставленные образцы и~обсуждение результатов, Занудятину~Г.\:Г. и авторов шаблона \tqt{Russian-Phd-LaTeX-Dissertation-Template} для работы в системе \LaTeX{} за~помощь в оформлении диссертации. Автор также благодарит много разных людей и~всех, кто сделал настоящую работу автора возможной. Соответствующие алфавиты следует поблагодарить за предоставленные буквы, шрифты "--- за начертания, а древних кхмеров "--- за цифры, включая \tqt{ноль}~(0).
      % Заключение
\include{Dissertation/acronyms}        % Список сокращений и условных обозначений
\include{Dissertation/dictionary}      % Словарь терминов
\clearpage                                  % В том числе гарантирует, что список литературы в оглавлении будет с правильным номером страницы
%\hypersetup{ urlcolor=black }               % Ссылки делаем чёрными
%\providecommand*{\BibDash}{}                % В стилях ugost2008 отключаем использование тире как разделителя
\urlstyle{rm}                               % ссылки URL обычным шрифтом
\ifdefmacro{\microtypesetup}{\microtypesetup{protrusion=false}}{} % не рекомендуется применять пакет микротипографики к автоматически генерируемому списку литературы
\insertbibliofull                           % Подключаем Bib-базы: все статьи единым списком
% Режим с подсписками
%\insertbiblioexternal                      % Подключаем Bib-базы: статьи, не являющиеся статьями автора по теме диссертации
% Для вывода выберите и расскомментируйте одно из двух
%\insertbiblioauthor                        % Подключаем Bib-базы: работы автора единым списком 
%\insertbiblioauthorgrouped                 % Подключаем Bib-базы: работы автора сгруппированные (ВАК, WoS, Scopus и т.д.)
\ifdefmacro{\microtypesetup}{\microtypesetup{protrusion=true}}{}
\urlstyle{tt}                               % возвращаем установки шрифта ссылок URL
%\hypersetup{ urlcolor={urlcolor} }          % Восстанавливаем цвет ссылок
      % Список литературы
\include{Dissertation/lists}           % Списки таблиц и изображений (иллюстративный материал)

%%% Настройки для приложений
\appendix
% Оформление заголовков приложений ближе к ГОСТ:
\setlength{\midchapskip}{20pt}
\renewcommand*{\afterchapternum}{\par\nobreak\vskip \midchapskip}
\renewcommand\thechapter{\Asbuk{chapter}} % Чтобы приложения русскими буквами нумеровались

\chapter{Примеры вставки листингов программного кода}\label{app:A}

Для крупных листингов есть два способа. Первый красивый, но в нём могут быть
проблемы с поддержкой кириллицы (у вас может встречаться в~комментариях
и печатаемых сообщениях), он представлен на листинге~\ref{lst:hwbeauty}.
\begin{ListingEnv}[!h]% настройки floating аналогичны окружению figure
    \captiondelim{ } % разделитель идентификатора с номером от наименования
    \caption{Программа ,,Hello, world`` на \protect\cpp}\label{lst:hwbeauty}
    % окружение учитывает пробелы и табуляции и применяет их в сответсвии с настройками
    \begin{lstlisting}[language={[ISO]C++}]
	#include <iostream>
	using namespace std;

	int main() //кириллица в комментариях при xelatex и lualatex имеет проблемы с пробелами
	{
		cout << "Hello, world" << endl; //latin letters in commentaries
		system("pause");
		return 0;
	}
    \end{lstlisting}
\end{ListingEnv}%
Второй не~такой красивый, но без ограничений (см.~листинг~\ref{lst:hwplain}).
\begin{ListingEnv}[!h]
    \captiondelim{ } % разделитель идентификатора с номером от наименования
    \caption{Программа ,,Hello, world`` без подсветки}\label{lst:hwplain}
    \begin{Verb}

        #include <iostream>
        using namespace std;

        int main() //кириллица в комментариях
        {
            cout << "Привет, мир" << endl;
        }
    \end{Verb}
\end{ListingEnv}

Можно использовать первый для вставки небольших фрагментов
внутри текста, а второй для вставки полного
кода в приложении, если таковое имеется.

Если нужно вставить совсем короткий пример кода (одна или две строки),
то~выделение  линейками и нумерация может смотреться чересчур громоздко.
В таких случаях можно использовать окружения \texttt{lstlisting} или
\texttt{Verb} без \texttt{ListingEnv}. Приведём такой пример
с указанием языка программирования, отличного от~заданного по умолчанию:
\begin{lstlisting}[language=Haskell]
fibs = 0 : 1 : zipWith (+) fibs (tail fibs)
\end{lstlisting}
Такое решение~--- со вставкой нумерованных листингов покрупнее
и вставок без выделения для маленьких фрагментов~--- выбрано,
например, в книге Эндрю Таненбаума и Тодда Остина по архитектуре
%компьютера~\autocite{TanAus2013} (см.~рис.~\ref{fig:tan-aus}).

Наконец, для оформления идентификаторов внутри строк
(функция \lstinline{main} и~тому подобное) используется
\texttt{lstinline} или, самое простое, моноширинный текст
(\texttt{\textbackslash texttt}).

Пример~\ref{lst:internal3}, иллюстрирующий подключение переопределённого
языка. Может быть полезным, если подсветка кода работает криво. Без
дополнительного окружения, с подписью и ссылкой, реализованной встроенным
средством.
\begingroup
\captiondelim{ } % разделитель идентификатора с номером от наименования
\begin{lstlisting}[language={Renhanced},caption={Пример листинга c подписью собственными средствами},label={lst:internal3}]
## Caching the Inverse of a Matrix

## Matrix inversion is usually a costly computation and there may be some
## benefit to caching the inverse of a matrix rather than compute it repeatedly
## This is a pair of functions that cache the inverse of a matrix.

## makeCacheMatrix creates a special "matrix" object that can cache its inverse

makeCacheMatrix <- function(x = matrix()) {#кириллица в комментариях при xelatex и lualatex имеет проблемы с пробелами
    i <- NULL
    set <- function(y) {
        x <<- y
        i <<- NULL
    }
    get <- function() x
    setSolved <- function(solve) i <<- solve
    getSolved <- function() i
    list(set = set, get = get,
    setSolved = setSolved,
    getSolved = getSolved)

}


## cacheSolve computes the inverse of the special "matrix" returned by
## makeCacheMatrix above. If the inverse has already been calculated (and the
## matrix has not changed), then the cachesolve should retrieve the inverse from
## the cache.

cacheSolve <- function(x, ...) {
    ## Return a matrix that is the inverse of 'x'
    i <- x$getSolved()
    if(!is.null(i)) {
        message("getting cached data")
        return(i)
    }
    data <- x$get()
    i <- solve(data, ...)
    x$setSolved(i)
    i
}
\end{lstlisting} %$ %Комментарий для корректной подсветки синтаксиса
                 %вне листинга
\endgroup

Листинг~\ref{lst:external1} подгружается из внешнего файла. Приходится
загружать без окружения дополнительного. Иначе по страницам не переносится.
\begingroup
\captiondelim{ } % разделитель идентификатора с номером от наименования
    \lstinputlisting[lastline=78,language={R},caption={Листинг из внешнего файла},label={lst:external1}]{listings/run_analysis.R}
\endgroup

\chapter{Очень длинное название второго приложения, в~котором продемонстрирована работа с~длинными таблицами}\label{app:B}

\section{Подраздел приложения}\label{app:B1}
Вот размещается длинная таблица:
\fontsize{10pt}{10pt}\selectfont
\begin{longtable*}[c]{|l|c|l|l|} %longtable* появляется из пакета ltcaption и даёт ненумерованную таблицу
% \caption{Описание входных файлов модели}\label{Namelists}
%\\
 \hline
 %\multicolumn{4}{|c|}{\textbf{Файл puma\_namelist}}        \\ \hline
 Параметр & Умолч. & Тип & Описание               \\ \hline
                                              \endfirsthead   \hline
 \multicolumn{4}{|c|}{\small\slshape (продолжение)}        \\ \hline
 Параметр & Умолч. & Тип & Описание               \\ \hline
                                              \endhead        \hline
% \multicolumn{4}{|c|}{\small\slshape (окончание)}        \\ \hline
% Параметр & Умолч. & Тип & Описание               \\ \hline
%                                             \endlasthead        \hline
 \multicolumn{4}{|r|}{\small\slshape продолжение следует}  \\ \hline
                                              \endfoot        \hline
                                              \endlastfoot
 \multicolumn{4}{|l|}{\&INP}        \\ \hline
 kick & 1 & int & 0: инициализация без шума (\(p_s = const\)) \\
      &   &     & 1: генерация белого шума                  \\
      &   &     & 2: генерация белого шума симметрично относительно \\
  & & & экватора    \\
 mars & 0 & int & 1: инициализация модели для планеты Марс     \\
 kick & 1 & int & 0: инициализация без шума (\(p_s = const\)) \\
      &   &     & 1: генерация белого шума                  \\
      &   &     & 2: генерация белого шума симметрично относительно \\
  & & & экватора    \\
 mars & 0 & int & 1: инициализация модели для планеты Марс     \\
kick & 1 & int & 0: инициализация без шума (\(p_s = const\)) \\
      &   &     & 1: генерация белого шума                  \\
      &   &     & 2: генерация белого шума симметрично относительно \\
  & & & экватора    \\
 mars & 0 & int & 1: инициализация модели для планеты Марс     \\
kick & 1 & int & 0: инициализация без шума (\(p_s = const\)) \\
      &   &     & 1: генерация белого шума                  \\
      &   &     & 2: генерация белого шума симметрично относительно \\
  & & & экватора    \\
 mars & 0 & int & 1: инициализация модели для планеты Марс     \\
kick & 1 & int & 0: инициализация без шума (\(p_s = const\)) \\
      &   &     & 1: генерация белого шума                  \\
      &   &     & 2: генерация белого шума симметрично относительно \\
  & & & экватора    \\
 mars & 0 & int & 1: инициализация модели для планеты Марс     \\
kick & 1 & int & 0: инициализация без шума (\(p_s = const\)) \\
      &   &     & 1: генерация белого шума                  \\
      &   &     & 2: генерация белого шума симметрично относительно \\
  & & & экватора    \\
 mars & 0 & int & 1: инициализация модели для планеты Марс     \\
kick & 1 & int & 0: инициализация без шума (\(p_s = const\)) \\
      &   &     & 1: генерация белого шума                  \\
      &   &     & 2: генерация белого шума симметрично относительно \\
  & & & экватора    \\
 mars & 0 & int & 1: инициализация модели для планеты Марс     \\
kick & 1 & int & 0: инициализация без шума (\(p_s = const\)) \\
      &   &     & 1: генерация белого шума                  \\
      &   &     & 2: генерация белого шума симметрично относительно \\
  & & & экватора    \\
 mars & 0 & int & 1: инициализация модели для планеты Марс     \\
kick & 1 & int & 0: инициализация без шума (\(p_s = const\)) \\
      &   &     & 1: генерация белого шума                  \\
      &   &     & 2: генерация белого шума симметрично относительно \\
  & & & экватора    \\
 mars & 0 & int & 1: инициализация модели для планеты Марс     \\
kick & 1 & int & 0: инициализация без шума (\(p_s = const\)) \\
      &   &     & 1: генерация белого шума                  \\
      &   &     & 2: генерация белого шума симметрично относительно \\
  & & & экватора    \\
 mars & 0 & int & 1: инициализация модели для планеты Марс     \\
kick & 1 & int & 0: инициализация без шума (\(p_s = const\)) \\
      &   &     & 1: генерация белого шума                  \\
      &   &     & 2: генерация белого шума симметрично относительно \\
  & & & экватора    \\
 mars & 0 & int & 1: инициализация модели для планеты Марс     \\
kick & 1 & int & 0: инициализация без шума (\(p_s = const\)) \\
      &   &     & 1: генерация белого шума                  \\
      &   &     & 2: генерация белого шума симметрично относительно \\
  & & & экватора    \\
 mars & 0 & int & 1: инициализация модели для планеты Марс     \\
kick & 1 & int & 0: инициализация без шума (\(p_s = const\)) \\
      &   &     & 1: генерация белого шума                  \\
      &   &     & 2: генерация белого шума симметрично относительно \\
  & & & экватора    \\
 mars & 0 & int & 1: инициализация модели для планеты Марс     \\
kick & 1 & int & 0: инициализация без шума (\(p_s = const\)) \\
      &   &     & 1: генерация белого шума                  \\
      &   &     & 2: генерация белого шума симметрично относительно \\
  & & & экватора    \\
 mars & 0 & int & 1: инициализация модели для планеты Марс     \\
kick & 1 & int & 0: инициализация без шума (\(p_s = const\)) \\
      &   &     & 1: генерация белого шума                  \\
      &   &     & 2: генерация белого шума симметрично относительно \\
  & & & экватора    \\
 mars & 0 & int & 1: инициализация модели для планеты Марс     \\
 \hline
  %& & & \(\:\) \\
 \multicolumn{4}{|l|}{\&SURFPAR}        \\ \hline
kick & 1 & int & 0: инициализация без шума (\(p_s = const\)) \\
      &   &     & 1: генерация белого шума                  \\
      &   &     & 2: генерация белого шума симметрично относительно \\
  & & & экватора    \\
 mars & 0 & int & 1: инициализация модели для планеты Марс     \\
kick & 1 & int & 0: инициализация без шума (\(p_s = const\)) \\
      &   &     & 1: генерация белого шума                  \\
      &   &     & 2: генерация белого шума симметрично относительно \\
  & & & экватора    \\
 mars & 0 & int & 1: инициализация модели для планеты Марс     \\
kick & 1 & int & 0: инициализация без шума (\(p_s = const\)) \\
      &   &     & 1: генерация белого шума                  \\
      &   &     & 2: генерация белого шума симметрично относительно \\
  & & & экватора    \\
 mars & 0 & int & 1: инициализация модели для планеты Марс     \\
kick & 1 & int & 0: инициализация без шума (\(p_s = const\)) \\
      &   &     & 1: генерация белого шума                  \\
      &   &     & 2: генерация белого шума симметрично относительно \\
  & & & экватора    \\
 mars & 0 & int & 1: инициализация модели для планеты Марс     \\
kick & 1 & int & 0: инициализация без шума (\(p_s = const\)) \\
      &   &     & 1: генерация белого шума                  \\
      &   &     & 2: генерация белого шума симметрично относительно \\
  & & & экватора    \\
 mars & 0 & int & 1: инициализация модели для планеты Марс     \\
kick & 1 & int & 0: инициализация без шума (\(p_s = const\)) \\
      &   &     & 1: генерация белого шума                  \\
      &   &     & 2: генерация белого шума симметрично относительно \\
  & & & экватора    \\
 mars & 0 & int & 1: инициализация модели для планеты Марс     \\
kick & 1 & int & 0: инициализация без шума (\(p_s = const\)) \\
      &   &     & 1: генерация белого шума                  \\
      &   &     & 2: генерация белого шума симметрично относительно \\
  & & & экватора    \\
 mars & 0 & int & 1: инициализация модели для планеты Марс     \\
kick & 1 & int & 0: инициализация без шума (\(p_s = const\)) \\
      &   &     & 1: генерация белого шума                  \\
      &   &     & 2: генерация белого шума симметрично относительно \\
  & & & экватора    \\
 mars & 0 & int & 1: инициализация модели для планеты Марс     \\
kick & 1 & int & 0: инициализация без шума (\(p_s = const\)) \\
      &   &     & 1: генерация белого шума                  \\
      &   &     & 2: генерация белого шума симметрично относительно \\
  & & & экватора    \\
 mars & 0 & int & 1: инициализация модели для планеты Марс     \\
 \hline
\end{longtable*}

\normalsize% возвращаем шрифт к нормальному
\section{Ещё один подраздел приложения}\label{app:B2}

Нужно больше подразделов приложения!
Конвынёры витюпырата но нам, тебиквюэ мэнтётюм позтюлант ед про. Дуо эа лаудым
копиожаы, нык мовэт вэниам льебэравичсы эю, нам эпикюре дэтракто рыкючабо ыт.

Пример длинной таблицы с записью продолжения по ГОСТ 2.105:

\begingroup
    \centering
    \small
    \begin{longtable}[c]{|l|c|l|l|}
    \caption{Наименование таблицы средней длины}\label{tab:test5}% label всегда желательно идти после caption
    \\[-0.45\onelineskip]
    \hline
    Параметр & Умолч. & Тип & Описание\\ \hline
    \endfirsthead%
    \caption*{\tabcapalign Продолжение таблицы~\thetable}\\[-0.45\onelineskip]
    \hline
    Параметр & Умолч. & Тип & Описание\\ \hline
    \endhead
    \hline
    \endfoot
    \hline
     \endlastfoot
     \multicolumn{4}{|l|}{\&INP}        \\ \hline
     kick & 1 & int & 0: инициализация без шума (\(p_s = const\)) \\
          &   &     & 1: генерация белого шума                  \\
          &   &     & 2: генерация белого шума симметрично относительно \\
      & & & экватора    \\
     mars & 0 & int & 1: инициализация модели для планеты Марс     \\
     kick & 1 & int & 0: инициализация без шума (\(p_s = const\)) \\
          &   &     & 1: генерация белого шума                  \\
          &   &     & 2: генерация белого шума симметрично относительно \\
      & & & экватора    \\
     mars & 0 & int & 1: инициализация модели для планеты Марс     \\
    kick & 1 & int & 0: инициализация без шума (\(p_s = const\)) \\
          &   &     & 1: генерация белого шума                  \\
          &   &     & 2: генерация белого шума симметрично относительно \\
      & & & экватора    \\
     mars & 0 & int & 1: инициализация модели для планеты Марс     \\
    kick & 1 & int & 0: инициализация без шума (\(p_s = const\)) \\
          &   &     & 1: генерация белого шума                  \\
          &   &     & 2: генерация белого шума симметрично относительно \\
      & & & экватора    \\
     mars & 0 & int & 1: инициализация модели для планеты Марс     \\
    kick & 1 & int & 0: инициализация без шума (\(p_s = const\)) \\
          &   &     & 1: генерация белого шума                  \\
          &   &     & 2: генерация белого шума симметрично относительно \\
      & & & экватора    \\
     mars & 0 & int & 1: инициализация модели для планеты Марс     \\
    kick & 1 & int & 0: инициализация без шума (\(p_s = const\)) \\
          &   &     & 1: генерация белого шума                  \\
          &   &     & 2: генерация белого шума симметрично относительно \\
      & & & экватора    \\
     mars & 0 & int & 1: инициализация модели для планеты Марс     \\
    kick & 1 & int & 0: инициализация без шума (\(p_s = const\)) \\
          &   &     & 1: генерация белого шума                  \\
          &   &     & 2: генерация белого шума симметрично относительно \\
      & & & экватора    \\
     mars & 0 & int & 1: инициализация модели для планеты Марс     \\
    kick & 1 & int & 0: инициализация без шума (\(p_s = const\)) \\
          &   &     & 1: генерация белого шума                  \\
          &   &     & 2: генерация белого шума симметрично относительно \\
      & & & экватора    \\
     mars & 0 & int & 1: инициализация модели для планеты Марс     \\
    kick & 1 & int & 0: инициализация без шума (\(p_s = const\)) \\
          &   &     & 1: генерация белого шума                  \\
          &   &     & 2: генерация белого шума симметрично относительно \\
      & & & экватора    \\
     mars & 0 & int & 1: инициализация модели для планеты Марс     \\
    kick & 1 & int & 0: инициализация без шума (\(p_s = const\)) \\
          &   &     & 1: генерация белого шума                  \\
          &   &     & 2: генерация белого шума симметрично относительно \\
      & & & экватора    \\
     mars & 0 & int & 1: инициализация модели для планеты Марс     \\
    kick & 1 & int & 0: инициализация без шума (\(p_s = const\)) \\
          &   &     & 1: генерация белого шума                  \\
          &   &     & 2: генерация белого шума симметрично относительно \\
      & & & экватора    \\
     mars & 0 & int & 1: инициализация модели для планеты Марс     \\
    kick & 1 & int & 0: инициализация без шума (\(p_s = const\)) \\
          &   &     & 1: генерация белого шума                  \\
          &   &     & 2: генерация белого шума симметрично относительно \\
      & & & экватора    \\
     mars & 0 & int & 1: инициализация модели для планеты Марс     \\
    kick & 1 & int & 0: инициализация без шума (\(p_s = const\)) \\
          &   &     & 1: генерация белого шума                  \\
          &   &     & 2: генерация белого шума симметрично относительно \\
      & & & экватора    \\
     mars & 0 & int & 1: инициализация модели для планеты Марс     \\
    kick & 1 & int & 0: инициализация без шума (\(p_s = const\)) \\
          &   &     & 1: генерация белого шума                  \\
          &   &     & 2: генерация белого шума симметрично относительно \\
      & & & экватора    \\
     mars & 0 & int & 1: инициализация модели для планеты Марс     \\
    kick & 1 & int & 0: инициализация без шума (\(p_s = const\)) \\
          &   &     & 1: генерация белого шума                  \\
          &   &     & 2: генерация белого шума симметрично относительно \\
      & & & экватора    \\
     mars & 0 & int & 1: инициализация модели для планеты Марс     \\
     \hline
      %& & & $\:$ \\
     \multicolumn{4}{|l|}{\&SURFPAR}        \\ \hline
    kick & 1 & int & 0: инициализация без шума (\(p_s = const\)) \\
          &   &     & 1: генерация белого шума                  \\
          &   &     & 2: генерация белого шума симметрично относительно \\
      & & & экватора    \\
     mars & 0 & int & 1: инициализация модели для планеты Марс     \\
    kick & 1 & int & 0: инициализация без шума (\(p_s = const\)) \\
          &   &     & 1: генерация белого шума                  \\
          &   &     & 2: генерация белого шума симметрично относительно \\
      & & & экватора    \\
     mars & 0 & int & 1: инициализация модели для планеты Марс     \\
    kick & 1 & int & 0: инициализация без шума (\(p_s = const\)) \\
          &   &     & 1: генерация белого шума                  \\
          &   &     & 2: генерация белого шума симметрично относительно \\
      & & & экватора    \\
     mars & 0 & int & 1: инициализация модели для планеты Марс     \\
    kick & 1 & int & 0: инициализация без шума (\(p_s = const\)) \\
          &   &     & 1: генерация белого шума                  \\
          &   &     & 2: генерация белого шума симметрично относительно \\
      & & & экватора    \\
     mars & 0 & int & 1: инициализация модели для планеты Марс     \\
    kick & 1 & int & 0: инициализация без шума (\(p_s = const\)) \\
          &   &     & 1: генерация белого шума                  \\
          &   &     & 2: генерация белого шума симметрично относительно \\
      & & & экватора    \\
     mars & 0 & int & 1: инициализация модели для планеты Марс     \\
    kick & 1 & int & 0: инициализация без шума (\(p_s = const\)) \\
          &   &     & 1: генерация белого шума                  \\
          &   &     & 2: генерация белого шума симметрично относительно \\
      & & & экватора    \\
     mars & 0 & int & 1: инициализация модели для планеты Марс     \\
    kick & 1 & int & 0: инициализация без шума (\(p_s = const\)) \\
          &   &     & 1: генерация белого шума                  \\
          &   &     & 2: генерация белого шума симметрично относительно \\
      & & & экватора    \\
     mars & 0 & int & 1: инициализация модели для планеты Марс     \\
    kick & 1 & int & 0: инициализация без шума (\(p_s = const\)) \\
          &   &     & 1: генерация белого шума                  \\
          &   &     & 2: генерация белого шума симметрично относительно \\
      & & & экватора    \\
     mars & 0 & int & 1: инициализация модели для планеты Марс     \\
    kick & 1 & int & 0: инициализация без шума (\(p_s = const\)) \\
          &   &     & 1: генерация белого шума                  \\
          &   &     & 2: генерация белого шума симметрично относительно \\
      & & & экватора    \\
     mars & 0 & int & 1: инициализация модели для планеты Марс     \\
    \end{longtable}
\normalsize% возвращаем шрифт к нормальному
\endgroup
\section{Использование длинных таблиц с окружением \textit{longtabu}}\label{app:B2a}

В таблице~\ref{tab:test-functions} более книжный вариант
длинной таблицы, используя окружение \verb!longtabu! и разнообразные
\verb!toprule! \verb!midrule! \verb!bottomrule! из~пакета
\verb!booktabs!. Чтобы визуально таблица смотрелась лучше, можно
использовать следующие параметры: в самом начале задаётся расстояние
между строчками с~помощью \verb!arraystretch!. Таблица задаётся на
всю ширину, \verb!longtabu! позволяет делить ширину колонок
пропорционально "--- тут три колонки в~пропорции 1.1:1:4 "--- для каждой
колонки первый параметр в~описании \verb!X[]!. Кроме того, в~таблице
убраны отступы слева и справа с~помощью \verb!@{}!
в~преамбуле таблицы. К~первому и~второму столбцу применяется
модификатор

\verb!>{\setlength{\baselineskip}{0.7\baselineskip}}!,

\noindent который уменьшает межстрочный интервал в для текста таблиц (иначе
заголовок второго столбца значительно шире, а двухстрочное имя
сливается с~окружающими). Для первой и второй колонки текст в ячейках
выравниваются по~центру как по~вертикали, так и по горизонтали "---
задаётся буквами \verb!m!~и~\verb!c!~в~описании столбца \verb!X[]!.

Так как формулы большие "--- используется окружение \verb!alignedat!,
чтобы отступ был одинаковый у всех формул "--- он сделан для всех, хотя
для большей части можно было и не использовать.  Чтобы формулы
занимали поменьше места в~каждом столбце формулы (где надо)
используется \verb!\textstyle! "--- он~делает дроби меньше, у~знаков
суммы и произведения "--- индексы сбоку. Иногда формулы слишком большая,
сливается со следующей, поэтому после неё ставится небольшой
дополнительный отступ \verb!\vspace*{2ex}!  Для штрафных функций "---
размер фигурных скобок задан вручную \verb!\Big\{!, т.\:к. не~умеет
\verb!alignedat! работать с~\verb!\left! и~\verb!\right! через
несколько строк/колонок.

В примечании к таблице наоборот, окружение \verb!cases! даёт слишком
большие промежутки между вариантами, чтобы их уменьшить, в конце
каждой строчки окружения использовался отрицательный дополнительный
отступ \verb!\\[-0.5em]!.

\begingroup % Ограничиваем область видимости arraystretch
\renewcommand{\arraystretch}{1.6}%% Увеличение расстояния между рядами, для улучшения восприятия.
\begin{longtabu} to \textwidth
{%
@{}>{\setlength{\baselineskip}{0.7\baselineskip}}X[1.1mc]%
>{\setlength{\baselineskip}{0.7\baselineskip}}X[1.1mc]%
X[4]@{}%
}
    \caption{Тестовые функции для оптимизации, \(D\) "---
      размерность. Для всех функций значение в точке глобального
      минимума равно нулю.\label{tab:test-functions}}\\% label всегда желательно идти после caption

    \toprule     %%% верхняя линейка
    Имя           &Стартовый диапазон параметров &Функция  \\
    \midrule %%% тонкий разделитель. Отделяет названия столбцов. Обязателен по ГОСТ 2.105 пункт 4.4.5
    \endfirsthead

    \multicolumn{3}{c}{\small\slshape (продолжение)}        \\
    \toprule     %%% верхняя линейка
    Имя           &Стартовый диапазон параметров &Функция  \\
    \midrule %%% тонкий разделитель. Отделяет названия столбцов. Обязателен по ГОСТ 2.105 пункт 4.4.5
    \endhead

    \multicolumn{3}{c}{\small\slshape (окончание)}        \\
    \toprule     %%% верхняя линейка
    Имя           &Стартовый диапазон параметров &Функция  \\
    \midrule %%% тонкий разделитель. Отделяет названия столбцов. Обязателен по ГОСТ 2.105 пункт 4.4.5
    \endlasthead

    \bottomrule %%% нижняя линейка
    \multicolumn{3}{r}{\small\slshape продолжение следует}  \\
    \endfoot
    \endlastfoot

    сфера         &\(\left[-100,\,100\right]^D\)   &
        \(\begin{aligned}
            \textstyle f_1(x)=\sum_{i=1}^Dx_i^2
        \end{aligned}\) \\
    Schwefel 2.22 &\(\left[-10,\,10\right]^D\)     &
        \(\begin{aligned}
            \textstyle f_2(x)=\sum_{i=1}^D|x_i|+\prod_{i=1}^D|x_i|
        \end{aligned}\) \\
    Schwefel 1.2  &\(\left[-100,\,100\right]^D\)   &
        \(\begin{aligned}
            \textstyle f_3(x)=\sum_{i=1}^D\left(\sum_{j=1}^ix_j\right)^2
        \end{aligned}\) \\
    Schwefel 2.21 &\(\left[-100,\,100\right]^D\)   &
        \(\begin{aligned}
            \textstyle f_4(x)=\max_i\!\left\{\left|x_i\right|\right\}
        \end{aligned}\) \\
    Rosenbrock    &\(\left[-30,\,30\right]^D\)     &
        \(\begin{aligned}
            \textstyle f_5(x)=
            \sum_{i=1}^{D-1}
            \left[100\!\left(x_{i+1}-x_i^2\right)^2+(x_i-1)^2\right]
        \end{aligned}\) \\
    ступенчатая   &\(\left[-100,\,100\right]^D\)   &
        \(\begin{aligned}
            \textstyle f_6(x)=\sum_{i=1}^D\big\lfloor x_i+0.5\big\rfloor^2
        \end{aligned}\) \\
    зашумлённая квартическая &\(\left[-1.28,\,1.28\right]^D\) &
        \(\begin{aligned}
            \textstyle f_7(x)=\sum_{i=1}^Dix_i^4+rand[0,1)
        \end{aligned}\)\vspace*{2ex}\\
    Schwefel 2.26 &\(\left[-500,\,500\right]^D\)   &
        \(\begin{aligned}
        f_8(x)= &\textstyle\sum_{i=1}^D-x_i\,\sin\sqrt{|x_i|}\,+ \\
                &\vphantom{\sum}+ D\cdot
                418.98288727243369
        \end{aligned}\)\\
    Rastrigin     &\(\left[-5.12,\,5.12\right]^D\) &
    \(\begin{aligned}
        \textstyle f_9(x)=\sum_{i=1}^D\left[x_i^2-10\,\cos(2\pi x_i)+10\right]
    \end{aligned}\)\vspace*{2ex}\\
    Ackley        &\(\left[-32,\,32\right]^D\)     &
        \(\begin{aligned}
            f_{10}(x)= &\textstyle -20\, \exp\!\left(
                            -0.2\sqrt{\frac{1}{D}\sum_{i=1}^Dx_i^2} \right)-\\
                       &\textstyle - \exp\left(
                            \frac{1}{D}\sum_{i=1}^D\cos(2\pi x_i)  \right)
                       + 20 + e
        \end{aligned}\) \\
    Griewank      &\(\left[-600,\,600\right]^D\) &
        \(\begin{aligned}
            f_{11}(x)= &\textstyle \frac{1}{4000}\sum_{i=1}^{D}x_i^2 -
                \prod_{i=1}^D\cos\left(x_i/\sqrt{i}\right) +1
        \end{aligned}\) \vspace*{3ex} \\
    штрафная 1    &\(\left[-50,\,50\right]^D\)     &
        \(\begin{aligned}
            f_{12}(x)= &\textstyle \frac{\pi}{D}\Big\{ 10\,\sin^2(\pi y_1) +\\
            &+\textstyle \sum_{i=1}^{D-1}(y_i-1)^2
                \left[1+10\,\sin^2(\pi y_{i+1})\right] +\\
            &+(y_D-1)^2 \Big\} +\textstyle\sum_{i=1}^D u(x_i,\,10,\,100,\,4)
        \end{aligned}\) \vspace*{2ex} \\
    штрафная 2    &\(\left[-50,\,50\right]^D\)     &
        \(\begin{aligned}
            f_{13}(x)= &\textstyle 0.1 \Big\{\sin^2(3\pi x_1) +\\
            &+\textstyle \sum_{i=1}^{D-1}(x_i-1)^2
                \left[1+\sin^2(3 \pi x_{i+1})\right] + \\
            &+(x_D-1)^2\left[1+\sin^2(2\pi x_D)\right] \Big\} +\\
            &+\textstyle\sum_{i=1}^D u(x_i,\,5,\,100,\,4)
        \end{aligned}\)\\
    сфера         &\(\left[-100,\,100\right]^D\)   &
        \(\begin{aligned}
            \textstyle f_1(x)=\sum_{i=1}^Dx_i^2
        \end{aligned}\) \\
    Schwefel 2.22 &\(\left[-10,\,10\right]^D\)     &
        \(\begin{aligned}
            \textstyle f_2(x)=\sum_{i=1}^D|x_i|+\prod_{i=1}^D|x_i|
        \end{aligned}\) \\
    Schwefel 1.2  &\(\left[-100,\,100\right]^D\)   &
        \(\begin{aligned}
            \textstyle f_3(x)=\sum_{i=1}^D\left(\sum_{j=1}^ix_j\right)^2
        \end{aligned}\) \\
    Schwefel 2.21 &\(\left[-100,\,100\right]^D\)   &
        \(\begin{aligned}
            \textstyle f_4(x)=\max_i\!\left\{\left|x_i\right|\right\}
        \end{aligned}\) \\
    Rosenbrock    &\(\left[-30,\,30\right]^D\)     &
        \(\begin{aligned}
            \textstyle f_5(x)=
            \sum_{i=1}^{D-1}
            \left[100\!\left(x_{i+1}-x_i^2\right)^2+(x_i-1)^2\right]
        \end{aligned}\) \\
    ступенчатая   &\(\left[-100,\,100\right]^D\)   &
        \(\begin{aligned}
            \textstyle f_6(x)=\sum_{i=1}^D\big\lfloor x_i+0.5\big\rfloor^2
        \end{aligned}\) \\
    зашумлённая квартическая &\(\left[-1.28,\,1.28\right]^D\) &
        \(\begin{aligned}
            \textstyle f_7(x)=\sum_{i=1}^Dix_i^4+rand[0,1)
        \end{aligned}\)\vspace*{2ex}\\
    Schwefel 2.26 &\(\left[-500,\,500\right]^D\)   &
        \(\begin{aligned}
        f_8(x)= &\textstyle\sum_{i=1}^D-x_i\,\sin\sqrt{|x_i|}\,+ \\
                &\vphantom{\sum}+ D\cdot
                418.98288727243369
        \end{aligned}\)\\
    Rastrigin     &\(\left[-5.12,\,5.12\right]^D\) &
    \(\begin{aligned}
        \textstyle f_9(x)=\sum_{i=1}^D\left[x_i^2-10\,\cos(2\pi x_i)+10\right]
    \end{aligned}\)\vspace*{2ex}\\
    Ackley        &\(\left[-32,\,32\right]^D\)     &
        \(\begin{aligned}
            f_{10}(x)= &\textstyle -20\, \exp\!\left(
                            -0.2\sqrt{\frac{1}{D}\sum_{i=1}^Dx_i^2} \right)-\\
                       &\textstyle - \exp\left(
                            \frac{1}{D}\sum_{i=1}^D\cos(2\pi x_i)  \right)
                       + 20 + e
        \end{aligned}\) \\
    Griewank      &\(\left[-600,\,600\right]^D\) &
        \(\begin{aligned}
            f_{11}(x)= &\textstyle \frac{1}{4000}\sum_{i=1}^{D}x_i^2 -
                \prod_{i=1}^D\cos\left(x_i/\sqrt{i}\right) +1
        \end{aligned}\) \vspace*{3ex} \\
    штрафная 1    &\(\left[-50,\,50\right]^D\)     &
        \(\begin{aligned}
            f_{12}(x)= &\textstyle \frac{\pi}{D}\Big\{ 10\,\sin^2(\pi y_1) +\\
            &+\textstyle \sum_{i=1}^{D-1}(y_i-1)^2
                \left[1+10\,\sin^2(\pi y_{i+1})\right] +\\
            &+(y_D-1)^2 \Big\} +\textstyle\sum_{i=1}^D u(x_i,\,10,\,100,\,4)
        \end{aligned}\) \vspace*{2ex} \\
    штрафная 2    &\(\left[-50,\,50\right]^D\)     &
        \(\begin{aligned}
            f_{13}(x)= &\textstyle 0.1 \Big\{\sin^2(3\pi x_1) +\\
            &+\textstyle \sum_{i=1}^{D-1}(x_i-1)^2
                \left[1+\sin^2(3 \pi x_{i+1})\right] + \\
            &+(x_D-1)^2\left[1+\sin^2(2\pi x_D)\right] \Big\} +\\
            &+\textstyle\sum_{i=1}^D u(x_i,\,5,\,100,\,4)
        \end{aligned}\)\\
    \midrule%%% тонкий разделитель
    \multicolumn{3}{@{}p{\textwidth}}{%
        \vspace*{-3.5ex}% этим подтягиваем повыше
        \hspace*{2.5em}% абзацный отступ - требование ГОСТ 2.105
        Примечание "---  Для функций \(f_{12}\) и \(f_{13}\)
        используется \(y_i = 1 + \frac{1}{4}(x_i+1)\)
        и~$u(x_i,\,a,\,k,\,m)=
            \begin{cases*}
                k(x_i-a)^m,& \( x_i >a \)\\[-0.5em]
                0,& \( -a\leq x_i \leq a \)\\[-0.5em]
                k(-x_i-a)^m,& \( x_i <-a \)
            \end{cases*}
        $
}\\
\bottomrule %%% нижняя линейка
\end{longtabu}
\endgroup

\section{Форматирование внутри таблиц}\label{app:B3}

В таблице~\ref{tab:other-row} пример с чересстрочным
форматированием. В~файле \verb+userstyles.tex+  задаётся счётчик
\verb+\newcounter{rowcnt}+ который увеличивается на~1 после каждой
строчки (как указано в преамбуле таблицы). Кроме того, задаётся
условный макрос \verb+\altshape+ который выдаёт одно
из~двух типов форматирования в~зависимости от чётности счётчика.

В таблице~\ref{tab:other-row} каждая чётная строчка "--- синяя,
нечётная "--- с наклоном и~слегка поднята вверх. Визуально это приводит
к тому, что среднее значение и~среднеквадратичное изменение
группируются и хорошо выделяются взглядом в~таблице. Сохраняется
возможность отдельные значения в таблице выделить цветом или
шрифтом. К первому и второму столбцу форматирование не применяется
по~сути таблицы, к шестому общее форматирование не~применяется для
наглядности.

Так как заголовок таблицы тоже считается за строчку, то перед ним (для
первого, промежуточного и финального варианта) счётчик обнуляется,
а~в~\verb+\altshape+ для нулевого значения счётчика форматирования
не~применяется.

\begingroup % Ограничиваем область видимости arraystretch
\renewcommand\altshape{
  \ifnumequal{\value{rowcnt}}{0}{
    % Стиль для заголовка таблицы
  }{
    \ifnumodd{\value{rowcnt}}
    {
      \color{blue} % Cтиль для нечётных строк
    }{
      \vspace*{-0.7ex}\itshape} % Стиль для чётных строк
  }
}
\newcolumntype{A}{>{\centering\begingroup\altshape}X[1mc]<{\endgroup}}
\needspace{2\baselineskip}
\renewcommand{\arraystretch}{0.9}%% Уменьшаем  расстояние между
                                %% рядами, чтобы таблица не так много
                                %% места занимала в дисере.
\begin{longtabu} to \textwidth {@{}X[0.27ml]@{}X[0.7mc]@{}A@{}A@{}A@{}X[0.98mc]@{}>{\setlength{\baselineskip}{0.7\baselineskip}}A@{}A<{\stepcounter{rowcnt}}@{}}
% \begin{longtabu} to \textwidth {@{}X[0.2ml]X[1mc]X[1mc]X[1mc]X[1mc]X[1mc]>{\setlength{\baselineskip}{0.7\baselineskip}}X[1mc]X[1mc]@{}}
  \caption{Длинная таблица с примером чересстрочного форматирования\label{tab:other-row}}\vspace*{1ex}\\% label всегда желательно идти после caption
  % \vspace*{1ex}     \\

  \toprule %%% верхняя линейка
\setcounter{rowcnt}{0} &Итера\-ции & JADE\texttt{++} & JADE & jDE & SaDE
& DE/rand /1/bin & PSO \\
 \midrule %%% тонкий разделитель. Отделяет названия столбцов. Обязателен по ГОСТ 2.105 пункт 4.4.5
 \endfirsthead

 \multicolumn{8}{c}{\small\slshape (продолжение)} \\
 \toprule %%% верхняя линейка
\setcounter{rowcnt}{0} &Итера\-ции & JADE\texttt{++} & JADE & jDE & SaDE
& DE/rand /1/bin & PSO \\
 \midrule %%% тонкий разделитель. Отделяет названия столбцов. Обязателен по ГОСТ 2.105 пункт 4.4.5
 \endhead

 \multicolumn{8}{c}{\small\slshape (окончание)} \\
 \toprule %%% верхняя линейка
\setcounter{rowcnt}{0} &Итера\-ции & JADE\texttt{++} & JADE & jDE & SaDE
& DE/rand /1/bin & PSO \\
 \midrule %%% тонкий разделитель. Отделяет названия столбцов. Обязателен по ГОСТ 2.105 пункт 4.4.5
 \endlasthead

 \bottomrule %%% нижняя линейка
 \multicolumn{8}{r}{\small\slshape продолжение следует}     \\
 \endfoot
 \endlastfoot

f1  & 1500 & \textbf{1.8E-60}   & 1.3E-54   & 2.5E-28   & 4.5E-20   & 9.8E-14   & 9.6E-42   \\\nopagebreak
    &      & (8.4E-60) & (9.2E-54) & {\color{red}(3.5E-28)} & (6.9E-20) & (8.4E-14) & (2.7E-41) \\
f2  & 2000 & 1.8E-25   & 3.9E-22   & 1.5E-23   & 1.9E-14   & 1.6E-09   & 9.3E-21   \\\nopagebreak
    &      & (8.8E-25) & (2.7E-21) & (1.0E-23) & (1.1E-14) & (1.1E-09) & (6.3E-20) \\
f3  & 5000 & 5.7E-61   & 6.0E-87   & 5.2E-14   & {\color{green}9.0E-37}   & 6.6E-11   & 2.5E-19   \\\nopagebreak
    &      & (2.7E-60) & (1.9E-86) & (1.1E-13) & (5.4E-36) & (8.8E-11) & (3.9E-19) \\
f4  & 5000 & 8.2E-24   & 4.3E-66   & 1.4E-15   & 7.4E-11   & 4.2E-01   & 4.4E-14   \\\nopagebreak
    &      & (4.0E-23) & (1.2E-65) & (1.0E-15) & (1.8E-10) & (1.1E+00) & (9.3E-14) \\
f5  & 3000 & 8.0E-02   & 3.2E-01   & 1.3E+01   & 2.1E+01   & 2.1E+00   & 2.5E+01   \\\nopagebreak
    &      & (5.6E-01) & (1.1E+00) & (1.4E+01) & (7.8E+00) & (1.5E+00) & (3.2E+01) \\
f6  & 100  & 2.9E+00   & 5.6E+00   & 1.0E+03   & 9.3E+02   & 4.7E+03   & 4.5E+01   \\\nopagebreak
    &      & (1.2E+00) & (1.6E+00) & (2.2E+02) & (1.8E+02) & (1.1E+03) & (2.4E+01) \\
f7  & 3000 & 6.4E-04   & 6.8E-04   & 3.3E-03   & 4.8E-03   & 4.7E-03   & 2.5E-03   \\\nopagebreak
    &      & (2.5E-04) & (2.5E-04) & (8.5E-04) & (1.2E-03) & (1.2E-03) & (1.4E-03) \\
f8  & 1000 & 3.3E-05   & 7.1E+00   & 7.9E-11   & 4.7E+00   & 5.9E+03   & 2.4E+03   \\\nopagebreak
    &      & (2.3E-05) & (2.8E+01) & (1.3E-10) & (3.3E+01) & (1.1E+03) & (6.7E+02) \\
f9  & 1000 & 1.0E-04   & 1.4E-04   & 1.5E-04   & 1.2E-03   & 1.8E+02   & 5.2E+01   \\\nopagebreak
    &      & (6.0E-05) & (6.5E-05) & (2.0E-04) & (6.5E-04) & (1.3E+01) & (1.6E+01) \\
f10 & 500  & 8.2E-10   & 3.0E-09   & 3.5E-04   & 2.7E-03   & 1.1E-01   & 4.6E-01   \\\nopagebreak
    &      & (6.9E-10) & (2.2E-09) & (1.0E-04) & (5.1E-04) & (3.9E-02) & (6.6E-01) \\
f11 & 500  & 9.9E-08   & 2.0E-04   & 1.9E-05   & 7.8E-04  & 2.0E-01   & 1.3E-02   \\\nopagebreak
    &      & (6.0E-07) & (1.4E-03) & (5.8E-05) & (1.2E-03)  & (1.1E-01) & (1.7E-02) \\
f12 & 500  & 4.6E-17   & 3.8E-16   & 1.6E-07   & 1.9E-05   & 1.2E-02   & 1.9E-01   \\\nopagebreak
    &      & (1.9E-16) & (8.3E-16) & (1.5E-07) & (9.2E-06) & (1.0E-02) & (3.9E-01) \\
f13 & 500  & 2.0E-16   & 1.2E-15   & 1.5E-06   & 6.1E-05   & 7.5E-02   & 2.9E-03   \\\nopagebreak
    &      & (6.5E-16) & (2.8E-15) & (9.8E-07) & (2.0E-05) & (3.8E-02) & (4.8E-03) \\
f1  & 1500 & \textbf{1.8E-60}   & 1.3E-54   & 2.5E-28   & 4.5E-20   & 9.8E-14   & 9.6E-42   \\\nopagebreak
    &      & (8.4E-60) & (9.2E-54) & {\color{red}(3.5E-28)} & (6.9E-20) & (8.4E-14) & (2.7E-41) \\
f2  & 2000 & 1.8E-25   & 3.9E-22   & 1.5E-23   & 1.9E-14   & 1.6E-09   & 9.3E-21   \\\nopagebreak
    &      & (8.8E-25) & (2.7E-21) & (1.0E-23) & (1.1E-14) & (1.1E-09) & (6.3E-20) \\
f3  & 5000 & 5.7E-61   & 6.0E-87   & 5.2E-14   & 9.0E-37   & 6.6E-11   & 2.5E-19   \\\nopagebreak
    &      & (2.7E-60) & (1.9E-86) & (1.1E-13) & (5.4E-36) & (8.8E-11) & (3.9E-19) \\
f4  & 5000 & 8.2E-24   & 4.3E-66   & 1.4E-15   & 7.4E-11   & 4.2E-01   & 4.4E-14   \\\nopagebreak
    &      & (4.0E-23) & (1.2E-65) & (1.0E-15) & (1.8E-10) & (1.1E+00) & (9.3E-14) \\
f5  & 3000 & 8.0E-02   & 3.2E-01   & 1.3E+01   & 2.1E+01   & 2.1E+00   & 2.5E+01   \\\nopagebreak
    &      & (5.6E-01) & (1.1E+00) & (1.4E+01) & (7.8E+00) & (1.5E+00) & (3.2E+01) \\
f6  & 100  & 2.9E+00   & 5.6E+00   & 1.0E+03   & 9.3E+02   & 4.7E+03   & 4.5E+01   \\\nopagebreak
    &      & (1.2E+00) & (1.6E+00) & (2.2E+02) & (1.8E+02) & (1.1E+03) & (2.4E+01) \\
f7  & 3000 & 6.4E-04   & 6.8E-04   & 3.3E-03   & 4.8E-03   & 4.7E-03   & 2.5E-03   \\\nopagebreak
    &      & (2.5E-04) & (2.5E-04) & (8.5E-04) & (1.2E-03) & (1.2E-03) & (1.4E-03) \\
f8  & 1000 & 3.3E-05   & 7.1E+00   & 7.9E-11   & 4.7E+00   & 5.9E+03   & 2.4E+03   \\\nopagebreak
    &      & (2.3E-05) & (2.8E+01) & (1.3E-10) & (3.3E+01) & (1.1E+03) & (6.7E+02) \\
f9  & 1000 & 1.0E-04   & 1.4E-04   & 1.5E-04   & 1.2E-03   & 1.8E+02   & 5.2E+01   \\\nopagebreak
    &      & (6.0E-05) & (6.5E-05) & (2.0E-04) & (6.5E-04) & (1.3E+01) & (1.6E+01) \\
f10 & 500  & 8.2E-10   & 3.0E-09   & 3.5E-04   & 2.7E-03   & 1.1E-01   & 4.6E-01   \\\nopagebreak
    &      & (6.9E-10) & (2.2E-09) & (1.0E-04) & (5.1E-04) & (3.9E-02) & (6.6E-01) \\
f11 & 500  & 9.9E-08   & 2.0E-04   & 1.9E-05   & 7.8E-04  & 2.0E-01   & 1.3E-02   \\\nopagebreak
    &      & (6.0E-07) & (1.4E-03) & (5.8E-05) & (1.2E-03)  & (1.1E-01) & (1.7E-02) \\
f12 & 500  & 4.6E-17   & 3.8E-16   & 1.6E-07   & 1.9E-05   & 1.2E-02   & 1.9E-01   \\\nopagebreak
    &      & (1.9E-16) & (8.3E-16) & (1.5E-07) & (9.2E-06) & (1.0E-02) & (3.9E-01) \\
f13 & 500  & 2.0E-16   & 1.2E-15   & 1.5E-06   & 6.1E-05   & 7.5E-02   & 2.9E-03   \\\nopagebreak
    &      & (6.5E-16) & (2.8E-15) & (9.8E-07) & (2.0E-05) & (3.8E-02) & (4.8E-03) \\
\bottomrule %%% нижняя линейка
\end{longtabu} \endgroup

\section{Стандартные префиксы ссылок}\label{app:B4}

Общепринятым является следующий формат ссылок: \texttt{<prefix>:<label>}.
Например, \verb+\label{fig:knuth}+; \verb+\ref{tab:test1}+; \verb+label={lst:external1}+.
В~таблице~\ref{tab:tab_pref} приведены стандартные префиксы для различных типов ссылок.

\begingroup
    \centering
    % \small
    \begin{longtable}[c]{|c|c|}
    \caption{Стандартные префиксы ссылок}\label{tab:tab_pref}% label всегда желательно идти после caption
    \\[-0.45\onelineskip]
    \hline
    \textbf{Префикс} & \textbf{Описание} \\ \hline
    \endfirsthead%
    \caption*{\tabcapalign Продолжение таблицы~\thetable}\\[-0.45\onelineskip]
    \hline
    \textbf{Префикс} & \textbf{Описание} \\ \hline
    \endhead
    \hline
    \endfoot
    \hline
    \endlastfoot
    ch:     & Глава             \\
    sec:    & Секция            \\
    subsec: & Подсекция         \\
    fig:    & Рисунок           \\
    tab:    & Таблица           \\
    eq:     & Уравнение         \\
    lst:    & Листинг программы \\
    itm:    & Элемент списка    \\
    alg:    & Алгоритм          \\
    app:    & Секция приложения \\
    \end{longtable}
% \normalsize% возвращаем шрифт к нормальному
\endgroup

Для упорядочивания ссылок можно использовать разделительные символы.
Например, \verb+\label{fig:scheemes/my_scheeme}+ или \\ \verb+\label{lst:dts/linked_list}+.

\section{Очередной подраздел приложения}~\label{app:B5}

Нужно больше подразделов приложения!

\section{И ещё один подраздел приложения}~\label{app:B6}

Нужно больше подразделов приложения!

\chapter{Потенциалы внутреннего вращения молекул}

Расчёт проводился в приближении RI-MP2/aug-cc-pVTZ без явного учёта симметрии.

\twocolumn
  \begin{table}
    \caption{\label{tab:InternalRotation:Butane}Внутреннее вращение \emph{н}-бутана~\cmpd{Butane}}\centering
    \pgfplotstabletypeset[column type=r,columns={Tau,Etot,Erel}]\ScanButane
  \end{table}
  
  \begin{table}
    \caption{\label{tab:InternalRotation:MeOEt}Внутреннее вращение метилэтилового эфира~\cmpd{MeOEt}}\centering
    \pgfplotstabletypeset[column type=r,columns={Tau,Etot,Erel}]\ScanMeOEt
  \end{table}
  
  \begin{table}
    \caption{\label{tab:InternalRotation:MeSEt}Внутреннее вращение метилэтилсульфида~\cmpd{MeSEt}}\centering
    \pgfplotstabletypeset[column type=r,columns={Tau,Etot,Erel}]\ScanMeSEt
  \end{table}
  
  \begin{table}
    \caption{\label{tab:InternalRotation:Butanone}Внутреннее вращение бутанона-2~\cmpd{Butanone}}\centering
    \pgfplotstabletypeset[column type=r,columns={Tau,Etot,Erel}]\ScanMeCOEt
  \end{table}
  
  \begin{table}
    \caption{\label{tab:InternalRotation:Isopentane}Внутреннее вращение изопентана~\cmpd{Isopentane}}\centering
    \pgfplotstabletypeset[column type=r,columns={Tau,Etot,Erel}]\ScanIsoPentane
  \end{table}
  
  \begin{table}
    \caption{\label{tab:InternalRotation:Me2CHOMe}Внутреннее вращение метил-\emph{изо}-пропилового эфира~\cmpd{Me2CHOMe}}\centering
    \pgfplotstabletypeset[column type=r,columns={Tau,Etot,Erel}]\ScanIPrOMe
  \end{table}
  
  \begin{table}
    \caption{\label{tab:InternalRotation:Me2CHSMe}Внутреннее вращение метил-\emph{изо}-пропилсульфида~\cmpd{Me2CHSMe}}\centering
    \pgfplotstabletypeset[column type=r,columns={Tau,Etot,Erel}]\ScanIPrSMe
  \end{table}
  
  \begin{table}
    \caption{\label{tab:InternalRotation:Me2NEt}Внутреннее вращение \ce{N},\ce{N}-диметилэтиламина~\cmpd{Me2NEt}}\centering
    \pgfplotstabletypeset[column type=r,columns={Tau,Etot,Erel}]\ScanMeMeNEt
  \end{table}
\onecolumn
        % Приложения

\end{document}
